% oss-plantilla: oss-2026.11
% Preámbulo corporativo fijo. No editar en sd-NN.tex: exportar_latex.py verificar lo compara byte a byte.
% Para cambiar el formato, editar esta plantilla u oss.sty y subir la versión de la primera línea.
\PassOptionsToPackage{unicode}{hyperref}
\PassOptionsToPackage{hyphens}{url}
\documentclass[11pt]{article}
\usepackage{amsmath,amssymb}
\usepackage{iftex}
\ifPDFTeX
  % Prism y otros editores en la nube: pdfLaTeX con Helvetica como familia por defecto.
  \usepackage[T1]{fontenc}
  \usepackage[utf8]{inputenc}
  \usepackage{textcomp}
  \usepackage[scaled=0.92]{helvet}
  \renewcommand{\familydefault}{\sfdefault}
  \DeclareUnicodeCharacter{2264}{\ensuremath{\leq}}
  \DeclareUnicodeCharacter{2265}{\ensuremath{\geq}}
  \DeclareUnicodeCharacter{2260}{\ensuremath{\neq}}
  \DeclareUnicodeCharacter{2248}{\ensuremath{\approx}}
  \DeclareUnicodeCharacter{2190}{\ensuremath{\leftarrow}}
  \DeclareUnicodeCharacter{2192}{\ensuremath{\rightarrow}}
  \DeclareUnicodeCharacter{25CF}{\ensuremath{\bullet}}
  \DeclareUnicodeCharacter{2713}{\ensuremath{\checkmark}}
\else
  % XeLaTeX/LuaLaTeX (compilación oficial local): DejaVu Sans.
  \usepackage{fontspec}
  \defaultfontfeatures{Ligatures=TeX,Scale=MatchLowercase}
  \setmainfont{DejaVu Sans}
  \setsansfont{DejaVu Sans}
  \setmonofont{DejaVu Sans Mono}
\fi
\usepackage[letterpaper,margin=20mm]{geometry}
\usepackage{microtype}
\usepackage{parskip}
\usepackage{xcolor}
\usepackage{longtable,booktabs,array}
\usepackage{calc}
\usepackage{etoolbox}
\makeatletter
\patchcmd\longtable{\par}{\if@noskipsec\mbox{}\fi\par}{}{}
\makeatother
\usepackage{footnotehyper}
\makesavenoteenv{longtable}
\usepackage{graphicx}
\makeatletter
\newsavebox\pandoc@box
\newcommand*\pandocbounded[1]{%
  \sbox\pandoc@box{#1}%
  \Gscale@div\@tempa{\textheight}{\dimexpr\ht\pandoc@box+\dp\pandoc@box\relax}%
  \Gscale@div\@tempb{\linewidth}{\wd\pandoc@box}%
  \ifdim\@tempb\p@<\@tempa\p@\let\@tempa\@tempb\fi
  \ifdim\@tempa\p@<\p@\scalebox{\@tempa}{\usebox\pandoc@box}%
  \else\usebox{\pandoc@box}%
  \fi}
\def\maxwidth{\ifdim\Gin@nat@width>\linewidth\linewidth\else\Gin@nat@width\fi}
\def\maxheight{\ifdim\Gin@nat@height>\textheight\textheight\else\Gin@nat@height\fi}
\makeatother
\setkeys{Gin}{width=\maxwidth,height=\maxheight,keepaspectratio}
\makeatletter
\def\fps@figure{htbp}
\makeatother
\usepackage[normalem]{ulem}
\providecommand{\st}[1]{\sout{#1}}
\usepackage{fancyvrb}
\setlength{\emergencystretch}{3em}
\providecommand{\tightlist}{%
  \setlength{\itemsep}{0pt}\setlength{\parskip}{0pt}}
\providecommand{\passthrough}[1]{#1}
\setcounter{secnumdepth}{-\maxdimen}
\ifPDFTeX
  \usepackage[spanish]{babel}
\else
  \usepackage[bidi=default]{babel}
  \babelprovide[main,import]{spanish}
\fi
\let\LanguageShortHands\languageshorthands
\def\languageshorthands#1{}
\usepackage{bookmark}
\usepackage{xurl}
\urlstyle{same}
\hypersetup{pdflang={es-CL},pdfcreator={Only Simple Solutions - exportar_latex.py}}
\usepackage{oss}
\begin{document}

\hypersetup{pageanchor=false}
\ossCover{Anexos}{Subdocumento 2}
\hypersetup{pageanchor=true}
\makeatletter\renewcommand{\@pnumwidth}{2.2em}\makeatother
\tableofcontents
\clearpage

\renewcommand{\tablename}{Tabla}
\renewcommand{\listtablename}{Lista de tablas}
\makeatletter\renewcommand*\l@table{\@dottedtocline{1}{0em}{4.2em}}\makeatother
\listoftables
\clearpage

\section{A. Anexos}\label{a-anexos}

\subsection{A.1. Listado de Requerimientos Funcionales (RF)}\label{a-1-listado-de-requerimientos-funcionales-rf}

Este anexo consolida, la totalidad de los requerimientos funcionales (RF), no funcionales (RNF) y las obligaciones del proponente (OP)

\subsubsection{A.1.1. Listado de Requerimientos Funcionales: Ventas en Tienda}\label{a-1-1-listado-de-requerimientos-funcionales-ventas-en-tienda}

A continuación, se detalla el listado de requerimientos funcionales (RF) del módulo/categoría ``Ventas en Tienda''.

\renewcommand{\thetable}{A.1.1}
\begin{longtable}[]{@{}
  >{\raggedright\arraybackslash}p{(\columnwidth - 4\tabcolsep) * \real{0.1766}}
  >{\raggedright\arraybackslash}p{(\columnwidth - 4\tabcolsep) * \real{0.8234}}@{}}
\caption{Requerimientos Funcionales - Ventas en Tienda.}\label{tab:a-1-1}\\
\toprule\noalign{}
\begin{minipage}[b]{\linewidth}\raggedright
\textbf{ID}
\end{minipage} & \begin{minipage}[b]{\linewidth}\raggedright
\textbf{Descripción del Requerimiento}
\end{minipage} \\
\midrule\noalign{}
\endfirsthead
\toprule\noalign{}
\begin{minipage}[b]{\linewidth}\raggedright
\textbf{ID}
\end{minipage} & \begin{minipage}[b]{\linewidth}\raggedright
\textbf{Descripción del Requerimiento}
\end{minipage} \\
\midrule\noalign{}
\endhead
\bottomrule\noalign{}
\endlastfoot
\textbf{RF-001} & El sistema deberá permitir la venta física de la unidad aunque exista una reserva temporal activa en el canal digital. \\
\end{longtable}

\subsubsection{A.1.2. Listado de Requerimientos Funcionales: Seguridad y Accesos}\label{a-1-2-listado-de-requerimientos-funcionales-seguridad-y-accesos}

A continuación, se detalla el listado de requerimientos funcionales (RF) del módulo/categoría ``Seguridad y Accesos''.

\renewcommand{\thetable}{A.1.2}
\begin{longtable}[]{@{}
  >{\raggedright\arraybackslash}p{(\columnwidth - 4\tabcolsep) * \real{0.0979}}
  >{\raggedright\arraybackslash}p{(\columnwidth - 4\tabcolsep) * \real{0.9021}}@{}}
\caption{Requerimientos Funcionales - Seguridad y Accesos.}\label{tab:a-1-2}\\
\toprule\noalign{}
\begin{minipage}[b]{\linewidth}\raggedright
\textbf{ID}
\end{minipage} & \begin{minipage}[b]{\linewidth}\raggedright
\textbf{Descripción del Requerimiento}
\end{minipage} \\
\midrule\noalign{}
\endfirsthead
\toprule\noalign{}
\begin{minipage}[b]{\linewidth}\raggedright
\textbf{ID}
\end{minipage} & \begin{minipage}[b]{\linewidth}\raggedright
\textbf{Descripción del Requerimiento}
\end{minipage} \\
\midrule\noalign{}
\endhead
\bottomrule\noalign{}
\endlastfoot
\textbf{RF-002} & El sistema deberá exigir credenciales individuales al usuario antes de habilitar cualquier acción. \\
\addlinespace[2pt]
\textbf{RF-003} & El sistema debe impedir el acceso anónimo a una terminal. \\
\addlinespace[2pt]
\textbf{RF-004} & El sistema debe impedir el acceso mediante credencial compartida. \\
\addlinespace[2pt]
\textbf{RF-005} & El sistema deberá permitir el traspaso de sesión nominativa entre usuarios en una terminal compartida. \\
\addlinespace[2pt]
\textbf{RF-006} & El sistema deberá cerrar automáticamente la sesión por inactividad en la terminal compartida. \\
\addlinespace[2pt]
\textbf{RF-007} & El sistema deberá impedir la ejecución de la función de originación a un usuario sin capacitación normativa acreditada y vigente. \\
\addlinespace[2pt]
\textbf{RF-008} & El sistema deberá impedir la ejecución de la función de repactación a un usuario sin capacitación normativa acreditada y vigente. \\
\addlinespace[2pt]
\textbf{RF-009} & El sistema deberá revocar automáticamente la habilitación funcional del usuario al vencer su acreditación de capacitación. \\
\addlinespace[2pt]
\textbf{RF-010} & El sistema deberá permitir al proveedor patrocinar la habilitación temporal de acceso de su repositor externo, con fecha de término declarada. \\
\addlinespace[2pt]
\textbf{RF-011} & El sistema deberá caducar automáticamente la habilitación de acceso del personal externo al cumplirse su vigencia. \\
\addlinespace[2pt]
\textbf{RF-012} & El sistema deberá registrar de forma individualizada el acceso y la actividad del personal externo. \\
\addlinespace[2pt]
\textbf{RF-013} & El sistema deberá revocar la totalidad de los accesos y credenciales del trabajador a partir del término efectivo de su vínculo. \\
\addlinespace[2pt]
\textbf{RF-014} & El sistema debe conciliar los accesos vigentes contra la nómina activa. \\
\addlinespace[2pt]
\textbf{RF-015} & El sistema debe informar al oficial de seguridad los accesos huérfanos detectados en la conciliación. \\
\end{longtable}

\subsubsection{A.1.3. Listado de Requerimientos Funcionales: Precios y Etiquetado}\label{a-1-3-listado-de-requerimientos-funcionales-precios-y-etiquetado}

A continuación, se detalla el listado de requerimientos funcionales (RF) del módulo/categoría ``Precios y Etiquetado''.

\renewcommand{\thetable}{A.1.3}
\begin{longtable}[]{@{}
  >{\raggedright\arraybackslash}p{(\columnwidth - 4\tabcolsep) * \real{0.0979}}
  >{\raggedright\arraybackslash}p{(\columnwidth - 4\tabcolsep) * \real{0.9021}}@{}}
\caption{Requerimientos Funcionales - Precios y Etiquetado.}\label{tab:a-1-3}\\
\toprule\noalign{}
\begin{minipage}[b]{\linewidth}\raggedright
\textbf{ID}
\end{minipage} & \begin{minipage}[b]{\linewidth}\raggedright
\textbf{Descripción del Requerimiento}
\end{minipage} \\
\midrule\noalign{}
\endfirsthead
\toprule\noalign{}
\begin{minipage}[b]{\linewidth}\raggedright
\textbf{ID}
\end{minipage} & \begin{minipage}[b]{\linewidth}\raggedright
\textbf{Descripción del Requerimiento}
\end{minipage} \\
\midrule\noalign{}
\endhead
\bottomrule\noalign{}
\endlastfoot
\textbf{RF-016} & El sistema deberá propagar el cambio de precio a las 380 líneas de caja. \\
\addlinespace[2pt]
\textbf{RF-017} & El sistema deberá propagar el cambio de precio al canal digital. \\
\addlinespace[2pt]
\textbf{RF-018} & El sistema deberá propagar el cambio de precio a los puntos de exhibición física de cada tienda afectada. \\
\addlinespace[2pt]
\textbf{RF-019} & El sistema deberá permitir al reponedor registrar la referencia y la tienda de la etiqueta cambiada, con el instante del cambio. \\
\addlinespace[2pt]
\textbf{RF-020} & El sistema deberá registrar la identidad individual del ejecutor del cambio de etiqueta. \\
\addlinespace[2pt]
\textbf{RF-021} & El sistema deberá permitir al ejecutivo de cumplimiento recuperar el precio publicado de una referencia para una fecha, hora y canal determinados. \\
\addlinespace[2pt]
\textbf{RF-022} & El sistema deberá obtener el precio vigente de la referencia según el último precio propagado a la etiqueta de esa tienda en la fecha de la venta. \\
\addlinespace[2pt]
\textbf{RF-023} & El sistema deberá cobrar el menor de ambos precios para el consumidor. \\
\addlinespace[2pt]
\textbf{RF-024} & El sistema deberá generar un registro de incidente de discrepancia de precio consultable. \\
\addlinespace[2pt]
\textbf{RF-027} & El sistema debe generar una alerta operativa al jefe de tienda identificando la etiqueta electrónica no confirmada. \\
\addlinespace[2pt]
\textbf{RF-028} & El sistema deberá impedir la venta de la referencia al precio nuevo mientras su punto de exhibición no confirme la actualización, salvo que el artículo se encuentre bloqueado para venta. \\
\addlinespace[2pt]
\textbf{RF-029} & El sistema deberá generar un indicador diario de puntos de exhibición desactualizados con identificación individual de cada punto. \\
\addlinespace[2pt]
\textbf{RF-030} & El sistema deberá agrupar los cambios de precio de sala en ventanas parametrizadas fuera del horario de atención. \\
\addlinespace[2pt]
\textbf{RF-032} & El sistema deberá evaluar la vigencia de cada promoción al instante de emisión del documento. \\
\addlinespace[2pt]
\textbf{RF-033} & El sistema deberá evaluar la aplicabilidad de la promoción según el canal de la venta. \\
\addlinespace[2pt]
\textbf{RF-034} & El sistema deberá evaluar la aplicabilidad de la promoción según la tienda en que se realiza la venta. \\
\end{longtable}

\subsubsection{A.1.4. Listado de Requerimientos Funcionales: Pedidos y Cumplimiento}\label{a-1-4-listado-de-requerimientos-funcionales-pedidos-y-cumplimiento}

A continuación, se detalla el listado de requerimientos funcionales (RF) del módulo/categoría ``Pedidos y Cumplimiento''.

\renewcommand{\thetable}{A.1.4}
\begin{longtable}[]{@{}
  >{\raggedright\arraybackslash}p{(\columnwidth - 4\tabcolsep) * \real{0.0979}}
  >{\raggedright\arraybackslash}p{(\columnwidth - 4\tabcolsep) * \real{0.9021}}@{}}
\caption{Requerimientos Funcionales - Pedidos y Cumplimiento.}\label{tab:a-1-4}\\
\toprule\noalign{}
\begin{minipage}[b]{\linewidth}\raggedright
\textbf{ID}
\end{minipage} & \begin{minipage}[b]{\linewidth}\raggedright
\textbf{Descripción del Requerimiento}
\end{minipage} \\
\midrule\noalign{}
\endfirsthead
\toprule\noalign{}
\begin{minipage}[b]{\linewidth}\raggedright
\textbf{ID}
\end{minipage} & \begin{minipage}[b]{\linewidth}\raggedright
\textbf{Descripción del Requerimiento}
\end{minipage} \\
\midrule\noalign{}
\endhead
\bottomrule\noalign{}
\endlastfoot
\textbf{RF-035} & El sistema deberá crear una reserva temporal para la unidad agregada al carro. \\
\addlinespace[2pt]
\textbf{RF-036} & El sistema deberá permitir parametrizar el período de vigencia de la reserva temporal. \\
\addlinespace[2pt]
\textbf{RF-037} & El sistema deberá mantener activa la reserva mientras el cliente mantenga actividad, conforme al período parametrizado. \\
\addlinespace[2pt]
\textbf{RF-038} & El sistema deberá impedir que una misma unidad sea comprometida simultáneamente en más de una operación de venta digital. \\
\addlinespace[2pt]
\textbf{RF-039} & El sistema debe liberar la reserva temporal al expirar su período de vigencia. \\
\addlinespace[2pt]
\textbf{RF-040} & El sistema debe devolver la unidad liberada al disponible para vender. \\
\addlinespace[2pt]
\textbf{RF-041} & El sistema deberá liberar de inmediato la reserva digital de esa unidad, sin esperar la expiración del período. \\
\addlinespace[2pt]
\textbf{RF-042} & El sistema deberá aplicar el período de vigencia de reserva reducido declarado como parámetro para el evento anual. \\
\addlinespace[2pt]
\textbf{RF-043} & El sistema deberá preautorizar el medio de pago al aceptar el pedido, sin capturar el cobro. \\
\addlinespace[2pt]
\textbf{RF-044} & El sistema deberá verificar la existencia física real de la unidad en el punto asignado antes de habilitar la captura del cobro. \\
\addlinespace[2pt]
\textbf{RF-045} & El sistema deberá capturar el cobro únicamente al registrarse el evento de confirmación de la preparación física. \\
\addlinespace[2pt]
\textbf{RF-046} & El sistema deberá determinar automáticamente la alternativa de resolución aplicable según el motor de reglas (cancelar, sustituir, derivar a tercero o entrega diferida). \\
\addlinespace[2pt]
\textbf{RF-047} & El sistema deberá requerir la aceptación explícita del cliente antes de capturar el cobro. \\
\addlinespace[2pt]
\textbf{RF-048} & El sistema deberá anular la preautorización de pago del pedido. \\
\addlinespace[2pt]
\textbf{RF-049} & El sistema deberá cancelar el pedido registrando el motivo de la cancelación. \\
\addlinespace[2pt]
\textbf{RF-050} & El sistema deberá generar la conciliación diaria de preautorizaciones vencidas sin captura. \\
\addlinespace[2pt]
\textbf{RF-051} & El sistema deberá notificar al cliente el cambio de estado de su pedido antes de efectuar cualquier cobro definitivo. \\
\addlinespace[2pt]
\textbf{RF-052} & El sistema deberá permitir a todos los actores consultar el estado del pedido desde una única fuente de verdad. \\
\addlinespace[2pt]
\textbf{RF-053} & El sistema deberá permitir marcar como no elegible para cumplimiento digital el stock de exhibición de las categorías declaradas. \\
\addlinespace[2pt]
\textbf{RF-054} & El sistema debe calcular la base de comisión reconociendo al vendedor y a la tienda de origen de la unidad. \\
\addlinespace[2pt]
\textbf{RF-055} & El sistema debe condicionar la transmisión de la base de comisión al movimiento real de inventario verificado en bodega. \\
\addlinespace[2pt]
\textbf{RF-056} & El sistema deberá reasignar el pedido a otro punto de despacho con existencia disponible. \\
\addlinespace[2pt]
\textbf{RF-057} & El sistema deberá impedir una segunda reasignación del mismo pedido. \\
\addlinespace[2pt]
\textbf{RF-058} & El sistema deberá acotar la reasignación a la ventana de tiempo parametrizada. \\
\addlinespace[2pt]
\textbf{RF-059} & El sistema deberá ofrecer al cliente un producto equivalente. \\
\addlinespace[2pt]
\textbf{RF-060} & El sistema deberá ofrecer al cliente la espera compensada con la compensación parametrizada. \\
\addlinespace[2pt]
\textbf{RF-061} & El sistema deberá ofrecer al cliente la liberación del pedido con anulación de la preautorización. \\
\addlinespace[2pt]
\textbf{RF-062} & El sistema debe monitorear el tiempo restante de cada pedido respecto de su fecha prometida de entrega. \\
\addlinespace[2pt]
\textbf{RF-063} & El sistema debe priorizar la preparación del pedido al alcanzar el umbral parametrizado de tiempo restante. \\
\addlinespace[2pt]
\textbf{RF-064} & El sistema debe permitir al cliente no autenticado consultar el precio vigente de una referencia. \\
\addlinespace[2pt]
\textbf{RF-065} & El sistema debe permitir al cliente no autenticado consultar la disponibilidad de la referencia por tienda. \\
\addlinespace[2pt]
\textbf{RF-066} & El sistema debe permitir al cliente no autenticado consultar la disponibilidad de la referencia para despacho. \\
\addlinespace[2pt]
\textbf{RF-067} & El sistema debe permitir al cliente autenticado consultar sus compras. \\
\addlinespace[2pt]
\textbf{RF-068} & El sistema debe permitir al cliente autenticado consultar sus devoluciones. \\
\addlinespace[2pt]
\textbf{RF-069} & El sistema debe permitir al cliente autenticado consultar su estado de cuenta. \\
\addlinespace[2pt]
\textbf{RF-070} & El sistema debe permitir al cliente autenticado consultar sus documentos. \\
\addlinespace[2pt]
\textbf{RF-071} & El sistema deberá calcular la base de comisión considerando el canal de origen y el canal de cumplimiento cuando estos difieran. \\
\addlinespace[2pt]
\textbf{RF-072} & El sistema debe permitir al ejecutivo de mesón consultar el estado real del pedido. \\
\addlinespace[2pt]
\textbf{RF-073} & El sistema debe permitir al ejecutivo de mesón consultar el precio aplicado al pedido. \\
\addlinespace[2pt]
\textbf{RF-074} & El sistema debe permitir al ejecutivo de mesón ofrecer al cliente las opciones de resolución definidas en RF-059, RF-060 y RF-061. \\
\addlinespace[2pt]
\textbf{RF-075} & El sistema deberá excluir al centro de distribución de Concepción como origen de promesa de entrega en el canal digital mientras no esté acreditado su sistema de gestión de almacenes. \\
\addlinespace[2pt]
\textbf{RF-076} & El sistema deberá calcular el costo total de servir de cada punto de despacho candidato, incorporando costo logístico, costo de oportunidad de sala y plazo comprometido. \\
\addlinespace[2pt]
\textbf{RF-077} & El sistema deberá seleccionar como punto de despacho aquel de menor costo total de servir. \\
\addlinespace[2pt]
\textbf{RF-078} & El sistema deberá registrar el valor de cada término del costo total de servir que fundamento la selección. \\
\addlinespace[2pt]
\textbf{RF-079} & El sistema deberá calcular la fecha prometida de entrega en función del punto de despacho, la capacidad de preparación y el transportista asignado. \\
\addlinespace[2pt]
\textbf{RF-080} & El sistema deberá impedir la publicación de un plazo fijo de catálogo como fecha prometida de entrega. \\
\addlinespace[2pt]
\textbf{RF-081} & El sistema deberá calcular el cumplimiento de la promesa de entrega por pedido individual. \\
\end{longtable}

\subsubsection{A.1.5. Listado de Requerimientos Funcionales: Operación de Tienda}\label{a-1-5-listado-de-requerimientos-funcionales-operacion-de-tienda}

A continuación, se detalla el listado de requerimientos funcionales (RF) del módulo/categoría ``Operación de Tienda''.

\renewcommand{\thetable}{A.1.5}
\begin{longtable}[]{@{}
  >{\raggedright\arraybackslash}p{(\columnwidth - 4\tabcolsep) * \real{0.0979}}
  >{\raggedright\arraybackslash}p{(\columnwidth - 4\tabcolsep) * \real{0.9021}}@{}}
\caption{Requerimientos Funcionales - Operación de Tienda.}\label{tab:a-1-5}\\
\toprule\noalign{}
\begin{minipage}[b]{\linewidth}\raggedright
\textbf{ID}
\end{minipage} & \begin{minipage}[b]{\linewidth}\raggedright
\textbf{Descripción del Requerimiento}
\end{minipage} \\
\midrule\noalign{}
\endfirsthead
\toprule\noalign{}
\begin{minipage}[b]{\linewidth}\raggedright
\textbf{ID}
\end{minipage} & \begin{minipage}[b]{\linewidth}\raggedright
\textbf{Descripción del Requerimiento}
\end{minipage} \\
\midrule\noalign{}
\endhead
\bottomrule\noalign{}
\endlastfoot
\textbf{RF-082} & El sistema deberá detectar la pérdida del enlace externo y conmutar automáticamente la tienda a modo desconectado. \\
\addlinespace[2pt]
\textbf{RF-083} & El sistema deberá permitir al cajero registrar una venta en modo desconectado. \\
\addlinespace[2pt]
\textbf{RF-084} & El sistema deberá permitir al cajero cobrar la venta en modo desconectado. \\
\addlinespace[2pt]
\textbf{RF-085} & El sistema deberá aplicar las promociones vigentes en modo desconectado, evaluadas según RF-032, RF-033 y RF-034. \\
\addlinespace[2pt]
\textbf{RF-086} & El sistema deberá emitir el documento de venta en contingencia utilizando folios previamente asignados por el sistema de gestión empresarial. \\
\addlinespace[2pt]
\textbf{RF-087} & El sistema deberá permitir al vendedor consultar la existencia local de la tienda en modo desconectado. \\
\addlinespace[2pt]
\textbf{RF-088} & El sistema deberá detectar el restablecimiento del enlace y conmutar la tienda a modo conectado. \\
\addlinespace[2pt]
\textbf{RF-089} & El sistema deberá permitir el otorgamiento de crédito en modo desconectado exclusivamente contra cupo preaprobado vigente almacenado en caché local. \\
\addlinespace[2pt]
\textbf{RF-090} & El sistema deberá aplicar el tope de monto parametrizado a cada operación de crédito cursada en modo desconectado. \\
\addlinespace[2pt]
\textbf{RF-091} & El sistema deberá aplicar el tope parametrizado de número de operaciones de crédito en modo desconectado. \\
\addlinespace[2pt]
\textbf{RF-092} & El sistema deberá impedir la apertura de una tarjeta nueva en modo desconectado. \\
\addlinespace[2pt]
\textbf{RF-093} & El sistema deberá impedir la ampliación de cupo en modo desconectado. \\
\addlinespace[2pt]
\textbf{RF-094} & El sistema deberá marcar la operación como cursada en modo desconectado para su posterior validación. \\
\addlinespace[2pt]
\textbf{RF-095} & El sistema debe reconciliar hacia los sistemas centrales la totalidad de las ventas registradas en modo desconectado. \\
\addlinespace[2pt]
\textbf{RF-096} & El sistema debe reconciliar hacia los sistemas centrales la totalidad de los documentos emitidos en modo desconectado. \\
\addlinespace[2pt]
\textbf{RF-097} & El sistema deberá procesar la reconciliación de forma idempotente, impidiendo la duplicación de ventas o documentos ante reintentos. \\
\addlinespace[2pt]
\textbf{RF-098} & El sistema deberá resolver el conflicto de existencia comprometida aplicando la regla de precedencia declarada, produciendo el mismo resultado ante repetición. \\
\addlinespace[2pt]
\textbf{RF-099} & El sistema deberá generar un informe de excepciones de la reconciliación, identificando cada conflicto resuelto y su regla aplicada. \\
\addlinespace[2pt]
\textbf{RF-100} & El sistema deberá enrutar la emisión de todo documento tributario hacia el sistema de gestión empresarial como único emisor. \\
\addlinespace[2pt]
\textbf{RF-101} & El sistema deberá permitir al vendedor de piso consultar la disponibilidad para vender de una referencia. \\
\end{longtable}

\subsubsection{A.1.6. Listado de Requerimientos Funcionales: Marketplace}\label{a-1-6-listado-de-requerimientos-funcionales-marketplace}

A continuación, se detalla el listado de requerimientos funcionales (RF) del módulo/categoría ``Marketplace''.

\renewcommand{\thetable}{A.1.6}
\begin{longtable}[]{@{}
  >{\raggedright\arraybackslash}p{(\columnwidth - 4\tabcolsep) * \real{0.0979}}
  >{\raggedright\arraybackslash}p{(\columnwidth - 4\tabcolsep) * \real{0.9021}}@{}}
\caption{Requerimientos Funcionales - Marketplace.}\label{tab:a-1-6}\\
\toprule\noalign{}
\begin{minipage}[b]{\linewidth}\raggedright
\textbf{ID}
\end{minipage} & \begin{minipage}[b]{\linewidth}\raggedright
\textbf{Descripción del Requerimiento}
\end{minipage} \\
\midrule\noalign{}
\endfirsthead
\toprule\noalign{}
\begin{minipage}[b]{\linewidth}\raggedright
\textbf{ID}
\end{minipage} & \begin{minipage}[b]{\linewidth}\raggedright
\textbf{Descripción del Requerimiento}
\end{minipage} \\
\midrule\noalign{}
\endhead
\bottomrule\noalign{}
\endlastfoot
\textbf{RF-102} & El sistema debe permitir al vendedor de marketplace declarar la existencia de sus referencias. \\
\addlinespace[2pt]
\textbf{RF-103} & El sistema debe permitir al vendedor de marketplace actualizar la existencia previamente declarada. \\
\addlinespace[2pt]
\textbf{RF-104} & El sistema deberá registrar cada actualización de existencia declarada con su fecha y hora. \\
\addlinespace[2pt]
\textbf{RF-105} & El sistema deberá calcular los indicadores de nivel de servicio por vendedor externo según las reglas publicadas. \\
\addlinespace[2pt]
\textbf{RF-106} & El sistema deberá notificar al vendedor de marketplace la recepción de la devolución en el momento en que se registra. \\
\addlinespace[2pt]
\textbf{RF-107} & El sistema deberá permitir al ejecutivo registrar que parte asume el costo de la devolución. \\
\addlinespace[2pt]
\textbf{RF-108} & El sistema deberá permitir al ejecutivo registrar la prestación entregada al cliente (reparación, reposición o devolución del precio). \\
\addlinespace[2pt]
\textbf{RF-109} & El sistema deberá publicar únicamente la existencia declarada en la última actualización vigente del vendedor. \\
\addlinespace[2pt]
\textbf{RF-110} & El sistema debe permitir al vendedor de marketplace consultar el estado de cada pedido intermediado. \\
\addlinespace[2pt]
\textbf{RF-111} & El sistema debe permitir al vendedor de marketplace consultar las devoluciones de sus pedidos. \\
\addlinespace[2pt]
\textbf{RF-112} & El sistema debe permitir al vendedor de marketplace consultar su evaluación de desempeño vigente. \\
\addlinespace[2pt]
\textbf{RF-113} & El sistema debe identificar visiblemente al vendedor de la unidad en el catálogo, la ficha de producto y el proceso de compra. \\
\addlinespace[2pt]
\textbf{RF-114} & El sistema debe identificar visiblemente a quien despacha la unidad en el catálogo, la ficha de producto y el proceso de compra. \\
\addlinespace[2pt]
\textbf{RF-115} & El sistema debe mostrar las condiciones de devolución aplicables a la unidad. \\
\addlinespace[2pt]
\textbf{RF-116} & El sistema debe mostrar las condiciones de garantía legal aplicables a la unidad. \\
\addlinespace[2pt]
\textbf{RF-117} & El sistema deberá impedir que un pedido intermediado comprometa existencia propia de la compañía. \\
\addlinespace[2pt]
\textbf{RF-118} & El sistema deberá impedir que una unidad de inventario propio sea asignada al cumplimiento de un pedido de marketplace. \\
\addlinespace[2pt]
\textbf{RF-119} & El sistema deberá registrar el acuse de conocimiento de las reglas de evaluación por parte de cada vendedor externo. \\
\addlinespace[2pt]
\textbf{RF-120} & El sistema deberá impedir la aplicación de una regla de evaluación a un vendedor que no tenga acuse de conocimiento previo de esa regla. \\
\addlinespace[2pt]
\textbf{RF-121} & El sistema deberá determinar la consecuencia escalonada aplicable según la matriz de escalamiento parametrizada (advertencia, restricción de publicación, suspensión). \\
\addlinespace[2pt]
\textbf{RF-122} & El sistema deberá requerir la ejecución de la sanción por el rol nominado en la matriz de escalamiento. \\
\addlinespace[2pt]
\textbf{RF-123} & El sistema deberá registrar la trazabilidad de la sanción aplicada, indicando regla incumplida, consecuencia, rol ejecutor e instante. \\
\addlinespace[2pt]
\textbf{RF-124} & El sistema deberá despublicar automáticamente la oferta cuyo stock declarado haya superado el plazo de vigencia sin actualización. \\
\addlinespace[2pt]
\textbf{RF-125} & El sistema deberá calcular e informar al sistema de gestión empresarial la base de comisión de marketplace sobre la venta efectivamente cumplida \\
\addlinespace[2pt]
\textbf{RF-126} & El sistema deberá notificar al sistema de gestión empresarial la anulación de la base de comisión calculada cuando ocurra una devolución o cancelación. \\
\end{longtable}

\subsubsection{A.1.7. Listado de Requerimientos Funcionales: Inventario y Disponibilidad}\label{a-1-7-listado-de-requerimientos-funcionales-inventario-y-disponibilidad}

A continuación, se detalla el listado de requerimientos funcionales (RF) del módulo/categoría ``Inventario y Disponibilidad''.

\renewcommand{\thetable}{A.1.7}
\begin{longtable}[]{@{}
  >{\raggedright\arraybackslash}p{(\columnwidth - 4\tabcolsep) * \real{0.0979}}
  >{\raggedright\arraybackslash}p{(\columnwidth - 4\tabcolsep) * \real{0.9021}}@{}}
\caption{Requerimientos Funcionales - Inventario y Disponibilidad.}\label{tab:a-1-7}\\
\toprule\noalign{}
\begin{minipage}[b]{\linewidth}\raggedright
\textbf{ID}
\end{minipage} & \begin{minipage}[b]{\linewidth}\raggedright
\textbf{Descripción del Requerimiento}
\end{minipage} \\
\midrule\noalign{}
\endfirsthead
\toprule\noalign{}
\begin{minipage}[b]{\linewidth}\raggedright
\textbf{ID}
\end{minipage} & \begin{minipage}[b]{\linewidth}\raggedright
\textbf{Descripción del Requerimiento}
\end{minipage} \\
\midrule\noalign{}
\endhead
\bottomrule\noalign{}
\endlastfoot
\textbf{RF-127} & El sistema deberá calcular la existencia disponible para vender restando de la existencia registrada las reservas vigentes, el comprometido no despachado y el colchón de confianza aplicable. \\
\addlinespace[2pt]
\textbf{RF-128} & El sistema deberá determinar el valor del colchón de confianza aplicable en función de la categoría del artículo. \\
\addlinespace[2pt]
\textbf{RF-129} & El sistema deberá determinar el valor del colchón de confianza aplicable en función del punto de existencia. \\
\addlinespace[2pt]
\textbf{RF-130} & El sistema deberá registrar la traza del cálculo del disponible, incluyendo el valor de cada término de la fórmula y la versión de parámetros aplicada. \\
\addlinespace[2pt]
\textbf{RF-131} & El sistema deberá permitir parametrizar la frecuencia de conteo cíclico por categoría. \\
\addlinespace[2pt]
\textbf{RF-132} & El sistema deberá permitir parametrizar el método de conteo por categoría. \\
\addlinespace[2pt]
\textbf{RF-133} & El sistema deberá permitir parametrizar el criterio de gatillo de recuento extraordinario por categoría. \\
\addlinespace[2pt]
\textbf{RF-134} & El sistema deberá generar la programación de conteos cíclicos según la frecuencia parametrizada de cada categoría. \\
\addlinespace[2pt]
\textbf{RF-135} & El sistema debe permitir registrar el resultado del conteo cíclico. \\
\addlinespace[2pt]
\textbf{RF-136} & El sistema deberá señalar automáticamente los pasillos o referencias que requieren recuento extraordinario. \\
\addlinespace[2pt]
\textbf{RF-137} & El sistema deberá calcular la exactitud de inventario resultante por categoría. \\
\addlinespace[2pt]
\textbf{RF-138} & El sistema deberá permitir clasificar cada diferencia detectada en uno de los componentes de merma definidos (pérdida física, daño no dado de baja, error de recepción, unidad mal ubicada, devolución mal reintegrada, error de digitación). \\
\addlinespace[2pt]
\textbf{RF-139} & El sistema deberá impedir el cierre de un ajuste de inventario que no tenga asignado un componente de merma. \\
\addlinespace[2pt]
\textbf{RF-140} & El sistema deberá cuantificar el monto y la proporción de cada componente de merma sobre el total del período. \\
\addlinespace[2pt]
\textbf{RF-141} & El sistema deberá generar el informe mensual de merma cuya suma de componentes sea igual al total registrado en el período. \\
\addlinespace[2pt]
\textbf{RF-142} & El sistema deberá exigir al analista comercial completar los atributos obligatorios de la referencia (categoría, dimensiones, fragilidad, volumen, restricción de despacho) antes de guardar el registro como apto para uso operativo. \\
\addlinespace[2pt]
\textbf{RF-143} & El sistema deberá impedir la publicación de esa referencia en el canal digital mientras no cuente con los atributos obligatorios completos. \\
\addlinespace[2pt]
\textbf{RF-144} & El sistema deberá generar automáticamente la propuesta diaria de reposición utilizando el disponible para vender calculado en RF-127. \\
\addlinespace[2pt]
\textbf{RF-145} & El sistema deberá permitir al planificador de logística ajustar la propuesta de reposición antes de confirmar el envío. \\
\addlinespace[2pt]
\textbf{RF-146} & El sistema deberá registrar el ajuste realizado indicando usuario, valor propuesto, valor confirmado e instante. \\
\addlinespace[2pt]
\textbf{RF-147} & El sistema debe generar un reporte periódico de referencias con atributos obligatorios incompletos. \\
\addlinespace[2pt]
\textbf{RF-148} & El sistema debe generar un reporte periódico de publicaciones fallidas indicando la causa del error. \\
\addlinespace[2pt]
\textbf{RF-149} & El sistema debe permitir al gerente de logística definir el valor del colchón de confianza por categoría. \\
\addlinespace[2pt]
\textbf{RF-150} & El sistema debe permitir al gerente de logística definir el valor del colchón de confianza por punto de existencia. \\
\addlinespace[2pt]
\textbf{RF-151} & El sistema deberá registrar quien definió o modificó el colchón, cuando, el valor anterior y el valor nuevo. \\
\addlinespace[2pt]
\textbf{RF-152} & El sistema deberá sugerir un valor de referencia del colchón calculado a partir del histórico de ventas y quiebres de la categoría. \\
\addlinespace[2pt]
\textbf{RF-153} & El sistema deberá permitir al vendedor de piso registrar las unidades que ingresan al probador. \\
\addlinespace[2pt]
\textbf{RF-154} & El sistema deberá excluir del disponible para vender las unidades registradas en probador. \\
\addlinespace[2pt]
\textbf{RF-155} & El sistema debe permitir al vendedor de piso confirmar el reingreso a sala de la unidad registrada en probador. \\
\addlinespace[2pt]
\textbf{RF-156} & El sistema debe permitir al vendedor de piso registrar la venta de la unidad registrada en probador. \\
\addlinespace[2pt]
\textbf{RF-157} & El sistema deberá impedir que cualquier canal de venta consuma el saldo bruto de inventario para publicar o comprometer existencia. \\
\addlinespace[2pt]
\textbf{RF-158} & El sistema debe calcular el porcentaje de error probable de la referencia en función de su categoría. \\
\addlinespace[2pt]
\textbf{RF-159} & El sistema debe calcular el porcentaje de error probable de la referencia en función de su punto de existencia. \\
\addlinespace[2pt]
\textbf{RF-160} & El sistema deberá mostrar al cliente el porcentaje de error probable calculado en RF-158 y RF-159 junto a la disponibilidad publicada. \\
\addlinespace[2pt]
\textbf{RF-161} & El sistema deberá mostrar al vendedor de piso el porcentaje de error probable calculado en RF-158 y RF-159. \\
\addlinespace[2pt]
\textbf{RF-162} & El sistema deberá generar una alerta interna dirigida al rol facultado cuando la exactitud de una categoría caiga bajo el umbral definido. \\
\addlinespace[2pt]
\textbf{RF-163} & El sistema deberá permitir exclusivamente al rol facultado suspender manualmente la publicación de una categoría. \\
\addlinespace[2pt]
\textbf{RF-164} & El sistema deberá abstenerse de suspender automáticamente la publicación de una categoría que no tenga una regla de degradación declarada y aprobada previamente. \\
\end{longtable}

\subsubsection{A.1.8. Listado de Requerimientos Funcionales: Gobernanza de Datos}\label{a-1-8-listado-de-requerimientos-funcionales-gobernanza-de-datos}

A continuación, se detalla el listado de requerimientos funcionales (RF) del módulo/categoría ``Gobernanza de Datos''.

\renewcommand{\thetable}{A.1.8}
\begin{longtable}[]{@{}
  >{\raggedright\arraybackslash}p{(\columnwidth - 4\tabcolsep) * \real{0.0979}}
  >{\raggedright\arraybackslash}p{(\columnwidth - 4\tabcolsep) * \real{0.9021}}@{}}
\caption{Requerimientos Funcionales - Gobernanza de Datos.}\label{tab:a-1-8}\\
\toprule\noalign{}
\begin{minipage}[b]{\linewidth}\raggedright
\textbf{ID}
\end{minipage} & \begin{minipage}[b]{\linewidth}\raggedright
\textbf{Descripción del Requerimiento}
\end{minipage} \\
\midrule\noalign{}
\endfirsthead
\toprule\noalign{}
\begin{minipage}[b]{\linewidth}\raggedright
\textbf{ID}
\end{minipage} & \begin{minipage}[b]{\linewidth}\raggedright
\textbf{Descripción del Requerimiento}
\end{minipage} \\
\midrule\noalign{}
\endhead
\bottomrule\noalign{}
\endlastfoot
\textbf{RF-165} & El sistema deberá registrar cada cruce de información ejecutado entre ámbitos, indicando el dato cruzado, la finalidad, la base de licitud, la autorización nominada y el instante. \\
\addlinespace[2pt]
\textbf{RF-166} & El sistema deberá bloquear todo intento de cruce de información que no corresponda a una interfaz declarada en el inventario de flujos autorizados. \\
\addlinespace[2pt]
\textbf{RF-167} & El sistema deberá registrar el intento de cruce bloqueado con el componente origen, el dato solicitado y el instante. \\
\addlinespace[2pt]
\textbf{RF-168} & El sistema debe permitir al oficial de cumplimiento mantener el inventario de interfaces de cruce declaradas. \\
\addlinespace[2pt]
\textbf{RF-169} & El sistema debe impedir el registro de una interfaz de cruce que no declare finalidad, base de licitud y autorización nominada. \\
\addlinespace[2pt]
\textbf{RF-170} & El sistema deberá excluir del catálogo de atributos disponibles en el motor de campañas todo atributo de origen financiero (mora, deuda, cupo utilizado, comportamiento de pago). \\
\addlinespace[2pt]
\textbf{RF-171} & El sistema deberá rechazar la ejecución de toda facilidad comercial que invoque un atributo de origen financiero sin un registro de cruce autorizado asociado. \\
\addlinespace[2pt]
\textbf{RF-172} & El sistema deberá rechazar la ejecución de todo proceso crediticio que invoque un atributo de origen retail sin un registro de cruce autorizado asociado. \\
\addlinespace[2pt]
\textbf{RF-173} & El sistema deberá asignar a la misma persona un identificador distinto en el ámbito retail y en el ámbito fiscalizado. \\
\addlinespace[2pt]
\textbf{RF-174} & El sistema deberá resolver la correspondencia entre identificadores exclusivamente a través de una tabla de correspondencia custodiada, con acceso nominado y registrado. \\
\addlinespace[2pt]
\textbf{RF-175} & El sistema deberá rechazar la persistencia de la clave del ámbito financiero como columna o atributo en cualquier entidad del ámbito retail. \\
\addlinespace[2pt]
\textbf{RF-176} & El sistema deberá exigir el registro de la evaluación de impacto sobre la frontera de datos antes de habilitar el estado 'aprobada' de una iniciativa. \\
\end{longtable}

\subsubsection{A.1.9. Listado de Requerimientos Funcionales: Evento de Alta Concurrencia}\label{a-1-9-listado-de-requerimientos-funcionales-evento-de-alta-concurrencia}

A continuación, se detalla el listado de requerimientos funcionales (RF) del módulo/categoría ``Evento de Alta Concurrencia''.

\renewcommand{\thetable}{A.1.9}
\begin{longtable}[]{@{}
  >{\raggedright\arraybackslash}p{(\columnwidth - 4\tabcolsep) * \real{0.0979}}
  >{\raggedright\arraybackslash}p{(\columnwidth - 4\tabcolsep) * \real{0.9021}}@{}}
\caption{Requerimientos Funcionales - Evento de Alta Concurrencia.}\label{tab:a-1-9}\\
\toprule\noalign{}
\begin{minipage}[b]{\linewidth}\raggedright
\textbf{ID}
\end{minipage} & \begin{minipage}[b]{\linewidth}\raggedright
\textbf{Descripción del Requerimiento}
\end{minipage} \\
\midrule\noalign{}
\endfirsthead
\toprule\noalign{}
\begin{minipage}[b]{\linewidth}\raggedright
\textbf{ID}
\end{minipage} & \begin{minipage}[b]{\linewidth}\raggedright
\textbf{Descripción del Requerimiento}
\end{minipage} \\
\midrule\noalign{}
\endhead
\bottomrule\noalign{}
\endlastfoot
\textbf{RF-177} & El sistema debe permitir declarar previamente el orden de degradación de los servicios. \\
\addlinespace[2pt]
\textbf{RF-178} & El sistema debe permitir parametrizar los criterios que activan cada nivel de degradación declarado. \\
\addlinespace[2pt]
\textbf{RF-179} & El sistema deberá suspender la publicación de la categoría indicada por el orden de degradación declarado. \\
\addlinespace[2pt]
\textbf{RF-180} & El sistema deberá reducir el límite de unidades por cliente al valor declarado. \\
\addlinespace[2pt]
\textbf{RF-181} & El sistema deberá desactivar temporalmente los medios de pago de mayor fricción operativa declarados. \\
\addlinespace[2pt]
\textbf{RF-182} & El sistema deberá monitorear en tiempo real la tasa de cancelación por falta de existencia por categoría. \\
\addlinespace[2pt]
\textbf{RF-183} & El sistema deberá activar la acción de degradación declarada para esa categoría en RF-026.1. \\
\addlinespace[2pt]
\textbf{RF-184} & El sistema deberá permitir parametrizar las cinco ventanas de congelamiento (1 nov - 6 ene; evento anual y su semana previa; evento de noviembre; semana del Día de la Madre; última semana de enero a primera de marzo). \\
\addlinespace[2pt]
\textbf{RF-185} & El sistema deberá impedir la ejecución del despliegue o intervención en producción durante la ventana de congelamiento. \\
\addlinespace[2pt]
\textbf{RF-186} & El sistema deberá registrar el intento de despliegue bloqueado, indicando componente, solicitante e instante. \\
\end{longtable}

\subsubsection{A.1.10. Listado de Requerimientos Funcionales: Devoluciones y Garantía Legal}\label{a-1-10-listado-de-requerimientos-funcionales-devoluciones-y-garantia-legal}

A continuación, se detalla el listado de requerimientos funcionales (RF) del módulo/categoría ``Devoluciones y Garantía Legal''.

\renewcommand{\thetable}{A.1.10}
\begin{longtable}[]{@{}
  >{\raggedright\arraybackslash}p{(\columnwidth - 4\tabcolsep) * \real{0.0979}}
  >{\raggedright\arraybackslash}p{(\columnwidth - 4\tabcolsep) * \real{0.9021}}@{}}
\caption{Requerimientos Funcionales - Devoluciones y Garantía Legal.}\label{tab:a-1-10}\\
\toprule\noalign{}
\begin{minipage}[b]{\linewidth}\raggedright
\textbf{ID}
\end{minipage} & \begin{minipage}[b]{\linewidth}\raggedright
\textbf{Descripción del Requerimiento}
\end{minipage} \\
\midrule\noalign{}
\endfirsthead
\toprule\noalign{}
\begin{minipage}[b]{\linewidth}\raggedright
\textbf{ID}
\end{minipage} & \begin{minipage}[b]{\linewidth}\raggedright
\textbf{Descripción del Requerimiento}
\end{minipage} \\
\midrule\noalign{}
\endhead
\bottomrule\noalign{}
\endlastfoot
\textbf{RF-187} & El sistema debe permitir al ejecutivo de mesón atender íntegramente el caso de garantía legal en la tienda. \\
\addlinespace[2pt]
\textbf{RF-188} & El sistema debe impedir que el flujo de atención exija la derivación del consumidor al fabricante, al servicio técnico o al vendedor de marketplace como condición para ser atendido. \\
\addlinespace[2pt]
\textbf{RF-189} & El sistema deberá permitir al ejecutivo registrar el motivo de la devolución (talla, color, expectativa, falla). \\
\addlinespace[2pt]
\textbf{RF-190} & El sistema deberá impedir el reingreso de una unidad devuelta al inventario disponible mientras no exista una decisión de aptitud registrada. \\
\addlinespace[2pt]
\textbf{RF-191} & El sistema deberá permitir al ejecutivo registrar la decisión de aptitud de la unidad devuelta, identificando al responsable de la decisión y su instante. \\
\addlinespace[2pt]
\textbf{RF-192} & El sistema deberá presentar al ejecutivo las tres opciones de garantía legal (reparación, reposición y devolución del precio) para que el consumidor elija. \\
\addlinespace[2pt]
\textbf{RF-193} & El sistema debe registrar la opción de garantía legal ofrecida al consumidor. \\
\addlinespace[2pt]
\textbf{RF-194} & El sistema debe registrar la opción de garantía legal efectivamente elegida por el consumidor. \\
\addlinespace[2pt]
\textbf{RF-195} & El sistema deberá permitir parametrizar el plazo de garantía legal por tipo de producto, sin requerir modificación de código. \\
\addlinespace[2pt]
\textbf{RF-196} & El sistema deberá registrar el hito de resolución al consumidor con su fecha propia. \\
\addlinespace[2pt]
\textbf{RF-197} & El sistema deberá registrar el hito de recuperación contra el tercero responsable con una fecha independiente de la resolución al consumidor. \\
\addlinespace[2pt]
\textbf{RF-198} & El sistema deberá impedir que el estado de la recuperación contra el tercero condicione el cierre de la resolución al consumidor. \\
\end{longtable}

\subsubsection{A.1.11. Listado de Requerimientos Funcionales: Crédito y Cobranza}\label{a-1-11-listado-de-requerimientos-funcionales-credito-y-cobranza}

A continuación, se detalla el listado de requerimientos funcionales (RF) del módulo/categoría ``Crédito y Cobranza''.

\renewcommand{\thetable}{A.1.11}
\begin{longtable}[]{@{}
  >{\raggedright\arraybackslash}p{(\columnwidth - 4\tabcolsep) * \real{0.0979}}
  >{\raggedright\arraybackslash}p{(\columnwidth - 4\tabcolsep) * \real{0.9021}}@{}}
\caption{Requerimientos Funcionales - Crédito y Cobranza.}\label{tab:a-1-11}\\
\toprule\noalign{}
\begin{minipage}[b]{\linewidth}\raggedright
\textbf{ID}
\end{minipage} & \begin{minipage}[b]{\linewidth}\raggedright
\textbf{Descripción del Requerimiento}
\end{minipage} \\
\midrule\noalign{}
\endfirsthead
\toprule\noalign{}
\begin{minipage}[b]{\linewidth}\raggedright
\textbf{ID}
\end{minipage} & \begin{minipage}[b]{\linewidth}\raggedright
\textbf{Descripción del Requerimiento}
\end{minipage} \\
\midrule\noalign{}
\endhead
\bottomrule\noalign{}
\endlastfoot
\textbf{RF-199} & El sistema deberá exigir la entrega completa de la información precontractual antes de habilitar la evaluación de la solicitud de crédito. \\
\addlinespace[2pt]
\textbf{RF-200} & El sistema deberá registrar la versión del documento precontractual entregado. \\
\addlinespace[2pt]
\textbf{RF-201} & El sistema deberá registrar el instante de entrega de la información precontractual. \\
\addlinespace[2pt]
\textbf{RF-202} & El sistema deberá registrar el instante de aceptación del cliente. \\
\addlinespace[2pt]
\textbf{RF-203} & El sistema deberá registrar el medio por el cual se entrego y se aceptó la información precontractual. \\
\addlinespace[2pt]
\textbf{RF-204} & El sistema deberá registrar el contenido exacto aceptado por el cliente o su huella de integridad verificable. \\
\addlinespace[2pt]
\textbf{RF-205} & El sistema deberá dejar constancia expresa de que el cliente recibió la información antes de aceptar, mediante un mecanismo distinto de la sola firma en papel archivado. \\
\addlinespace[2pt]
\textbf{RF-206} & El sistema deberá impedir el registro de la aceptación del crédito mientras no exista acreditación de entrega previa de la información precontractual completa. \\
\addlinespace[2pt]
\textbf{RF-207} & El sistema deberá registrar la evidencia del consentimiento expreso e informado del cliente mediante un mecanismo verificable de integridad. \\
\addlinespace[2pt]
\textbf{RF-208} & El sistema deberá impedir el registro de una modificación de condiciones del crédito que no tenga evidencia de consentimiento asociada. \\
\addlinespace[2pt]
\textbf{RF-209} & El sistema deberá reconstruir el acto de consentimiento presentando que se informó, en qué versión y que aceptó el cliente. \\
\addlinespace[2pt]
\textbf{RF-210} & El sistema deberá restaurar desde archivo frío los antecedentes de una operación de crédito de cualquier cohorte dentro del plazo de conservación. \\
\addlinespace[2pt]
\textbf{RF-211} & El sistema debe permitir al ejecutivo de cobranza iniciar una gestión de cobranza. \\
\addlinespace[2pt]
\textbf{RF-212} & El sistema debe registrar la gestión de cobranza ejecutada con su medio, su instante y su ejecutor. \\
\addlinespace[2pt]
\textbf{RF-213} & El sistema debe impedir la ejecución de una gestión de cobranza fuera de los límites normativos de horario y de medio de contacto. \\
\addlinespace[2pt]
\textbf{RF-214} & El sistema debe enlazar el expediente de la repactación con la gestión de cobranza que la originó. \\
\addlinespace[2pt]
\textbf{RF-215} & El sistema debe enlazar el expediente de la repactación con la evidencia de consentimiento registrada en RF-207. \\
\addlinespace[2pt]
\textbf{RF-216} & El sistema deberá generar un reporte de conciliación diaria de saldos durante todo el proceso de migración. \\
\addlinespace[2pt]
\textbf{RF-217} & El sistema deberá permitir al cliente no autenticado consultar la información precontractual del crédito. \\
\addlinespace[2pt]
\textbf{RF-218} & El sistema deberá permitir al cliente no autenticado utilizar un simulador de costo total del crédito. \\
\addlinespace[2pt]
\textbf{RF-219} & El sistema deberá mantener cargada la Tasa Máxima Convencional vigente por tipo y tramo de operación, con su fecha de vigencia. \\
\addlinespace[2pt]
\textbf{RF-220} & El sistema deberá impedir la originación de una operación cuya tasa supere la Tasa Máxima Convencional vigente a la fecha de la operación. \\
\addlinespace[2pt]
\textbf{RF-221} & El sistema deberá determinar el cupo exclusivamente a partir de las variables de la evaluación de capacidad de pago documentada. \\
\addlinespace[2pt]
\textbf{RF-222} & El sistema deberá excluir el monto de la venta en curso del conjunto de variables de asignación de cupo. \\
\addlinespace[2pt]
\textbf{RF-223} & El sistema deberá excluir la solicitud del vendedor del conjunto de variables de asignación de cupo. \\
\addlinespace[2pt]
\textbf{RF-224} & El sistema deberá impedir al vendedor alterar el resultado de una evaluación crediticia. \\
\addlinespace[2pt]
\textbf{RF-225} & El sistema deberá impedir la repetición de la evaluación crediticia del mismo cliente dentro de la ventana de enfriamiento parametrizada. \\
\addlinespace[2pt]
\textbf{RF-226} & El sistema deberá registrar cada intento de evaluación con la identidad del solicitante y el resultado obtenido. \\
\end{longtable}

\subsection{A.2. Listado de Requerimientos No Funcionales (RNF)}\label{a-2-listado-de-requerimientos-no-funcionales-rnf}

A continuación, se detalla el listado completo de requerimientos no funcionales (RNF) del Catálogo Atómico, agrupado por categoría, incluyendo su umbral/métrica y método de verificación. Total: 75 requerimientos.

\subsubsection{A.2.1. Listado de Requerimientos No Funcionales: Disponibilidad}\label{a-2-1-listado-de-requerimientos-no-funcionales-disponibilidad}

A continuación, se detalla el listado de requerimientos no funcionales (RNF) de la categoría ``Disponibilidad''.

\renewcommand{\thetable}{A.2.1}
\begin{longtable}[]{@{}
  >{\raggedright\arraybackslash}p{(\columnwidth - 8\tabcolsep) * \real{0.1027}}
  >{\raggedright\arraybackslash}p{(\columnwidth - 8\tabcolsep) * \real{0.4302}}
  >{\raggedright\arraybackslash}p{(\columnwidth - 8\tabcolsep) * \real{0.2087}}
  >{\raggedright\arraybackslash}p{(\columnwidth - 8\tabcolsep) * \real{0.2584}}@{}}
\caption{Requerimientos No Funcionales - Disponibilidad.}\label{tab:a-2-1}\\
\toprule\noalign{}
\begin{minipage}[b]{\linewidth}\raggedright
\textbf{ID}
\end{minipage} & \begin{minipage}[b]{\linewidth}\raggedright
\textbf{Descripción del Requerimiento}
\end{minipage} & \begin{minipage}[b]{\linewidth}\raggedright
\textbf{Umbral / Métrica}
\end{minipage} & \begin{minipage}[b]{\linewidth}\raggedright
\textbf{Método de verificación}
\end{minipage} \\
\midrule\noalign{}
\endfirsthead
\toprule\noalign{}
\begin{minipage}[b]{\linewidth}\raggedright
\textbf{ID}
\end{minipage} & \begin{minipage}[b]{\linewidth}\raggedright
\textbf{Descripción del Requerimiento}
\end{minipage} & \begin{minipage}[b]{\linewidth}\raggedright
\textbf{Umbral / Métrica}
\end{minipage} & \begin{minipage}[b]{\linewidth}\raggedright
\textbf{Método de verificación}
\end{minipage} \\
\midrule\noalign{}
\endhead
\bottomrule\noalign{}
\endlastfoot
\textbf{RNF-02} & La interfaz de declaración de existencia del vendedor de marketplace debe estar disponible de forma continua. & 24x7x365; disponibilidad \textgreater{}= 99,5 \% (supuesto fundamentado). & Medición mensual de disponibilidad de la interfaz de sincronización de catálogo. \\
\addlinespace[2pt]
\textbf{RNF-24} & Una tienda debe poder operar sin enlace externo durante el período mínimo definido. & ≥ 24 horas continuas en tienda & Prueba de desconexión controlada en tienda. \\
\addlinespace[2pt]
\textbf{RNF-25} & El centro de distribución principal debe poder operar sin enlace externo durante el período mínimo definido. & \textgreater{}= 4 horas continuas. & Prueba de desconexión controlada en CD. \\
\addlinespace[2pt]
\textbf{RNF-26} & La sincronización tras la reconexión no debe superar el plazo definido. & \textless{}= 30 minutos tras 8 horas de desconexión. & Prueba de reconexión cronometrada tras desconexión programada de 8 horas. \\
\addlinespace[2pt]
\textbf{RNF-27} & La sincronización tras la reconexión no debe producir pérdida de ventas ni de documentos. & 0 ventas y 0 documentos perdidos o duplicados. & Cuadratura del universo de transacciones de la ventana desconectada contra el sistema central. \\
\addlinespace[2pt]
\textbf{RNF-28} & La operación desconectada autónoma on-premise debe sostenerse durante el período mínimo exigido por las bases técnicas. & \textgreater{}= 24 horas (bases técnicas transversales), superior al mínimo de 8 h de RN-46. & Prueba de desconexión extendida. \\
\addlinespace[2pt]
\textbf{RNF-29} & La operación de tiendas y centros de distribución debe estar disponible en la ventana horaria definida. & 09:00 a 23:00, todos los días del año. & Monitoreo continuo de disponibilidad por servicio. \\
\addlinespace[2pt]
\textbf{RNF-30} & El canal digital debe estar disponible de forma continua. & 24x7x365; disponibilidad \textgreater{}= 99,9 \% & Medición mensual de disponibilidad contra el nivel de servicio comprometido. \\
\addlinespace[2pt]
\textbf{RNF-31} & Los servicios financieros que afectan pagos, estados de cuenta y bloqueo de tarjetas deben estar disponibles de forma continua. & 24x7x365; disponibilidad \textgreater{}= 99,9 \% & Medición mensual segregada por ámbito, informada a la filial emisora. \\
\end{longtable}

\subsubsection{A.2.2. Listado de Requerimientos No Funcionales: Desempeño}\label{a-2-2-listado-de-requerimientos-no-funcionales-desempeno}

A continuación, se detalla el listado de requerimientos no funcionales (RNF) de la categoría ``Desempeño''.

\renewcommand{\thetable}{A.2.2}
\begin{longtable}[]{@{}
  >{\raggedright\arraybackslash}p{(\columnwidth - 8\tabcolsep) * \real{0.1027}}
  >{\raggedright\arraybackslash}p{(\columnwidth - 8\tabcolsep) * \real{0.4302}}
  >{\raggedright\arraybackslash}p{(\columnwidth - 8\tabcolsep) * \real{0.2087}}
  >{\raggedright\arraybackslash}p{(\columnwidth - 8\tabcolsep) * \real{0.2584}}@{}}
\caption{Requerimientos No Funcionales - Desempeño.}\label{tab:a-2-2}\\
\toprule\noalign{}
\begin{minipage}[b]{\linewidth}\raggedright
\textbf{ID}
\end{minipage} & \begin{minipage}[b]{\linewidth}\raggedright
\textbf{Descripción del Requerimiento}
\end{minipage} & \begin{minipage}[b]{\linewidth}\raggedright
\textbf{Umbral / Métrica}
\end{minipage} & \begin{minipage}[b]{\linewidth}\raggedright
\textbf{Método de verificación}
\end{minipage} \\
\midrule\noalign{}
\endfirsthead
\toprule\noalign{}
\begin{minipage}[b]{\linewidth}\raggedright
\textbf{ID}
\end{minipage} & \begin{minipage}[b]{\linewidth}\raggedright
\textbf{Descripción del Requerimiento}
\end{minipage} & \begin{minipage}[b]{\linewidth}\raggedright
\textbf{Umbral / Métrica}
\end{minipage} & \begin{minipage}[b]{\linewidth}\raggedright
\textbf{Método de verificación}
\end{minipage} \\
\midrule\noalign{}
\endhead
\bottomrule\noalign{}
\endlastfoot
\textbf{RNF-03} & El registro del conteo cíclico no debe bloquear ni degradar la operación de venta de la sala. & 0 bloqueos de la línea de caja durante el conteo; degradación del tiempo de venta \textless{}= 5 \% respecto de la línea base de RNF-18. & Prueba de conteo cíclico ejecutado en horario de venta con medición simultánea del tiempo de transacción en caja. \\
\addlinespace[2pt]
\textbf{RNF-05} & La actualización del estado de un pedido debe reflejarse en todos los puntos de consulta dentro del plazo máximo definido. & Desfase máximo entre cambio real de estado y su reflejo en cualquier canal: \textless{}= 30 segundos & Auditoría cruzada de estado entre los cuatro canales (alineado con Subdocumento 3, Tabla 3.26). \\
\addlinespace[2pt]
\textbf{RNF-06} & El tiempo de evaluación de una solicitud de crédito en el punto de venta físico no debe exceder el umbral objetivo. & \textless{}= 8 segundos. & Prueba de tiempo de respuesta en punto de venta bajo carga representativa. \\
\addlinespace[2pt]
\textbf{RNF-16} & La consulta de disponibilidad en la ficha de producto del canal digital debe responder dentro del umbral definido. & \textless{}= 400 ms. & Prueba de tiempo de respuesta bajo carga representativa. \\
\addlinespace[2pt]
\textbf{RNF-17} & La confirmación de un pedido durante el evento anual debe completarse dentro del umbral definido. & \textless{}= 3 s. & Prueba de carga con volumetria del evento anual. \\
\addlinespace[2pt]
\textbf{RNF-18} & La venta completa en caja con medio de pago externo debe completarse dentro del umbral definido. & \textless{}= 25 s. & Prueba de tiempo de respuesta en linea de caja. \\
\addlinespace[2pt]
\textbf{RNF-19} & La propagación de un cambio de precio a las 380 líneas de caja y al canal digital debe completarse dentro del umbral definido. & \textless{}= 5 minutos. & Medicion de latencia de propagación por destino. \\
\addlinespace[2pt]
\textbf{RNF-20} & El registro de una devolución en el mesón debe completarse dentro del umbral definido. & \textless{}= 60 s. & Prueba de tiempo de respuesta en mesón de atención. \\
\addlinespace[2pt]
\textbf{RNF-21} & La consulta de disponibilidad desde una terminal compartida del piso de venta debe responder dentro del umbral definido. & \textless{}= 2 s. & Prueba de tiempo de respuesta en terminal de sala. \\
\addlinespace[2pt]
\textbf{RNF-58} & La recuperación del precio publicado de una referencia en una fecha, hora y canal arbitrarios debe completarse dentro del plazo objetivo. & \textless{}= 1 minuto, sobre una ventana de 3 años. & Consulta de auditoría sobre fecha/hora/canal aleatorios. \\
\end{longtable}

\subsubsection{A.2.3. Listado de Requerimientos No Funcionales: Consistencia de datos}\label{a-2-3-listado-de-requerimientos-no-funcionales-consistencia-de-datos}

A continuación, se detalla el listado de requerimientos no funcionales (RNF) de la categoría ``Consistencia de datos''.

\renewcommand{\thetable}{A.2.3}
\begin{longtable}[]{@{}
  >{\raggedright\arraybackslash}p{(\columnwidth - 8\tabcolsep) * \real{0.1027}}
  >{\raggedright\arraybackslash}p{(\columnwidth - 8\tabcolsep) * \real{0.4302}}
  >{\raggedright\arraybackslash}p{(\columnwidth - 8\tabcolsep) * \real{0.2087}}
  >{\raggedright\arraybackslash}p{(\columnwidth - 8\tabcolsep) * \real{0.2584}}@{}}
\caption{Requerimientos No Funcionales - Consistencia de datos.}\label{tab:a-2-3}\\
\toprule\noalign{}
\begin{minipage}[b]{\linewidth}\raggedright
\textbf{ID}
\end{minipage} & \begin{minipage}[b]{\linewidth}\raggedright
\textbf{Descripción del Requerimiento}
\end{minipage} & \begin{minipage}[b]{\linewidth}\raggedright
\textbf{Umbral / Métrica}
\end{minipage} & \begin{minipage}[b]{\linewidth}\raggedright
\textbf{Método de verificación}
\end{minipage} \\
\midrule\noalign{}
\endfirsthead
\toprule\noalign{}
\begin{minipage}[b]{\linewidth}\raggedright
\textbf{ID}
\end{minipage} & \begin{minipage}[b]{\linewidth}\raggedright
\textbf{Descripción del Requerimiento}
\end{minipage} & \begin{minipage}[b]{\linewidth}\raggedright
\textbf{Umbral / Métrica}
\end{minipage} & \begin{minipage}[b]{\linewidth}\raggedright
\textbf{Método de verificación}
\end{minipage} \\
\midrule\noalign{}
\endhead
\bottomrule\noalign{}
\endlastfoot
\textbf{RNF-04} & El estado de un pedido no debe presentar discrepancias entre canales de consulta. & 0\% de discrepancias en auditoría cruzada. & Prueba de sincronización cruzada entre canales tras cada cambio de estado. \\
\end{longtable}

\subsubsection{A.2.4. Listado de Requerimientos No Funcionales: Desempeño / Cumplimiento}\label{a-2-4-listado-de-requerimientos-no-funcionales-desempeno-cumplimiento}

A continuación, se detalla el listado de requerimientos no funcionales (RNF) de la categoría ``Desempeño / Cumplimiento''.

\renewcommand{\thetable}{A.2.4}
\begin{longtable}[]{@{}
  >{\raggedright\arraybackslash}p{(\columnwidth - 8\tabcolsep) * \real{0.1027}}
  >{\raggedright\arraybackslash}p{(\columnwidth - 8\tabcolsep) * \real{0.4302}}
  >{\raggedright\arraybackslash}p{(\columnwidth - 8\tabcolsep) * \real{0.2087}}
  >{\raggedright\arraybackslash}p{(\columnwidth - 8\tabcolsep) * \real{0.2584}}@{}}
\caption{Requerimientos No Funcionales - Desempeño / Cumplimiento.}\label{tab:a-2-4}\\
\toprule\noalign{}
\begin{minipage}[b]{\linewidth}\raggedright
\textbf{ID}
\end{minipage} & \begin{minipage}[b]{\linewidth}\raggedright
\textbf{Descripción del Requerimiento}
\end{minipage} & \begin{minipage}[b]{\linewidth}\raggedright
\textbf{Umbral / Métrica}
\end{minipage} & \begin{minipage}[b]{\linewidth}\raggedright
\textbf{Método de verificación}
\end{minipage} \\
\midrule\noalign{}
\endfirsthead
\toprule\noalign{}
\begin{minipage}[b]{\linewidth}\raggedright
\textbf{ID}
\end{minipage} & \begin{minipage}[b]{\linewidth}\raggedright
\textbf{Descripción del Requerimiento}
\end{minipage} & \begin{minipage}[b]{\linewidth}\raggedright
\textbf{Umbral / Métrica}
\end{minipage} & \begin{minipage}[b]{\linewidth}\raggedright
\textbf{Método de verificación}
\end{minipage} \\
\midrule\noalign{}
\endhead
\bottomrule\noalign{}
\endlastfoot
\textbf{RNF-07} & El tiempo total de originación de crédito en el mesón, con información precontractual acreditada, no debe superar el tiempo del proceso actual. & \textless{}= 3 minutos (tiempo actual declarado en el caso). & Medicion de tiempo de originación extremo a extremo en mesón, con y sin acreditación. \\
\end{longtable}

\subsubsection{A.2.5. Listado de Requerimientos No Funcionales: Portabilidad / Migración}\label{a-2-5-listado-de-requerimientos-no-funcionales-portabilidad-migracion}

A continuación, se detalla el listado de requerimientos no funcionales (RNF) de la categoría ``Portabilidad / Migración''.

\renewcommand{\thetable}{A.2.5}
\begin{longtable}[]{@{}
  >{\raggedright\arraybackslash}p{(\columnwidth - 8\tabcolsep) * \real{0.1027}}
  >{\raggedright\arraybackslash}p{(\columnwidth - 8\tabcolsep) * \real{0.4302}}
  >{\raggedright\arraybackslash}p{(\columnwidth - 8\tabcolsep) * \real{0.2087}}
  >{\raggedright\arraybackslash}p{(\columnwidth - 8\tabcolsep) * \real{0.2584}}@{}}
\caption{Requerimientos No Funcionales - Portabilidad / Migración.}\label{tab:a-2-5}\\
\toprule\noalign{}
\begin{minipage}[b]{\linewidth}\raggedright
\textbf{ID}
\end{minipage} & \begin{minipage}[b]{\linewidth}\raggedright
\textbf{Descripción del Requerimiento}
\end{minipage} & \begin{minipage}[b]{\linewidth}\raggedright
\textbf{Umbral / Métrica}
\end{minipage} & \begin{minipage}[b]{\linewidth}\raggedright
\textbf{Método de verificación}
\end{minipage} \\
\midrule\noalign{}
\endfirsthead
\toprule\noalign{}
\begin{minipage}[b]{\linewidth}\raggedright
\textbf{ID}
\end{minipage} & \begin{minipage}[b]{\linewidth}\raggedright
\textbf{Descripción del Requerimiento}
\end{minipage} & \begin{minipage}[b]{\linewidth}\raggedright
\textbf{Umbral / Métrica}
\end{minipage} & \begin{minipage}[b]{\linewidth}\raggedright
\textbf{Método de verificación}
\end{minipage} \\
\midrule\noalign{}
\endhead
\bottomrule\noalign{}
\endlastfoot
\textbf{RNF-08} & La migración de la cartera de crédito viva no debe generar pérdida de datos. & 0 registros perdidos respecto del universo de 620.000 clientes con saldo. & Conteo y cuadratura de universos origen y destino en cada corte de migración. \\
\addlinespace[2pt]
\textbf{RNF-09} & La migración de la cartera de crédito viva no debe interrumpir el servicio de cobro. & 0 interrupciones del servicio de cobro durante la ventana de migración. & Monitoreo continuo de disponibilidad del servicio de cobro durante cada etapa de coexistencia. \\
\addlinespace[2pt]
\textbf{RNF-10} & La migración de la cartera de crédito viva no debe producir divergencia de saldos. & 0 divergencias no conciliadas en la conciliación diaria; migración completa antes de 2029. & Conciliación diaria automatizada de saldos entre entorno origen y destino (RF-216). \\
\end{longtable}

\subsubsection{A.2.6. Listado de Requerimientos No Funcionales: Seguridad / Arquitectura}\label{a-2-6-listado-de-requerimientos-no-funcionales-seguridad-arquitectura}

A continuación, se detalla el listado de requerimientos no funcionales (RNF) de la categoría ``Seguridad / Arquitectura''.

\renewcommand{\thetable}{A.2.6}
\begin{longtable}[]{@{}
  >{\raggedright\arraybackslash}p{(\columnwidth - 8\tabcolsep) * \real{0.1027}}
  >{\raggedright\arraybackslash}p{(\columnwidth - 8\tabcolsep) * \real{0.4302}}
  >{\raggedright\arraybackslash}p{(\columnwidth - 8\tabcolsep) * \real{0.2087}}
  >{\raggedright\arraybackslash}p{(\columnwidth - 8\tabcolsep) * \real{0.2584}}@{}}
\caption{Requerimientos No Funcionales - Seguridad / Arquitectura.}\label{tab:a-2-6}\\
\toprule\noalign{}
\begin{minipage}[b]{\linewidth}\raggedright
\textbf{ID}
\end{minipage} & \begin{minipage}[b]{\linewidth}\raggedright
\textbf{Descripción del Requerimiento}
\end{minipage} & \begin{minipage}[b]{\linewidth}\raggedright
\textbf{Umbral / Métrica}
\end{minipage} & \begin{minipage}[b]{\linewidth}\raggedright
\textbf{Método de verificación}
\end{minipage} \\
\midrule\noalign{}
\endfirsthead
\toprule\noalign{}
\begin{minipage}[b]{\linewidth}\raggedright
\textbf{ID}
\end{minipage} & \begin{minipage}[b]{\linewidth}\raggedright
\textbf{Descripción del Requerimiento}
\end{minipage} & \begin{minipage}[b]{\linewidth}\raggedright
\textbf{Umbral / Métrica}
\end{minipage} & \begin{minipage}[b]{\linewidth}\raggedright
\textbf{Método de verificación}
\end{minipage} \\
\midrule\noalign{}
\endhead
\bottomrule\noalign{}
\endlastfoot
\textbf{RNF-11} & La separación lógica entre los datos del retail y los del negocio financiero debe estar implementada. & Separación lógica operativa en todos los componentes que tratan datos de ambos ámbitos. & Revisión de arquitectura desplegada y prueba de acceso denegado desde el ámbito retail. \\
\addlinespace[2pt]
\textbf{RNF-12} & La separación lógica entre ambos ámbitos debe estar documentada. & Documento de separación vigente y versionado, aprobado por la Contraloría interna. & Revisión documental contra el diseño desplegado. \\
\addlinespace[2pt]
\textbf{RNF-13} & La separación lógica entre ambos ámbitos debe ser verificada por auditoría independiente. & Informe de auditoría aprobado con 0 hallazgos críticos abiertos, previo a cualquier vista unificada. & Auditoría independiente con informe formal y seguimiento de hallazgos. \\
\addlinespace[2pt]
\textbf{RNF-14} & La separación física entre la infraestructura del retail y la de la filial emisora debe estar implementada y acreditada. & Informe técnico de separación física acreditado; prueba de penetración que demuestre que un componente del retail no alcanza datos de crédito. & Prueba de penetración e inspección de infraestructura. \\
\end{longtable}

\subsubsection{A.2.7. Listado de Requerimientos No Funcionales: Auditabilidad}\label{a-2-7-listado-de-requerimientos-no-funcionales-auditabilidad}

A continuación, se detalla el listado de requerimientos no funcionales (RNF) de la categoría ``Auditabilidad''.

\renewcommand{\thetable}{A.2.7}
\begin{longtable}[]{@{}
  >{\raggedright\arraybackslash}p{(\columnwidth - 8\tabcolsep) * \real{0.1027}}
  >{\raggedright\arraybackslash}p{(\columnwidth - 8\tabcolsep) * \real{0.4302}}
  >{\raggedright\arraybackslash}p{(\columnwidth - 8\tabcolsep) * \real{0.2087}}
  >{\raggedright\arraybackslash}p{(\columnwidth - 8\tabcolsep) * \real{0.2584}}@{}}
\caption{Requerimientos No Funcionales - Auditabilidad.}\label{tab:a-2-7}\\
\toprule\noalign{}
\begin{minipage}[b]{\linewidth}\raggedright
\textbf{ID}
\end{minipage} & \begin{minipage}[b]{\linewidth}\raggedright
\textbf{Descripción del Requerimiento}
\end{minipage} & \begin{minipage}[b]{\linewidth}\raggedright
\textbf{Umbral / Métrica}
\end{minipage} & \begin{minipage}[b]{\linewidth}\raggedright
\textbf{Método de verificación}
\end{minipage} \\
\midrule\noalign{}
\endfirsthead
\toprule\noalign{}
\begin{minipage}[b]{\linewidth}\raggedright
\textbf{ID}
\end{minipage} & \begin{minipage}[b]{\linewidth}\raggedright
\textbf{Descripción del Requerimiento}
\end{minipage} & \begin{minipage}[b]{\linewidth}\raggedright
\textbf{Umbral / Métrica}
\end{minipage} & \begin{minipage}[b]{\linewidth}\raggedright
\textbf{Método de verificación}
\end{minipage} \\
\midrule\noalign{}
\endhead
\bottomrule\noalign{}
\endlastfoot
\textbf{RNF-15} & Toda acción relevante ejecutada en el sistema debe quedar registrada indicando usuario individual, acción e instante. & 100\% de acciones relevantes con registro de auditoría. & Auditoría de logs sobre muestra de transacciones. \\
\addlinespace[2pt]
\textbf{RNF-38} & Debe registrarse todo acceso a la cartera de crédito, al comportamiento de pago y a los antecedentes de evaluación crediticia. & 100\% de accesos a información sensible con registro (usuario, dato, finalidad, instante). & Auditoría de logs de acceso sobre muestra representativa. \\
\end{longtable}

\subsubsection{A.2.8. Listado de Requerimientos No Funcionales: Desempeño / Capacidad}\label{a-2-8-listado-de-requerimientos-no-funcionales-desempeno-capacidad}

A continuación, se detalla el listado de requerimientos no funcionales (RNF) de la categoría ``Desempeño / Capacidad''.

\renewcommand{\thetable}{A.2.8}
\begin{longtable}[]{@{}
  >{\raggedright\arraybackslash}p{(\columnwidth - 8\tabcolsep) * \real{0.1027}}
  >{\raggedright\arraybackslash}p{(\columnwidth - 8\tabcolsep) * \real{0.4302}}
  >{\raggedright\arraybackslash}p{(\columnwidth - 8\tabcolsep) * \real{0.2087}}
  >{\raggedright\arraybackslash}p{(\columnwidth - 8\tabcolsep) * \real{0.2584}}@{}}
\caption{Requerimientos No Funcionales - Desempeño / Capacidad.}\label{tab:a-2-8}\\
\toprule\noalign{}
\begin{minipage}[b]{\linewidth}\raggedright
\textbf{ID}
\end{minipage} & \begin{minipage}[b]{\linewidth}\raggedright
\textbf{Descripción del Requerimiento}
\end{minipage} & \begin{minipage}[b]{\linewidth}\raggedright
\textbf{Umbral / Métrica}
\end{minipage} & \begin{minipage}[b]{\linewidth}\raggedright
\textbf{Método de verificación}
\end{minipage} \\
\midrule\noalign{}
\endfirsthead
\toprule\noalign{}
\begin{minipage}[b]{\linewidth}\raggedright
\textbf{ID}
\end{minipage} & \begin{minipage}[b]{\linewidth}\raggedright
\textbf{Descripción del Requerimiento}
\end{minipage} & \begin{minipage}[b]{\linewidth}\raggedright
\textbf{Umbral / Métrica}
\end{minipage} & \begin{minipage}[b]{\linewidth}\raggedright
\textbf{Método de verificación}
\end{minipage} \\
\midrule\noalign{}
\endhead
\bottomrule\noalign{}
\endlastfoot
\textbf{RNF-22} & La plataforma debe soportar la concurrencia y el volumen del peak digital del evento anual sin degradar los umbrales de NFR-06.x. & Perfil declarado por el PROPONENTE sobre 104.000 pedidos en 3 días (proyección 150.000), equivalente a 22 días de venta en línea. & Prueba de carga con perfil del evento anual al 120 \% del peak proyectado. \\
\addlinespace[2pt]
\textbf{RNF-23} & La plataforma debe soportar la concurrencia y el volumen del peak presencial de la campaña de noviembre y diciembre sin degradar los umbrales de NFR-06.x. & Perfil declarado por el PROPONENTE sobre 380 líneas de caja (430 proyectadas) en peak de sábado de diciembre. & Prueba de carga con perfil presencial sostenido, independiente de la prueba del peak digital. \\
\end{longtable}

\subsubsection{A.2.9. Listado de Requerimientos No Funcionales: Recuperabilidad}\label{a-2-9-listado-de-requerimientos-no-funcionales-recuperabilidad}

A continuación, se detalla el listado de requerimientos no funcionales (RNF) de la categoría ``Recuperabilidad''.

\renewcommand{\thetable}{A.2.9}
\begin{longtable}[]{@{}
  >{\raggedright\arraybackslash}p{(\columnwidth - 8\tabcolsep) * \real{0.1027}}
  >{\raggedright\arraybackslash}p{(\columnwidth - 8\tabcolsep) * \real{0.4302}}
  >{\raggedright\arraybackslash}p{(\columnwidth - 8\tabcolsep) * \real{0.2087}}
  >{\raggedright\arraybackslash}p{(\columnwidth - 8\tabcolsep) * \real{0.2584}}@{}}
\caption{Requerimientos No Funcionales - Recuperabilidad.}\label{tab:a-2-9}\\
\toprule\noalign{}
\begin{minipage}[b]{\linewidth}\raggedright
\textbf{ID}
\end{minipage} & \begin{minipage}[b]{\linewidth}\raggedright
\textbf{Descripción del Requerimiento}
\end{minipage} & \begin{minipage}[b]{\linewidth}\raggedright
\textbf{Umbral / Métrica}
\end{minipage} & \begin{minipage}[b]{\linewidth}\raggedright
\textbf{Método de verificación}
\end{minipage} \\
\midrule\noalign{}
\endfirsthead
\toprule\noalign{}
\begin{minipage}[b]{\linewidth}\raggedright
\textbf{ID}
\end{minipage} & \begin{minipage}[b]{\linewidth}\raggedright
\textbf{Descripción del Requerimiento}
\end{minipage} & \begin{minipage}[b]{\linewidth}\raggedright
\textbf{Umbral / Métrica}
\end{minipage} & \begin{minipage}[b]{\linewidth}\raggedright
\textbf{Método de verificación}
\end{minipage} \\
\midrule\noalign{}
\endhead
\bottomrule\noalign{}
\endlastfoot
\textbf{RNF-32} & La solución debe cumplir los objetivos de recuperación ante desastres exigidos. & RTO \textless{}= 4 h; RPO \textless{}= 15 min. & Ensayo de recuperación ante desastres con medición de RTO y RPO efectivos. \\
\addlinespace[2pt]
\textbf{RNF-68} & La restauración desde archivo frío de una operación de crédito de la cohorte más antigua debe completarse dentro del plazo objetivo. & \textless{}= 24 horas (supuesto fundamentado). & Prueba de restauración desde archivo frío. \\
\end{longtable}

\subsubsection{A.2.10. Listado de Requerimientos No Funcionales: Seguridad}\label{a-2-10-listado-de-requerimientos-no-funcionales-seguridad}

A continuación, se detalla el listado de requerimientos no funcionales (RNF) de la categoría ``Seguridad''.

\renewcommand{\thetable}{A.2.10}
\begin{longtable}[]{@{}
  >{\raggedright\arraybackslash}p{(\columnwidth - 8\tabcolsep) * \real{0.1027}}
  >{\raggedright\arraybackslash}p{(\columnwidth - 8\tabcolsep) * \real{0.4302}}
  >{\raggedright\arraybackslash}p{(\columnwidth - 8\tabcolsep) * \real{0.2087}}
  >{\raggedright\arraybackslash}p{(\columnwidth - 8\tabcolsep) * \real{0.2584}}@{}}
\caption{Requerimientos No Funcionales - Seguridad.}\label{tab:a-2-10}\\
\toprule\noalign{}
\begin{minipage}[b]{\linewidth}\raggedright
\textbf{ID}
\end{minipage} & \begin{minipage}[b]{\linewidth}\raggedright
\textbf{Descripción del Requerimiento}
\end{minipage} & \begin{minipage}[b]{\linewidth}\raggedright
\textbf{Umbral / Métrica}
\end{minipage} & \begin{minipage}[b]{\linewidth}\raggedright
\textbf{Método de verificación}
\end{minipage} \\
\midrule\noalign{}
\endfirsthead
\toprule\noalign{}
\begin{minipage}[b]{\linewidth}\raggedright
\textbf{ID}
\end{minipage} & \begin{minipage}[b]{\linewidth}\raggedright
\textbf{Descripción del Requerimiento}
\end{minipage} & \begin{minipage}[b]{\linewidth}\raggedright
\textbf{Umbral / Métrica}
\end{minipage} & \begin{minipage}[b]{\linewidth}\raggedright
\textbf{Método de verificación}
\end{minipage} \\
\midrule\noalign{}
\endhead
\bottomrule\noalign{}
\endlastfoot
\textbf{RNF-33} & Los datos de identificación de clientes, la cartera de crédito, el comportamiento de pago y los antecedentes de evaluación crediticia deben cifrarse a nivel de campo. & 100\% de campos sensibles cifrados. & Auditoría de esquema de datos y prueba de exfiltración controlada. \\
\addlinespace[2pt]
\textbf{RNF-34} & Los datos de medios de pago no deben almacenarse en claro, sino tokenizarse. & 0\% de datos de medios de pago almacenados sin tokenizar. & Auditoría de esquema y prueba de exfiltración controlada. \\
\addlinespace[2pt]
\textbf{RNF-35} & La red de cajas, la administrativa, la de videovigilancia y la inalámbrica de clientes deben estar segmentadas en las 22 tiendas. & 22/22 tiendas con segmentación verificada (hoy 9). & Auditoría de configuración de red por tienda. \\
\addlinespace[2pt]
\textbf{RNF-36} & La red y el ámbito de sistemas de la filial emisora deben estar separados de forma acreditada respecto del retail. & Separación acreditada mediante informe técnico. & Auditoría de arquitectura de red. \\
\addlinespace[2pt]
\textbf{RNF-37} & Debe existir segregación de funciones verificable entre originación, aprobación, modificación de condiciones y cobranza. & 0 usuarios con permisos combinados de originación + aprobación o de modificación + cobranza sobre un mismo caso. & Auditoría de matriz de roles y permisos. \\
\addlinespace[2pt]
\textbf{RNF-59} & La revocación total de accesos de un trabajador debe completarse dentro del plazo máximo de la política interna. & \textless{}= 4 horas desde el término efectivo del vínculo (supuesto fundamentado). & Auditoría de accesos vigentes contra nómina activa. \\
\end{longtable}

\subsubsection{A.2.11. Listado de Requerimientos No Funcionales: Seguridad / Cumplimiento}\label{a-2-11-listado-de-requerimientos-no-funcionales-seguridad-cumplimiento}

A continuación, se detalla el listado de requerimientos no funcionales (RNF) de la categoría ``Seguridad / Cumplimiento''.

\renewcommand{\thetable}{A.2.11}
\begin{longtable}[]{@{}
  >{\raggedright\arraybackslash}p{(\columnwidth - 8\tabcolsep) * \real{0.1027}}
  >{\raggedright\arraybackslash}p{(\columnwidth - 8\tabcolsep) * \real{0.4302}}
  >{\raggedright\arraybackslash}p{(\columnwidth - 8\tabcolsep) * \real{0.2087}}
  >{\raggedright\arraybackslash}p{(\columnwidth - 8\tabcolsep) * \real{0.2584}}@{}}
\caption{Requerimientos No Funcionales - Seguridad / Cumplimiento.}\label{tab:a-2-11}\\
\toprule\noalign{}
\begin{minipage}[b]{\linewidth}\raggedright
\textbf{ID}
\end{minipage} & \begin{minipage}[b]{\linewidth}\raggedright
\textbf{Descripción del Requerimiento}
\end{minipage} & \begin{minipage}[b]{\linewidth}\raggedright
\textbf{Umbral / Métrica}
\end{minipage} & \begin{minipage}[b]{\linewidth}\raggedright
\textbf{Método de verificación}
\end{minipage} \\
\midrule\noalign{}
\endfirsthead
\toprule\noalign{}
\begin{minipage}[b]{\linewidth}\raggedright
\textbf{ID}
\end{minipage} & \begin{minipage}[b]{\linewidth}\raggedright
\textbf{Descripción del Requerimiento}
\end{minipage} & \begin{minipage}[b]{\linewidth}\raggedright
\textbf{Umbral / Métrica}
\end{minipage} & \begin{minipage}[b]{\linewidth}\raggedright
\textbf{Método de verificación}
\end{minipage} \\
\midrule\noalign{}
\endhead
\bottomrule\noalign{}
\endlastfoot
\textbf{RNF-39} & Los actos de crédito de apertura de tarjeta y de repactación deben contar con un mecanismo de consentimiento de integridad verificable. & 100 \% de aperturas y repactaciones con evidencia de integridad verificable, en la modalidad que la normativa admita para cada canal. & Auditoría censal del período y verificación de la huella de integridad por muestreo. \\
\addlinespace[2pt]
\textbf{RNF-40} & La conformidad de recepción de mercadería y la constancia de devolución deben contar con un mecanismo de consentimiento de integridad verificable. & 100 \% de conformidades y constancias con evidencia de integridad verificable. & Verificación de la validez de la firma y de la integridad del documento por muestreo. \\
\end{longtable}

\subsubsection{A.2.12. Listado de Requerimientos No Funcionales: Cumplimiento}\label{a-2-12-listado-de-requerimientos-no-funcionales-cumplimiento}

A continuación, se detalla el listado de requerimientos no funcionales (RNF) de la categoría ``Cumplimiento''.

\renewcommand{\thetable}{A.2.12}
\begin{longtable}[]{@{}
  >{\raggedright\arraybackslash}p{(\columnwidth - 8\tabcolsep) * \real{0.1027}}
  >{\raggedright\arraybackslash}p{(\columnwidth - 8\tabcolsep) * \real{0.4302}}
  >{\raggedright\arraybackslash}p{(\columnwidth - 8\tabcolsep) * \real{0.2087}}
  >{\raggedright\arraybackslash}p{(\columnwidth - 8\tabcolsep) * \real{0.2584}}@{}}
\caption{Requerimientos No Funcionales - Cumplimiento.}\label{tab:a-2-12}\\
\toprule\noalign{}
\begin{minipage}[b]{\linewidth}\raggedright
\textbf{ID}
\end{minipage} & \begin{minipage}[b]{\linewidth}\raggedright
\textbf{Descripción del Requerimiento}
\end{minipage} & \begin{minipage}[b]{\linewidth}\raggedright
\textbf{Umbral / Métrica}
\end{minipage} & \begin{minipage}[b]{\linewidth}\raggedright
\textbf{Método de verificación}
\end{minipage} \\
\midrule\noalign{}
\endfirsthead
\toprule\noalign{}
\begin{minipage}[b]{\linewidth}\raggedright
\textbf{ID}
\end{minipage} & \begin{minipage}[b]{\linewidth}\raggedright
\textbf{Descripción del Requerimiento}
\end{minipage} & \begin{minipage}[b]{\linewidth}\raggedright
\textbf{Umbral / Métrica}
\end{minipage} & \begin{minipage}[b]{\linewidth}\raggedright
\textbf{Método de verificación}
\end{minipage} \\
\midrule\noalign{}
\endhead
\bottomrule\noalign{}
\endlastfoot
\textbf{RNF-41} & Los documentos tributarios y antecedentes de venta deben conservarse por el plazo definido. & 6 años. & Auditoría de políticas de retención por tipo de dato. \\
\addlinespace[2pt]
\textbf{RNF-42} & Los antecedentes del crédito deben conservarse por el plazo del crédito más el período posterior exigido. & Plazo del crédito + 6 años (sustituye la práctica actual de 90 días para grabaciones). & Auditoría de retención y prueba de restauración desde archivo frío. \\
\addlinespace[2pt]
\textbf{RNF-43} & La evidencia de consentimiento de originación y repactación debe conservarse por el mismo plazo que el crédito. & Plazo del crédito + 6 años; recuperable a 10 años según RN-35. & Auditoría de retención y prueba de recuperación. \\
\addlinespace[2pt]
\textbf{RNF-44} & La trazabilidad del precio publicado debe conservarse por el plazo definido. & 3 años. & Auditoría de retención. \\
\addlinespace[2pt]
\textbf{RNF-45} & Los movimientos de inventario y conteos cíclicos deben conservarse por el plazo definido. & 6 años. & Auditoría de retención. \\
\addlinespace[2pt]
\textbf{RNF-46} & Los datos de fidelización deben conservarse por el plazo definido. & Duración de la relación + 2 años. & Auditoría de retención. \\
\addlinespace[2pt]
\textbf{RNF-47} & Los registros de videovigilancia deben conservarse por el plazo definido. & 30 días. & Auditoría de retención. \\
\addlinespace[2pt]
\textbf{RNF-69} & La atención de garantía legal en mesón no debe derivar al consumidor a un tercero como condición de atención, en ninguna tienda. & 0 derivaciones como condición de atención, medido con cliente oculto en las 22 tiendas. & Programa de cliente oculto en las 22 tiendas. \\
\addlinespace[2pt]
\textbf{RNF-74} & Todo tratamiento de datos personales debe declarar su finalidad. & 100 \% de los flujos con finalidad declarada. & Revisión del registro de actividades de tratamiento. \\
\addlinespace[2pt]
\textbf{RNF-75} & Todo tratamiento de datos personales debe invocar una de las bases de licitud de la Ley N° 21.719. & 100 \% de los flujos con base de licitud declarada; 0 tratamientos con consentimiento presunto. & Auditoría de conformidad por el oficial de cumplimiento. \\
\addlinespace[2pt]
\textbf{RNF-76} & La compañía debe mantener el registro de actividades de tratamiento de datos personales. & Registro completo y vigente, actualizado ante cada cambio de flujo. & Revisión documental periódica contra el inventario de interfaces de RF-168. \\
\end{longtable}

\subsubsection{A.2.13. Listado de Requerimientos No Funcionales: Portabilidad / Escalabilidad}\label{a-2-13-listado-de-requerimientos-no-funcionales-portabilidad-escalabilidad}

A continuación, se detalla el listado de requerimientos no funcionales (RNF) de la categoría ``Portabilidad / Escalabilidad''.

\renewcommand{\thetable}{A.2.13}
\begin{longtable}[]{@{}
  >{\raggedright\arraybackslash}p{(\columnwidth - 8\tabcolsep) * \real{0.1027}}
  >{\raggedright\arraybackslash}p{(\columnwidth - 8\tabcolsep) * \real{0.4302}}
  >{\raggedright\arraybackslash}p{(\columnwidth - 8\tabcolsep) * \real{0.2087}}
  >{\raggedright\arraybackslash}p{(\columnwidth - 8\tabcolsep) * \real{0.2584}}@{}}
\caption{Requerimientos No Funcionales - Portabilidad / Escalabilidad.}\label{tab:a-2-13}\\
\toprule\noalign{}
\begin{minipage}[b]{\linewidth}\raggedright
\textbf{ID}
\end{minipage} & \begin{minipage}[b]{\linewidth}\raggedright
\textbf{Descripción del Requerimiento}
\end{minipage} & \begin{minipage}[b]{\linewidth}\raggedright
\textbf{Umbral / Métrica}
\end{minipage} & \begin{minipage}[b]{\linewidth}\raggedright
\textbf{Método de verificación}
\end{minipage} \\
\midrule\noalign{}
\endfirsthead
\toprule\noalign{}
\begin{minipage}[b]{\linewidth}\raggedright
\textbf{ID}
\end{minipage} & \begin{minipage}[b]{\linewidth}\raggedright
\textbf{Descripción del Requerimiento}
\end{minipage} & \begin{minipage}[b]{\linewidth}\raggedright
\textbf{Umbral / Métrica}
\end{minipage} & \begin{minipage}[b]{\linewidth}\raggedright
\textbf{Método de verificación}
\end{minipage} \\
\midrule\noalign{}
\endhead
\bottomrule\noalign{}
\endlastfoot
\textbf{RNF-48} & La solución debe admitir la apertura de una tienda nueva por parametrización, sin desarrollo a medida. & Alta de tienda por configuración en \textless{}= 10 días hábiles (supuesto fundamentado, compatible con una apertura comercial). & Prueba de aprovisionamiento de una tienda simulada. \\
\end{longtable}

\subsubsection{A.2.14. Listado de Requerimientos No Funcionales: Usabilidad}\label{a-2-14-listado-de-requerimientos-no-funcionales-usabilidad}

A continuación, se detalla el listado de requerimientos no funcionales (RNF) de la categoría ``Usabilidad''.

\renewcommand{\thetable}{A.2.14}
\begin{longtable}[]{@{}
  >{\raggedright\arraybackslash}p{(\columnwidth - 8\tabcolsep) * \real{0.1027}}
  >{\raggedright\arraybackslash}p{(\columnwidth - 8\tabcolsep) * \real{0.4302}}
  >{\raggedright\arraybackslash}p{(\columnwidth - 8\tabcolsep) * \real{0.2087}}
  >{\raggedright\arraybackslash}p{(\columnwidth - 8\tabcolsep) * \real{0.2584}}@{}}
\caption{Requerimientos No Funcionales - Usabilidad.}\label{tab:a-2-14}\\
\toprule\noalign{}
\begin{minipage}[b]{\linewidth}\raggedright
\textbf{ID}
\end{minipage} & \begin{minipage}[b]{\linewidth}\raggedright
\textbf{Descripción del Requerimiento}
\end{minipage} & \begin{minipage}[b]{\linewidth}\raggedright
\textbf{Umbral / Métrica}
\end{minipage} & \begin{minipage}[b]{\linewidth}\raggedright
\textbf{Método de verificación}
\end{minipage} \\
\midrule\noalign{}
\endfirsthead
\toprule\noalign{}
\begin{minipage}[b]{\linewidth}\raggedright
\textbf{ID}
\end{minipage} & \begin{minipage}[b]{\linewidth}\raggedright
\textbf{Descripción del Requerimiento}
\end{minipage} & \begin{minipage}[b]{\linewidth}\raggedright
\textbf{Umbral / Métrica}
\end{minipage} & \begin{minipage}[b]{\linewidth}\raggedright
\textbf{Método de verificación}
\end{minipage} \\
\midrule\noalign{}
\endhead
\bottomrule\noalign{}
\endlastfoot
\textbf{RNF-49} & Las interfaces de sala, caja y mesón financiero deben ser operables por personal recién incorporado con capacitación mínima. & Tiempo de capacitación operativa básica \textless{}= 4 horas (supuesto fundamentado), medido con personal de temporada. & Prueba de usabilidad con personal recién incorporado en condiciones reales. \\
\end{longtable}

\subsubsection{A.2.15. Listado de Requerimientos No Funcionales: Usabilidad / Cumplimiento}\label{a-2-15-listado-de-requerimientos-no-funcionales-usabilidad-cumplimiento}

A continuación, se detalla el listado de requerimientos no funcionales (RNF) de la categoría ``Usabilidad / Cumplimiento''.

\renewcommand{\thetable}{A.2.15}
\begin{longtable}[]{@{}
  >{\raggedright\arraybackslash}p{(\columnwidth - 8\tabcolsep) * \real{0.1027}}
  >{\raggedright\arraybackslash}p{(\columnwidth - 8\tabcolsep) * \real{0.4302}}
  >{\raggedright\arraybackslash}p{(\columnwidth - 8\tabcolsep) * \real{0.2087}}
  >{\raggedright\arraybackslash}p{(\columnwidth - 8\tabcolsep) * \real{0.2584}}@{}}
\caption{Requerimientos No Funcionales - Usabilidad / Cumplimiento.}\label{tab:a-2-15}\\
\toprule\noalign{}
\begin{minipage}[b]{\linewidth}\raggedright
\textbf{ID}
\end{minipage} & \begin{minipage}[b]{\linewidth}\raggedright
\textbf{Descripción del Requerimiento}
\end{minipage} & \begin{minipage}[b]{\linewidth}\raggedright
\textbf{Umbral / Métrica}
\end{minipage} & \begin{minipage}[b]{\linewidth}\raggedright
\textbf{Método de verificación}
\end{minipage} \\
\midrule\noalign{}
\endfirsthead
\toprule\noalign{}
\begin{minipage}[b]{\linewidth}\raggedright
\textbf{ID}
\end{minipage} & \begin{minipage}[b]{\linewidth}\raggedright
\textbf{Descripción del Requerimiento}
\end{minipage} & \begin{minipage}[b]{\linewidth}\raggedright
\textbf{Umbral / Métrica}
\end{minipage} & \begin{minipage}[b]{\linewidth}\raggedright
\textbf{Método de verificación}
\end{minipage} \\
\midrule\noalign{}
\endhead
\bottomrule\noalign{}
\endlastfoot
\textbf{RNF-50} & Toda comunicación relativa al crédito debe cumplir criterios verificables de lenguaje claro, validados con personas usuarias reales. & 100\% de plantillas validadas con pruebas de comprensión lectora con usuarios reales. & Prueba de comprensión lectora sobre muestra de plantillas. \\
\end{longtable}

\subsubsection{A.2.16. Listado de Requerimientos No Funcionales: Operabilidad}\label{a-2-16-listado-de-requerimientos-no-funcionales-operabilidad}

A continuación, se detalla el listado de requerimientos no funcionales (RNF) de la categoría ``Operabilidad''.

\renewcommand{\thetable}{A.2.16}
\begin{longtable}[]{@{}
  >{\raggedright\arraybackslash}p{(\columnwidth - 8\tabcolsep) * \real{0.1027}}
  >{\raggedright\arraybackslash}p{(\columnwidth - 8\tabcolsep) * \real{0.4302}}
  >{\raggedright\arraybackslash}p{(\columnwidth - 8\tabcolsep) * \real{0.2087}}
  >{\raggedright\arraybackslash}p{(\columnwidth - 8\tabcolsep) * \real{0.2584}}@{}}
\caption{Requerimientos No Funcionales - Operabilidad.}\label{tab:a-2-16}\\
\toprule\noalign{}
\begin{minipage}[b]{\linewidth}\raggedright
\textbf{ID}
\end{minipage} & \begin{minipage}[b]{\linewidth}\raggedright
\textbf{Descripción del Requerimiento}
\end{minipage} & \begin{minipage}[b]{\linewidth}\raggedright
\textbf{Umbral / Métrica}
\end{minipage} & \begin{minipage}[b]{\linewidth}\raggedright
\textbf{Método de verificación}
\end{minipage} \\
\midrule\noalign{}
\endfirsthead
\toprule\noalign{}
\begin{minipage}[b]{\linewidth}\raggedright
\textbf{ID}
\end{minipage} & \begin{minipage}[b]{\linewidth}\raggedright
\textbf{Descripción del Requerimiento}
\end{minipage} & \begin{minipage}[b]{\linewidth}\raggedright
\textbf{Umbral / Métrica}
\end{minipage} & \begin{minipage}[b]{\linewidth}\raggedright
\textbf{Método de verificación}
\end{minipage} \\
\midrule\noalign{}
\endhead
\bottomrule\noalign{}
\endlastfoot
\textbf{RNF-51} & Las notificaciones definidas funcionalmente deben emitirse y registrarse en el 100\% de los eventos que las gatillan. & 100\% de eventos definidos con notificación generada y registrada. & Prueba de disparo de eventos y verificación del registro. \\
\addlinespace[2pt]
\textbf{RNF-64} & El tiempo entre la detección de un quiebre y el aviso al cliente no debe superar el objetivo declarado. & \textless{}= 30 minutos (supuesto fundamentado). & Medición de la latencia entre el evento de quiebre y el envío de la notificación. \\
\addlinespace[2pt]
\textbf{RNF-65} & No debe mantenerse un cobro capturado sobre un pedido declarado no cumplible. & 0 casos de cobro sostenido sobre pedido no cumplible. & Auditoría de conciliación diaria de preautorizaciones y capturas (RF-050). \\
\addlinespace[2pt]
\textbf{RNF-66} & La latencia entre la recepción física de una devolución en tienda y la notificación al vendedor de marketplace no debe superar el objetivo declarado. & \textless{}= 15 minutos (supuesto fundamentado). & Medición de la latencia entre el registro de recepción y el envío de la notificación. \\
\addlinespace[2pt]
\textbf{RNF-67} & No debe permanecer mercadería de marketplace en bodega sin aviso a su propietario. & 0 unidades en bodega sin notificación asociada al cierre de cada día. & Conciliación diaria entre recepciones registradas y notificaciones emitidas. \\
\end{longtable}

\subsubsection{A.2.17. Listado de Requerimientos No Funcionales: Interoperabilidad}\label{a-2-17-listado-de-requerimientos-no-funcionales-interoperabilidad}

A continuación, se detalla el listado de requerimientos no funcionales (RNF) de la categoría ``Interoperabilidad''.

\renewcommand{\thetable}{A.2.17}
\begin{longtable}[]{@{}
  >{\raggedright\arraybackslash}p{(\columnwidth - 8\tabcolsep) * \real{0.1027}}
  >{\raggedright\arraybackslash}p{(\columnwidth - 8\tabcolsep) * \real{0.4302}}
  >{\raggedright\arraybackslash}p{(\columnwidth - 8\tabcolsep) * \real{0.2087}}
  >{\raggedright\arraybackslash}p{(\columnwidth - 8\tabcolsep) * \real{0.2584}}@{}}
\caption{Requerimientos No Funcionales - Interoperabilidad.}\label{tab:a-2-17}\\
\toprule\noalign{}
\begin{minipage}[b]{\linewidth}\raggedright
\textbf{ID}
\end{minipage} & \begin{minipage}[b]{\linewidth}\raggedright
\textbf{Descripción del Requerimiento}
\end{minipage} & \begin{minipage}[b]{\linewidth}\raggedright
\textbf{Umbral / Métrica}
\end{minipage} & \begin{minipage}[b]{\linewidth}\raggedright
\textbf{Método de verificación}
\end{minipage} \\
\midrule\noalign{}
\endfirsthead
\toprule\noalign{}
\begin{minipage}[b]{\linewidth}\raggedright
\textbf{ID}
\end{minipage} & \begin{minipage}[b]{\linewidth}\raggedright
\textbf{Descripción del Requerimiento}
\end{minipage} & \begin{minipage}[b]{\linewidth}\raggedright
\textbf{Umbral / Métrica}
\end{minipage} & \begin{minipage}[b]{\linewidth}\raggedright
\textbf{Método de verificación}
\end{minipage} \\
\midrule\noalign{}
\endhead
\bottomrule\noalign{}
\endlastfoot
\textbf{RNF-52} & La solución debe adoptar un estándar identificado y justificado para el catálogo e identificación de producto. & Estándar identificado por su denominación y justificado en la Oferta Técnica. & Prueba de interoperabilidad con al menos una contraparte real del dominio. \\
\addlinespace[2pt]
\textbf{RNF-53} & La solución debe adoptar un estándar identificado y justificado para el intercambio de órdenes y avisos de despacho con los proveedores. & Estándar identificado y justificado; 100 \% de las interfaces con proveedores mapeadas a él. & Prueba de intercambio con un proveedor piloto. \\
\addlinespace[2pt]
\textbf{RNF-54} & La solución debe adoptar un estándar identificado y justificado para el intercambio con los transportistas de última milla. & Estándar identificado y justificado; trazabilidad del estado del pedido hasta la entrega. & Prueba de integración con un transportista piloto. \\
\addlinespace[2pt]
\textbf{RNF-55} & La solución debe adoptar un estándar identificado y justificado para la sincronización de catálogo, existencia y pedido con los vendedores de marketplace. & Estándar identificado y justificado; 100 \% de vendedores sincronizados por él. & Prueba de sincronización con un vendedor piloto. \\
\addlinespace[2pt]
\textbf{RNF-56} & La solución debe adoptar el formato normativo vigente de documento tributario electrónico. & 100 \% de los documentos emitidos conforme al formato vigente de la autoridad tributaria. & Validación de esquema sobre una muestra de documentos emitidos. \\
\end{longtable}

\subsubsection{A.2.18. Listado de Requerimientos No Funcionales: Auditabilidad / Desempeño}\label{a-2-18-listado-de-requerimientos-no-funcionales-auditabilidad-desempeno}

A continuación, se detalla el listado de requerimientos no funcionales (RNF) de la categoría ``Auditabilidad / Desempeño''.

\renewcommand{\thetable}{A.2.18}
\begin{longtable}[]{@{}
  >{\raggedright\arraybackslash}p{(\columnwidth - 8\tabcolsep) * \real{0.1027}}
  >{\raggedright\arraybackslash}p{(\columnwidth - 8\tabcolsep) * \real{0.4302}}
  >{\raggedright\arraybackslash}p{(\columnwidth - 8\tabcolsep) * \real{0.2087}}
  >{\raggedright\arraybackslash}p{(\columnwidth - 8\tabcolsep) * \real{0.2584}}@{}}
\caption{Requerimientos No Funcionales - Auditabilidad / Desempeño.}\label{tab:a-2-18}\\
\toprule\noalign{}
\begin{minipage}[b]{\linewidth}\raggedright
\textbf{ID}
\end{minipage} & \begin{minipage}[b]{\linewidth}\raggedright
\textbf{Descripción del Requerimiento}
\end{minipage} & \begin{minipage}[b]{\linewidth}\raggedright
\textbf{Umbral / Métrica}
\end{minipage} & \begin{minipage}[b]{\linewidth}\raggedright
\textbf{Método de verificación}
\end{minipage} \\
\midrule\noalign{}
\endfirsthead
\toprule\noalign{}
\begin{minipage}[b]{\linewidth}\raggedright
\textbf{ID}
\end{minipage} & \begin{minipage}[b]{\linewidth}\raggedright
\textbf{Descripción del Requerimiento}
\end{minipage} & \begin{minipage}[b]{\linewidth}\raggedright
\textbf{Umbral / Métrica}
\end{minipage} & \begin{minipage}[b]{\linewidth}\raggedright
\textbf{Método de verificación}
\end{minipage} \\
\midrule\noalign{}
\endhead
\bottomrule\noalign{}
\endlastfoot
\textbf{RNF-57} & La evidencia completa de consentimiento de una operación arbitraria debe recuperarse dentro del plazo objetivo, a diez años. & Recuperación \textless{}= 5 minutos (supuesto fundamentado) sobre operaciones de hasta 10 años de antigüedad. & Simulacro de requerimiento de la autoridad fiscalizadora. \\
\end{longtable}

\subsubsection{A.2.19. Listado de Requerimientos No Funcionales: Usabilidad / Desempeño}\label{a-2-19-listado-de-requerimientos-no-funcionales-usabilidad-desempeno}

A continuación, se detalla el listado de requerimientos no funcionales (RNF) de la categoría ``Usabilidad / Desempeño''.

\renewcommand{\thetable}{A.2.19}
\begin{longtable}[]{@{}
  >{\raggedright\arraybackslash}p{(\columnwidth - 8\tabcolsep) * \real{0.1027}}
  >{\raggedright\arraybackslash}p{(\columnwidth - 8\tabcolsep) * \real{0.4302}}
  >{\raggedright\arraybackslash}p{(\columnwidth - 8\tabcolsep) * \real{0.2087}}
  >{\raggedright\arraybackslash}p{(\columnwidth - 8\tabcolsep) * \real{0.2584}}@{}}
\caption{Requerimientos No Funcionales - Usabilidad / Desempeño.}\label{tab:a-2-19}\\
\toprule\noalign{}
\begin{minipage}[b]{\linewidth}\raggedright
\textbf{ID}
\end{minipage} & \begin{minipage}[b]{\linewidth}\raggedright
\textbf{Descripción del Requerimiento}
\end{minipage} & \begin{minipage}[b]{\linewidth}\raggedright
\textbf{Umbral / Métrica}
\end{minipage} & \begin{minipage}[b]{\linewidth}\raggedright
\textbf{Método de verificación}
\end{minipage} \\
\midrule\noalign{}
\endfirsthead
\toprule\noalign{}
\begin{minipage}[b]{\linewidth}\raggedright
\textbf{ID}
\end{minipage} & \begin{minipage}[b]{\linewidth}\raggedright
\textbf{Descripción del Requerimiento}
\end{minipage} & \begin{minipage}[b]{\linewidth}\raggedright
\textbf{Umbral / Métrica}
\end{minipage} & \begin{minipage}[b]{\linewidth}\raggedright
\textbf{Método de verificación}
\end{minipage} \\
\midrule\noalign{}
\endhead
\bottomrule\noalign{}
\endlastfoot
\textbf{RNF-60} & El traspaso de sesión nominativa en una terminal compartida debe completarse dentro del plazo objetivo. & \textless{}= 5 segundos (supuesto fundamentado). & Medición de tiempo de traspaso en terminal de sala. \\
\end{longtable}

\subsubsection{A.2.20. Listado de Requerimientos No Funcionales: Efectividad de negocio}\label{a-2-20-listado-de-requerimientos-no-funcionales-efectividad-de-negocio}

A continuación, se detalla el listado de requerimientos no funcionales (RNF) de la categoría ``Efectividad de negocio''.

\renewcommand{\thetable}{A.2.20}
\begin{longtable}[]{@{}
  >{\raggedright\arraybackslash}p{(\columnwidth - 8\tabcolsep) * \real{0.1027}}
  >{\raggedright\arraybackslash}p{(\columnwidth - 8\tabcolsep) * \real{0.4302}}
  >{\raggedright\arraybackslash}p{(\columnwidth - 8\tabcolsep) * \real{0.2087}}
  >{\raggedright\arraybackslash}p{(\columnwidth - 8\tabcolsep) * \real{0.2584}}@{}}
\caption{Requerimientos No Funcionales - Efectividad de negocio.}\label{tab:a-2-20}\\
\toprule\noalign{}
\begin{minipage}[b]{\linewidth}\raggedright
\textbf{ID}
\end{minipage} & \begin{minipage}[b]{\linewidth}\raggedright
\textbf{Descripción del Requerimiento}
\end{minipage} & \begin{minipage}[b]{\linewidth}\raggedright
\textbf{Umbral / Métrica}
\end{minipage} & \begin{minipage}[b]{\linewidth}\raggedright
\textbf{Método de verificación}
\end{minipage} \\
\midrule\noalign{}
\endfirsthead
\toprule\noalign{}
\begin{minipage}[b]{\linewidth}\raggedright
\textbf{ID}
\end{minipage} & \begin{minipage}[b]{\linewidth}\raggedright
\textbf{Descripción del Requerimiento}
\end{minipage} & \begin{minipage}[b]{\linewidth}\raggedright
\textbf{Umbral / Métrica}
\end{minipage} & \begin{minipage}[b]{\linewidth}\raggedright
\textbf{Método de verificación}
\end{minipage} \\
\midrule\noalign{}
\endhead
\bottomrule\noalign{}
\endlastfoot
\textbf{RNF-61} & La tasa de cancelación de pedidos por falta de existencia no debe superar el umbral comprometido. & \textless{} 0,3\% anual. & Medición anual sobre el total de pedidos aceptados. \\
\addlinespace[2pt]
\textbf{RNF-62} & El cumplimiento de la promesa de entrega publicada debe alcanzar el umbral comprometido, medido por pedido. & \textgreater{}= 97\% (línea base 81\%), medido por pedido y no por promedio de canal. & Medición por pedido sobre el período. \\
\addlinespace[2pt]
\textbf{RNF-63} & La diferencia de inventario por categoría debe reducirse al umbral comprometido. & \textless{} 2\% (linea base 12,4\%). & Medición sobre resultados de conteo cíclico por categoría. \\
\end{longtable}

\subsubsection{A.2.21. Listado de Requerimientos No Funcionales: Seguridad / DevSecOps}\label{a-2-21-listado-de-requerimientos-no-funcionales-seguridad-devsecops}

A continuación, se detalla el listado de requerimientos no funcionales (RNF) de la categoría ``Seguridad / DevSecOps''.

\renewcommand{\thetable}{A.2.21}
\begin{longtable}[]{@{}
  >{\raggedright\arraybackslash}p{(\columnwidth - 8\tabcolsep) * \real{0.1027}}
  >{\raggedright\arraybackslash}p{(\columnwidth - 8\tabcolsep) * \real{0.4302}}
  >{\raggedright\arraybackslash}p{(\columnwidth - 8\tabcolsep) * \real{0.2087}}
  >{\raggedright\arraybackslash}p{(\columnwidth - 8\tabcolsep) * \real{0.2584}}@{}}
\caption{Requerimientos No Funcionales - Seguridad / DevSecOps.}\label{tab:a-2-21}\\
\toprule\noalign{}
\begin{minipage}[b]{\linewidth}\raggedright
\textbf{ID}
\end{minipage} & \begin{minipage}[b]{\linewidth}\raggedright
\textbf{Descripción del Requerimiento}
\end{minipage} & \begin{minipage}[b]{\linewidth}\raggedright
\textbf{Umbral / Métrica}
\end{minipage} & \begin{minipage}[b]{\linewidth}\raggedright
\textbf{Método de verificación}
\end{minipage} \\
\midrule\noalign{}
\endfirsthead
\toprule\noalign{}
\begin{minipage}[b]{\linewidth}\raggedright
\textbf{ID}
\end{minipage} & \begin{minipage}[b]{\linewidth}\raggedright
\textbf{Descripción del Requerimiento}
\end{minipage} & \begin{minipage}[b]{\linewidth}\raggedright
\textbf{Umbral / Métrica}
\end{minipage} & \begin{minipage}[b]{\linewidth}\raggedright
\textbf{Método de verificación}
\end{minipage} \\
\midrule\noalign{}
\endhead
\bottomrule\noalign{}
\endlastfoot
\textbf{RNF-70} & El ciclo de desarrollo debe producir un inventario de componentes (SBOM) por artefacto desplegado. & SBOM del 100 \% de los artefactos desplegados. & Verificación del SBOM en cada liberación. \\
\addlinespace[2pt]
\textbf{RNF-71} & La cadena de suministro de software debe alcanzar el nivel SLSA 3 o superior. & Nivel SLSA 3 o superior acreditado. & Atestación de la cadena de construcción por artefacto. \\
\addlinespace[2pt]
\textbf{RNF-72} & La capa expuesta debe cumplir el estándar OWASP ASVS en el nivel comprometido. & 0 hallazgos críticos abiertos en el nivel ASVS comprometido. & Análisis estático y dinámico por ciclo y prueba de seguridad ofensiva. \\
\addlinespace[2pt]
\textbf{RNF-73} & La arquitectura debe implementar un modelo Zero Trust conforme a NIST SP 800-207. & Verificación explícita en cada solicitud y denegación por omisión en el 100 \% de los puntos de control. & Revisión de arquitectura contra los principios de NIST SP 800-207. \\
\end{longtable}

\subsection{A.3. Listado de Obligaciones del Proponente (OP)}\label{a-3-listado-de-obligaciones-del-proponente-op}

A continuación, se detalla el listado de obligaciones del proponente (OP) extraídas del Catálogo Atómico: compromisos de la oferta que no constituyen funciones del sistema y por lo tanto no admiten caso de prueba de software. Total: 9 obligaciones.

\renewcommand{\thetable}{A.3.1}
\begin{longtable}[]{@{}
  >{\raggedright\arraybackslash}p{(\columnwidth - 4\tabcolsep) * \real{0.0865}}
  >{\raggedright\arraybackslash}p{(\columnwidth - 4\tabcolsep) * \real{0.9135}}@{}}
\caption{Obligación del PROPONENTE.}\label{tab:a-3-1}\\
\toprule\noalign{}
\begin{minipage}[b]{\linewidth}\raggedright
\textbf{ID}
\end{minipage} & \begin{minipage}[b]{\linewidth}\raggedright
\textbf{Descripción del Requerimiento}
\end{minipage} \\
\midrule\noalign{}
\endfirsthead
\toprule\noalign{}
\begin{minipage}[b]{\linewidth}\raggedright
\textbf{ID}
\end{minipage} & \begin{minipage}[b]{\linewidth}\raggedright
\textbf{Descripción del Requerimiento}
\end{minipage} \\
\midrule\noalign{}
\endhead
\bottomrule\noalign{}
\endlastfoot
\textbf{OP-01} & El PROPONENTE debe especificar la alternativa de etiquetas electrónicas frente a las demás opciones de resolución de la discrepancia de precio. \\
\addlinespace[2pt]
\textbf{OP-02} & El PROPONENTE debe costear la alternativa de etiquetas electrónicas sin asumir su adquisición dentro del alcance. \\
\addlinespace[2pt]
\textbf{OP-03} & El PROPONENTE debe especificar la cobertura de despliegue de las etiquetas electrónicas. \\
\addlinespace[2pt]
\textbf{OP-04} & El PROPONENTE debe especificar el hardware requerido por la alternativa de etiquetas electrónicas. \\
\addlinespace[2pt]
\textbf{OP-05} & El PROPONENTE debe presentar el costo de la alternativa separadamente del resto de la solución. \\
\addlinespace[2pt]
\textbf{OP-06} & El PROPONENTE debe especificar la cantidad de dispositivos móviles requeridos por tienda y por departamento. \\
\addlinespace[2pt]
\textbf{OP-07} & El PROPONENTE debe especificar las características técnicas de los dispositivos móviles requeridos. \\
\addlinespace[2pt]
\textbf{OP-08} & El PROPONENTE debe presentar el análisis de hacer o comprar del centro de distribución de Concepción con sus tres alternativas. \\
\addlinespace[2pt]
\textbf{OP-09} & El PROPONENTE debe declarar el impacto de la alternativa seleccionada sobre el cumplimiento de RN-15. \\
\end{longtable}

\subsection{A.4. Matriz de cumplimiento y trazabilidad}\label{a-4-matriz-de-cumplimiento-y-trazabilidad}

\renewcommand{\thetable}{A.4.1}
\begin{longtable}[]{@{}
  >{\raggedright\arraybackslash\footnotesize}p{(\columnwidth - 10\tabcolsep) * \real{0.1300}}
  >{\raggedright\arraybackslash\footnotesize}p{(\columnwidth - 10\tabcolsep) * \real{0.3100}}
  >{\raggedright\arraybackslash\footnotesize}p{(\columnwidth - 10\tabcolsep) * \real{0.0800}}
  >{\raggedright\arraybackslash\footnotesize}p{(\columnwidth - 10\tabcolsep) * \real{0.2600}}
  >{\raggedright\arraybackslash\footnotesize}p{(\columnwidth - 10\tabcolsep) * \real{0.2200}}@{}}
\caption{Matriz de cumplimiento y trazabilidad.}\label{tab:a-4-1}\\
\toprule\noalign{}
\begin{minipage}[b]{\linewidth}\raggedright
\textbf{ID Requerimiento}
\end{minipage} & \begin{minipage}[b]{\linewidth}\raggedright
\textbf{Descripción}
\end{minipage} & \begin{minipage}[b]{\linewidth}\raggedright
\textbf{Cumple}
\end{minipage} & \begin{minipage}[b]{\linewidth}\raggedright
\textbf{Componente que lo satisface}
\end{minipage} & \begin{minipage}[b]{\linewidth}\raggedright
\textbf{Sección de la propuesta}
\end{minipage} \\
\midrule\noalign{}
\endfirsthead
\toprule\noalign{}
\begin{minipage}[b]{\linewidth}\raggedright
\textbf{ID Requerimiento}
\end{minipage} & \begin{minipage}[b]{\linewidth}\raggedright
\textbf{Descripción}
\end{minipage} & \begin{minipage}[b]{\linewidth}\raggedright
\textbf{Cumple}
\end{minipage} & \begin{minipage}[b]{\linewidth}\raggedright
\textbf{Componente que lo satisface}
\end{minipage} & \begin{minipage}[b]{\linewidth}\raggedright
\textbf{Sección de la propuesta}
\end{minipage} \\
\midrule\noalign{}
\endhead
\bottomrule\noalign{}
\endlastfoot
RF-001 & El sistema deberá permitir la venta física de la unidad aunque exista una reserva temporal activa en el canal digital. & Sí & M-09 Orquestador de Pedidos (reserva temporal / prelación de canal presencial) & Subdoc.3 §3.2.6.1 (SUP-12, Tabla 3.1); §3.3.3 Tabla 3.9 (M-09) \\
\addlinespace[2pt]
RF-002 & El sistema deberá exigir credenciales individuales al usuario antes de habilitar cualquier acción. & Sí & M-23 Identidad y Accesos de Personas (IAM) — traspaso de sesión, revocación por finiquito, aprovisionamiento B2B & Subdoc.3 §3.2.6.6 Tabla 3.6 (SUP-19, SUP-24); §3.3.6 Tabla 3.12 (M-23) \textbar{} Subdoc.4 §4.1.7 (frente 2, IAM) \\
\addlinespace[2pt]
RF-003 & El sistema debe impedir el acceso anónimo a una terminal. & Sí & M-23 Identidad y Accesos de Personas (IAM) — traspaso de sesión, revocación por finiquito, aprovisionamiento B2B & Subdoc.3 §3.2.6.6 Tabla 3.6 (SUP-19, SUP-24); §3.3.6 Tabla 3.12 (M-23) \textbar{} Subdoc.4 §4.1.7 (frente 2, IAM) \\
\addlinespace[2pt]
RF-004 & El sistema debe impedir el acceso mediante credencial compartida. & Sí & M-23 Identidad y Accesos de Personas (IAM) — traspaso de sesión, revocación por finiquito, aprovisionamiento B2B & Subdoc.3 §3.2.6.6 Tabla 3.6 (SUP-19, SUP-24); §3.3.6 Tabla 3.12 (M-23) \textbar{} Subdoc.4 §4.1.7 (frente 2, IAM) \\
\addlinespace[2pt]
RF-005 & El sistema deberá permitir el traspaso de sesión nominativa entre usuarios en una terminal compartida. & Sí & M-23 Identidad y Accesos de Personas (IAM) — traspaso de sesión, revocación por finiquito, aprovisionamiento B2B & Subdoc.3 §3.2.6.6 Tabla 3.6 (SUP-19, SUP-24); §3.3.6 Tabla 3.12 (M-23) \textbar{} Subdoc.4 §4.1.7 (frente 2, IAM) \\
\addlinespace[2pt]
RF-006 & El sistema deberá cerrar automáticamente la sesión por inactividad en la terminal compartida. & Sí & M-23 Identidad y Accesos de Personas (IAM) — traspaso de sesión, revocación por finiquito, aprovisionamiento B2B & Subdoc.3 §3.2.6.6 Tabla 3.6 (SUP-19, SUP-24); §3.3.6 Tabla 3.12 (M-23) \textbar{} Subdoc.4 §4.1.7 (frente 2, IAM) \\
\addlinespace[2pt]
RF-007 & El sistema deberá impedir la ejecución de la función de originación a un usuario sin capacitación normativa acreditada y vigente. & Sí & M-23 Identidad y Accesos de Personas (IAM) — traspaso de sesión, revocación por finiquito, aprovisionamiento B2B & Subdoc.3 §3.2.6.6 Tabla 3.6 (SUP-19, SUP-24); §3.3.6 Tabla 3.12 (M-23) \textbar{} Subdoc.4 §4.1.7 (frente 2, IAM) \\
\addlinespace[2pt]
RF-008 & El sistema deberá impedir la ejecución de la función de repactación a un usuario sin capacitación normativa acreditada y vigente. & Sí & M-23 Identidad y Accesos de Personas (IAM) — traspaso de sesión, revocación por finiquito, aprovisionamiento B2B & Subdoc.3 §3.2.6.6 Tabla 3.6 (SUP-19, SUP-24); §3.3.6 Tabla 3.12 (M-23) \textbar{} Subdoc.4 §4.1.7 (frente 2, IAM) \\
\addlinespace[2pt]
RF-009 & El sistema deberá revocar automáticamente la habilitación funcional del usuario al vencer su acreditación de capacitación. & Sí & M-23 Identidad y Accesos de Personas (IAM) — traspaso de sesión, revocación por finiquito, aprovisionamiento B2B & Subdoc.3 §3.2.6.6 Tabla 3.6 (SUP-19, SUP-24); §3.3.6 Tabla 3.12 (M-23) \textbar{} Subdoc.4 §4.1.7 (frente 2, IAM) \\
\addlinespace[2pt]
RF-010 & El sistema deberá permitir al proveedor patrocinar la habilitación temporal de acceso de su repositor externo, con fecha de término declarada. & Sí & M-23 Identidad y Accesos de Personas (IAM) — traspaso de sesión, revocación por finiquito, aprovisionamiento B2B & Subdoc.3 §3.2.6.6 Tabla 3.6 (SUP-19, SUP-24); §3.3.6 Tabla 3.12 (M-23) \textbar{} Subdoc.4 §4.1.7 (frente 2, IAM) \\
\addlinespace[2pt]
RF-011 & El sistema deberá caducar automáticamente la habilitación de acceso del personal externo al cumplirse su vigencia. & Sí & M-23 Identidad y Accesos de Personas (IAM) — traspaso de sesión, revocación por finiquito, aprovisionamiento B2B & Subdoc.3 §3.2.6.6 Tabla 3.6 (SUP-19, SUP-24); §3.3.6 Tabla 3.12 (M-23) \textbar{} Subdoc.4 §4.1.7 (frente 2, IAM) \\
\addlinespace[2pt]
RF-012 & El sistema deberá registrar de forma individualizada el acceso y la actividad del personal externo. & Sí & M-23 Identidad y Accesos de Personas (IAM) — traspaso de sesión, revocación por finiquito, aprovisionamiento B2B & Subdoc.3 §3.2.6.6 Tabla 3.6 (SUP-19, SUP-24); §3.3.6 Tabla 3.12 (M-23) \textbar{} Subdoc.4 §4.1.7 (frente 2, IAM) \\
\addlinespace[2pt]
RF-013 & El sistema deberá revocar la totalidad de los accesos y credenciales del trabajador a partir del término efectivo de su vínculo. & Sí & M-23 Identidad y Accesos de Personas (IAM) — traspaso de sesión, revocación por finiquito, aprovisionamiento B2B & Subdoc.3 §3.2.6.6 Tabla 3.6 (SUP-19, SUP-24); §3.3.6 Tabla 3.12 (M-23) \textbar{} Subdoc.4 §4.1.7 (frente 2, IAM) \\
\addlinespace[2pt]
RF-014 & El sistema debe conciliar los accesos vigentes contra la nómina activa. & Sí & M-23 Identidad y Accesos de Personas (IAM) — traspaso de sesión, revocación por finiquito, aprovisionamiento B2B & Subdoc.3 §3.2.6.6 Tabla 3.6 (SUP-19, SUP-24); §3.3.6 Tabla 3.12 (M-23) \textbar{} Subdoc.4 §4.1.7 (frente 2, IAM) \\
\addlinespace[2pt]
RF-015 & El sistema debe informar al oficial de seguridad los accesos huérfanos detectados en la conciliación. & Sí & M-23 Identidad y Accesos de Personas (IAM) — traspaso de sesión, revocación por finiquito, aprovisionamiento B2B & Subdoc.3 §3.2.6.6 Tabla 3.6 (SUP-19, SUP-24); §3.3.6 Tabla 3.12 (M-23) \textbar{} Subdoc.4 §4.1.7 (frente 2, IAM) \\
\addlinespace[2pt]
RF-016 & El sistema deberá propagar el cambio de precio a las 380 líneas de caja. & Sí & M-05 Motor de Precios y Promociones / M-06 Gobierno de Etiquetado Físico / M-07 Repositorio de Trazabilidad de Precio & Subdoc.3 §3.2.6.2 Tablas 3.2 (SUP-08, SUP-09, SUP-10); §3.3.2 Tabla 3.8 \textbar{} Subdoc.4 §4.1.5.4 (evento PrecioPublicado) \\
\addlinespace[2pt]
RF-017 & El sistema deberá propagar el cambio de precio al canal digital. & Sí & M-05 Motor de Precios y Promociones / M-06 Gobierno de Etiquetado Físico / M-07 Repositorio de Trazabilidad de Precio & Subdoc.3 §3.2.6.2 Tablas 3.2 (SUP-08, SUP-09, SUP-10); §3.3.2 Tabla 3.8 \textbar{} Subdoc.4 §4.1.5.4 (evento PrecioPublicado) \\
\addlinespace[2pt]
RF-018 & El sistema deberá propagar el cambio de precio a los puntos de exhibición física de cada tienda afectada. & Sí & M-05 Motor de Precios y Promociones / M-06 Gobierno de Etiquetado Físico / M-07 Repositorio de Trazabilidad de Precio & Subdoc.3 §3.2.6.2 Tablas 3.2 (SUP-08, SUP-09, SUP-10); §3.3.2 Tabla 3.8 \textbar{} Subdoc.4 §4.1.5.4 (evento PrecioPublicado) \\
\addlinespace[2pt]
RF-019 & El sistema deberá permitir al reponedor registrar la referencia y la tienda de la etiqueta cambiada, con el instante del cambio. & Sí & M-05 Motor de Precios y Promociones / M-06 Gobierno de Etiquetado Físico / M-07 Repositorio de Trazabilidad de Precio & Subdoc.3 §3.2.6.2 Tablas 3.2 (SUP-08, SUP-09, SUP-10); §3.3.2 Tabla 3.8 \textbar{} Subdoc.4 §4.1.5.4 (evento PrecioPublicado) \\
\addlinespace[2pt]
RF-020 & El sistema deberá registrar la identidad individual del ejecutor del cambio de etiqueta. & Sí & M-05 Motor de Precios y Promociones / M-06 Gobierno de Etiquetado Físico / M-07 Repositorio de Trazabilidad de Precio & Subdoc.3 §3.2.6.2 Tablas 3.2 (SUP-08, SUP-09, SUP-10); §3.3.2 Tabla 3.8 \textbar{} Subdoc.4 §4.1.5.4 (evento PrecioPublicado) \\
\addlinespace[2pt]
RF-021 & El sistema deberá permitir al ejecutivo de cumplimiento recuperar el precio publicado de una referencia para una fecha, hora y canal determinados. & Sí & M-05 Motor de Precios y Promociones / M-06 Gobierno de Etiquetado Físico / M-07 Repositorio de Trazabilidad de Precio & Subdoc.3 §3.2.6.2 Tablas 3.2 (SUP-08, SUP-09, SUP-10); §3.3.2 Tabla 3.8 \textbar{} Subdoc.4 §4.1.5.4 (evento PrecioPublicado) \\
\addlinespace[2pt]
RF-022 & El sistema deberá obtener el precio vigente de la referencia según el último precio propagado a la etiqueta de esa tienda en la fecha de la venta. & Sí & M-05 Motor de Precios y Promociones / M-06 Gobierno de Etiquetado Físico / M-07 Repositorio de Trazabilidad de Precio & Subdoc.3 §3.2.6.2 Tablas 3.2 (SUP-08, SUP-09, SUP-10); §3.3.2 Tabla 3.8 \textbar{} Subdoc.4 §4.1.5.4 (evento PrecioPublicado) \\
\addlinespace[2pt]
RF-023 & El sistema deberá cobrar el menor de ambos precios para el consumidor. & Sí & M-05 Motor de Precios y Promociones / M-06 Gobierno de Etiquetado Físico / M-07 Repositorio de Trazabilidad de Precio & Subdoc.3 §3.2.6.2 Tablas 3.2 (SUP-08, SUP-09, SUP-10); §3.3.2 Tabla 3.8 \textbar{} Subdoc.4 §4.1.5.4 (evento PrecioPublicado) \\
\addlinespace[2pt]
RF-024 & El sistema deberá generar un registro de incidente de discrepancia de precio consultable. & Sí & M-05 Motor de Precios y Promociones / M-06 Gobierno de Etiquetado Físico / M-07 Repositorio de Trazabilidad de Precio & Subdoc.3 §3.2.6.2 Tablas 3.2 (SUP-08, SUP-09, SUP-10); §3.3.2 Tabla 3.8 \textbar{} Subdoc.4 §4.1.5.4 (evento PrecioPublicado) \\
\addlinespace[2pt]
RF-028 & El sistema deberá impedir la venta de la referencia al precio nuevo mientras su punto de exhibición no confirme la actualización, salvo que el artículo se encuentre bloqueado para venta. & Sí & M-05 Motor de Precios y Promociones / M-06 Gobierno de Etiquetado Físico / M-07 Repositorio de Trazabilidad de Precio & Subdoc.3 §3.2.6.2 Tablas 3.2 (SUP-08, SUP-09, SUP-10); §3.3.2 Tabla 3.8 \textbar{} Subdoc.4 §4.1.5.4 (evento PrecioPublicado) \\
\addlinespace[2pt]
RF-029 & El sistema deberá generar un indicador diario de puntos de exhibición desactualizados con identificación individual de cada punto. & Sí & M-05 Motor de Precios y Promociones / M-06 Gobierno de Etiquetado Físico / M-07 Repositorio de Trazabilidad de Precio & Subdoc.3 §3.2.6.2 Tablas 3.2 (SUP-08, SUP-09, SUP-10); §3.3.2 Tabla 3.8 \textbar{} Subdoc.4 §4.1.5.4 (evento PrecioPublicado) \\
\addlinespace[2pt]
RF-030 & El sistema deberá agrupar los cambios de precio de sala en ventanas parametrizadas fuera del horario de atención. & Sí & M-05 Motor de Precios y Promociones / M-06 Gobierno de Etiquetado Físico / M-07 Repositorio de Trazabilidad de Precio & Subdoc.3 §3.2.6.2 Tablas 3.2 (SUP-08, SUP-09, SUP-10); §3.3.2 Tabla 3.8 \textbar{} Subdoc.4 §4.1.5.4 (evento PrecioPublicado) \\
\addlinespace[2pt]
RF-032 & El sistema deberá evaluar la vigencia de cada promoción al instante de emisión del documento. & Sí & M-05 Motor de Precios y Promociones / M-06 Gobierno de Etiquetado Físico / M-07 Repositorio de Trazabilidad de Precio & Subdoc.3 §3.2.6.2 Tablas 3.2 (SUP-08, SUP-09, SUP-10); §3.3.2 Tabla 3.8 \textbar{} Subdoc.4 §4.1.5.4 (evento PrecioPublicado) \\
\addlinespace[2pt]
RF-033 & El sistema deberá evaluar la aplicabilidad de la promoción según el canal de la venta. & Sí & M-05 Motor de Precios y Promociones / M-06 Gobierno de Etiquetado Físico / M-07 Repositorio de Trazabilidad de Precio & Subdoc.3 §3.2.6.2 Tablas 3.2 (SUP-08, SUP-09, SUP-10); §3.3.2 Tabla 3.8 \textbar{} Subdoc.4 §4.1.5.4 (evento PrecioPublicado) \\
\addlinespace[2pt]
RF-034 & El sistema deberá evaluar la aplicabilidad de la promoción según la tienda en que se realiza la venta. & Sí & M-05 Motor de Precios y Promociones / M-06 Gobierno de Etiquetado Físico / M-07 Repositorio de Trazabilidad de Precio & Subdoc.3 §3.2.6.2 Tablas 3.2 (SUP-08, SUP-09, SUP-10); §3.3.2 Tabla 3.8 \textbar{} Subdoc.4 §4.1.5.4 (evento PrecioPublicado) \\
\addlinespace[2pt]
RF-035 & El sistema deberá crear una reserva temporal para la unidad agregada al carro. & Sí & M-09 Orquestador de Pedidos (OMS) / M-10 Atribución de Venta y Base de Comisión / M-11 Portal del Cliente & Subdoc.3 §3.2.6.3 Tabla 3.3 (SUP-11, SUP-12, SUP-13, SUP-14); §3.3.3 Tabla 3.9 \textbar{} Subdoc.4 §4.1.8.1 Flujo 2 \\
\addlinespace[2pt]
RF-036 & El sistema deberá permitir parametrizar el período de vigencia de la reserva temporal. & Sí & M-09 Orquestador de Pedidos (OMS) / M-10 Atribución de Venta y Base de Comisión / M-11 Portal del Cliente & Subdoc.3 §3.2.6.3 Tabla 3.3 (SUP-11, SUP-12, SUP-13, SUP-14); §3.3.3 Tabla 3.9 \textbar{} Subdoc.4 §4.1.8.1 Flujo 2 \\
\addlinespace[2pt]
RF-037 & El sistema deberá mantener activa la reserva mientras el cliente mantenga actividad, conforme al período parametrizado. & Sí & M-09 Orquestador de Pedidos (OMS) / M-10 Atribución de Venta y Base de Comisión / M-11 Portal del Cliente & Subdoc.3 §3.2.6.3 Tabla 3.3 (SUP-11, SUP-12, SUP-13, SUP-14); §3.3.3 Tabla 3.9 \textbar{} Subdoc.4 §4.1.8.1 Flujo 2 \\
\addlinespace[2pt]
RF-038 & El sistema deberá impedir que una misma unidad sea comprometida simultáneamente en más de una operación de venta digital. & Sí & M-09 Orquestador de Pedidos (OMS) / M-10 Atribución de Venta y Base de Comisión / M-11 Portal del Cliente & Subdoc.3 §3.2.6.3 Tabla 3.3 (SUP-11, SUP-12, SUP-13, SUP-14); §3.3.3 Tabla 3.9 \textbar{} Subdoc.4 §4.1.8.1 Flujo 2 \\
\addlinespace[2pt]
RF-039 & El sistema debe liberar la reserva temporal al expirar su período de vigencia. & Sí & M-09 Orquestador de Pedidos (OMS) / M-10 Atribución de Venta y Base de Comisión / M-11 Portal del Cliente & Subdoc.3 §3.2.6.3 Tabla 3.3 (SUP-11, SUP-12, SUP-13, SUP-14); §3.3.3 Tabla 3.9 \textbar{} Subdoc.4 §4.1.8.1 Flujo 2 \\
\addlinespace[2pt]
RF-040 & El sistema debe devolver la unidad liberada al disponible para vender. & Sí & M-09 Orquestador de Pedidos (OMS) / M-10 Atribución de Venta y Base de Comisión / M-11 Portal del Cliente & Subdoc.3 §3.2.6.3 Tabla 3.3 (SUP-11, SUP-12, SUP-13, SUP-14); §3.3.3 Tabla 3.9 \textbar{} Subdoc.4 §4.1.8.1 Flujo 2 \\
\addlinespace[2pt]
RF-041 & El sistema deberá liberar de inmediato la reserva digital de esa unidad, sin esperar la expiración del período. & Sí & M-09 Orquestador de Pedidos (OMS) / M-10 Atribución de Venta y Base de Comisión / M-11 Portal del Cliente & Subdoc.3 §3.2.6.3 Tabla 3.3 (SUP-11, SUP-12, SUP-13, SUP-14); §3.3.3 Tabla 3.9 \textbar{} Subdoc.4 §4.1.8.1 Flujo 2 \\
\addlinespace[2pt]
RF-042 & El sistema deberá aplicar el período de vigencia de reserva reducido declarado como parámetro para el evento anual. & Sí & M-09 Orquestador de Pedidos (OMS) / M-10 Atribución de Venta y Base de Comisión / M-11 Portal del Cliente & Subdoc.3 §3.2.6.3 Tabla 3.3 (SUP-11, SUP-12, SUP-13, SUP-14); §3.3.3 Tabla 3.9 \textbar{} Subdoc.4 §4.1.8.1 Flujo 2 \\
\addlinespace[2pt]
RF-043 & El sistema deberá preautorizar el medio de pago al aceptar el pedido, sin capturar el cobro. & Sí & M-09 Orquestador de Pedidos (OMS) / M-10 Atribución de Venta y Base de Comisión / M-11 Portal del Cliente & Subdoc.3 §3.2.6.3 Tabla 3.3 (SUP-11, SUP-12, SUP-13, SUP-14); §3.3.3 Tabla 3.9 \textbar{} Subdoc.4 §4.1.8.1 Flujo 2 \\
\addlinespace[2pt]
RF-044 & El sistema deberá verificar la existencia física real de la unidad en el punto asignado antes de habilitar la captura del cobro. & Sí & M-09 Orquestador de Pedidos (OMS) / M-10 Atribución de Venta y Base de Comisión / M-11 Portal del Cliente & Subdoc.3 §3.2.6.3 Tabla 3.3 (SUP-11, SUP-12, SUP-13, SUP-14); §3.3.3 Tabla 3.9 \textbar{} Subdoc.4 §4.1.8.1 Flujo 2 \\
\addlinespace[2pt]
RF-045 & El sistema deberá capturar el cobro únicamente al registrarse el evento de confirmación de la preparacion física. & Sí & M-09 Orquestador de Pedidos (OMS) / M-10 Atribución de Venta y Base de Comisión / M-11 Portal del Cliente & Subdoc.3 §3.2.6.3 Tabla 3.3 (SUP-11, SUP-12, SUP-13, SUP-14); §3.3.3 Tabla 3.9 \textbar{} Subdoc.4 §4.1.8.1 Flujo 2 \\
\addlinespace[2pt]
RF-046 & El sistema deberá determinar automáticamente la alternativa de resolución aplicable según el motor de reglas (cancelar, sustituir, derivar a tercero o entrega diferida). & Sí & M-09 Orquestador de Pedidos (OMS) / M-10 Atribución de Venta y Base de Comisión / M-11 Portal del Cliente & Subdoc.3 §3.2.6.3 Tabla 3.3 (SUP-11, SUP-12, SUP-13, SUP-14); §3.3.3 Tabla 3.9 \textbar{} Subdoc.4 §4.1.8.1 Flujo 2 \\
\addlinespace[2pt]
RF-047 & El sistema deberá requerir la aceptación explícita del cliente antes de capturar el cobro. & Sí & M-09 Orquestador de Pedidos (OMS) / M-10 Atribución de Venta y Base de Comisión / M-11 Portal del Cliente & Subdoc.3 §3.2.6.3 Tabla 3.3 (SUP-11, SUP-12, SUP-13, SUP-14); §3.3.3 Tabla 3.9 \textbar{} Subdoc.4 §4.1.8.1 Flujo 2 \\
\addlinespace[2pt]
RF-048 & El sistema deberá anular la preautorización de pago del pedido. & Sí & M-09 Orquestador de Pedidos (OMS) / M-10 Atribución de Venta y Base de Comisión / M-11 Portal del Cliente & Subdoc.3 §3.2.6.3 Tabla 3.3 (SUP-11, SUP-12, SUP-13, SUP-14); §3.3.3 Tabla 3.9 \textbar{} Subdoc.4 §4.1.8.1 Flujo 2 \\
\addlinespace[2pt]
RF-049 & El sistema deberá cancelar el pedido registrando el motivo de la cancelación. & Sí & M-09 Orquestador de Pedidos (OMS) / M-10 Atribución de Venta y Base de Comisión / M-11 Portal del Cliente & Subdoc.3 §3.2.6.3 Tabla 3.3 (SUP-11, SUP-12, SUP-13, SUP-14); §3.3.3 Tabla 3.9 \textbar{} Subdoc.4 §4.1.8.1 Flujo 2 \\
\addlinespace[2pt]
RF-050 & El sistema deberá generar la conciliación diaria de preautorizaciones vencidas sin captura. & Sí & M-09 Orquestador de Pedidos (OMS) / M-10 Atribución de Venta y Base de Comisión / M-11 Portal del Cliente & Subdoc.3 §3.2.6.3 Tabla 3.3 (SUP-11, SUP-12, SUP-13, SUP-14); §3.3.3 Tabla 3.9 \textbar{} Subdoc.4 §4.1.8.1 Flujo 2 \\
\addlinespace[2pt]
RF-051 & El sistema deberá notificar al cliente el cambio de estado de su pedido antes de efectuar cualquier cobro definitivo. & Sí & M-09 Orquestador de Pedidos (OMS) / M-10 Atribución de Venta y Base de Comisión / M-11 Portal del Cliente & Subdoc.3 §3.2.6.3 Tabla 3.3 (SUP-11, SUP-12, SUP-13, SUP-14); §3.3.3 Tabla 3.9 \textbar{} Subdoc.4 §4.1.8.1 Flujo 2 \\
\addlinespace[2pt]
RF-052 & El sistema deberá permitir a todos los actores consultar el estado del pedido desde una única fuente de verdad. & Sí & M-09 Orquestador de Pedidos (OMS) / M-10 Atribución de Venta y Base de Comisión / M-11 Portal del Cliente & Subdoc.3 §3.2.6.3 Tabla 3.3 (SUP-11, SUP-12, SUP-13, SUP-14); §3.3.3 Tabla 3.9 \textbar{} Subdoc.4 §4.1.8.1 Flujo 2 \\
\addlinespace[2pt]
RF-053 & El sistema deberá permitir marcar como no elegible para cumplimiento digital el stock de exhibición de las categorías declaradas. & Sí & M-09 Orquestador de Pedidos (OMS) / M-10 Atribución de Venta y Base de Comisión / M-11 Portal del Cliente & Subdoc.3 §3.2.6.3 Tabla 3.3 (SUP-11, SUP-12, SUP-13, SUP-14); §3.3.3 Tabla 3.9 \textbar{} Subdoc.4 §4.1.8.1 Flujo 2 \\
\addlinespace[2pt]
RF-054 & El sistema debe calcular la base de comisión reconociendo al vendedor y a la tienda de origen de la unidad. & Sí & M-09 Orquestador de Pedidos (OMS) / M-10 Atribución de Venta y Base de Comisión / M-11 Portal del Cliente & Subdoc.3 §3.2.6.3 Tabla 3.3 (SUP-11, SUP-12, SUP-13, SUP-14); §3.3.3 Tabla 3.9 \textbar{} Subdoc.4 §4.1.8.1 Flujo 2 \\
\addlinespace[2pt]
RF-055 & El sistema debe condicionar la transmisión de la base de comisión al movimiento real de inventario verificado en bodega. & Sí & M-09 Orquestador de Pedidos (OMS) / M-10 Atribución de Venta y Base de Comisión / M-11 Portal del Cliente & Subdoc.3 §3.2.6.3 Tabla 3.3 (SUP-11, SUP-12, SUP-13, SUP-14); §3.3.3 Tabla 3.9 \textbar{} Subdoc.4 §4.1.8.1 Flujo 2 \\
\addlinespace[2pt]
RF-056 & El sistema deberá reasignar el pedido a otro punto de despacho con existencia disponible. & Sí & M-09 Orquestador de Pedidos (OMS) / M-10 Atribución de Venta y Base de Comisión / M-11 Portal del Cliente & Subdoc.3 §3.2.6.3 Tabla 3.3 (SUP-11, SUP-12, SUP-13, SUP-14); §3.3.3 Tabla 3.9 \textbar{} Subdoc.4 §4.1.8.1 Flujo 2 \\
\addlinespace[2pt]
RF-057 & El sistema deberá impedir una segunda reasignación del mismo pedido. & Sí & M-09 Orquestador de Pedidos (OMS) / M-10 Atribución de Venta y Base de Comisión / M-11 Portal del Cliente & Subdoc.3 §3.2.6.3 Tabla 3.3 (SUP-11, SUP-12, SUP-13, SUP-14); §3.3.3 Tabla 3.9 \textbar{} Subdoc.4 §4.1.8.1 Flujo 2 \\
\addlinespace[2pt]
RF-058 & El sistema deberá acotar la reasignación a la ventana de tiempo parametrizada. & Sí & M-09 Orquestador de Pedidos (OMS) / M-10 Atribución de Venta y Base de Comisión / M-11 Portal del Cliente & Subdoc.3 §3.2.6.3 Tabla 3.3 (SUP-11, SUP-12, SUP-13, SUP-14); §3.3.3 Tabla 3.9 \textbar{} Subdoc.4 §4.1.8.1 Flujo 2 \\
\addlinespace[2pt]
RF-059 & El sistema deberá ofrecer al cliente un producto equivalente. & Sí & M-09 Orquestador de Pedidos (OMS) / M-10 Atribución de Venta y Base de Comisión / M-11 Portal del Cliente & Subdoc.3 §3.2.6.3 Tabla 3.3 (SUP-11, SUP-12, SUP-13, SUP-14); §3.3.3 Tabla 3.9 \textbar{} Subdoc.4 §4.1.8.1 Flujo 2 \\
\addlinespace[2pt]
RF-060 & El sistema deberá ofrecer al cliente la espera compensada con la compensación parametrizada. & Sí & M-09 Orquestador de Pedidos (OMS) / M-10 Atribución de Venta y Base de Comisión / M-11 Portal del Cliente & Subdoc.3 §3.2.6.3 Tabla 3.3 (SUP-11, SUP-12, SUP-13, SUP-14); §3.3.3 Tabla 3.9 \textbar{} Subdoc.4 §4.1.8.1 Flujo 2 \\
\addlinespace[2pt]
RF-061 & El sistema deberá ofrecer al cliente la liberacion del pedido con anulación de la preautorización. & Sí & M-09 Orquestador de Pedidos (OMS) / M-10 Atribución de Venta y Base de Comisión / M-11 Portal del Cliente & Subdoc.3 §3.2.6.3 Tabla 3.3 (SUP-11, SUP-12, SUP-13, SUP-14); §3.3.3 Tabla 3.9 \textbar{} Subdoc.4 §4.1.8.1 Flujo 2 \\
\addlinespace[2pt]
RF-062 & El sistema debe monitorear el tiempo restante de cada pedido respecto de su fecha prometida de entrega. & Sí & M-09 Orquestador de Pedidos (OMS) / M-10 Atribución de Venta y Base de Comisión / M-11 Portal del Cliente & Subdoc.3 §3.2.6.3 Tabla 3.3 (SUP-11, SUP-12, SUP-13, SUP-14); §3.3.3 Tabla 3.9 \textbar{} Subdoc.4 §4.1.8.1 Flujo 2 \\
\addlinespace[2pt]
RF-063 & El sistema debe priorizar la preparación del pedido al alcanzar el umbral parametrizado de tiempo restante. & Sí & M-09 Orquestador de Pedidos (OMS) / M-10 Atribución de Venta y Base de Comisión / M-11 Portal del Cliente & Subdoc.3 §3.2.6.3 Tabla 3.3 (SUP-11, SUP-12, SUP-13, SUP-14); §3.3.3 Tabla 3.9 \textbar{} Subdoc.4 §4.1.8.1 Flujo 2 \\
\addlinespace[2pt]
RF-064 & El sistema debe permitir al cliente no autenticado consultar el precio vigente de una referencia. & Sí & M-09 Orquestador de Pedidos (OMS) / M-10 Atribución de Venta y Base de Comisión / M-11 Portal del Cliente & Subdoc.3 §3.2.6.3 Tabla 3.3 (SUP-11, SUP-12, SUP-13, SUP-14); §3.3.3 Tabla 3.9 \textbar{} Subdoc.4 §4.1.8.1 Flujo 2 \\
\addlinespace[2pt]
RF-065 & El sistema debe permitir al cliente no autenticado consultar la disponibilidad de la referencia por tienda. & Sí & M-09 Orquestador de Pedidos (OMS) / M-10 Atribución de Venta y Base de Comisión / M-11 Portal del Cliente & Subdoc.3 §3.2.6.3 Tabla 3.3 (SUP-11, SUP-12, SUP-13, SUP-14); §3.3.3 Tabla 3.9 \textbar{} Subdoc.4 §4.1.8.1 Flujo 2 \\
\addlinespace[2pt]
RF-066 & El sistema debe permitir al cliente no autenticado consultar la disponibilidad de la referencia para despacho. & Sí & M-09 Orquestador de Pedidos (OMS) / M-10 Atribución de Venta y Base de Comisión / M-11 Portal del Cliente & Subdoc.3 §3.2.6.3 Tabla 3.3 (SUP-11, SUP-12, SUP-13, SUP-14); §3.3.3 Tabla 3.9 \textbar{} Subdoc.4 §4.1.8.1 Flujo 2 \\
\addlinespace[2pt]
RF-067 & El sistema debe permitir al cliente autenticado consultar sus compras. & Sí & M-09 Orquestador de Pedidos (OMS) / M-10 Atribución de Venta y Base de Comisión / M-11 Portal del Cliente & Subdoc.3 §3.2.6.3 Tabla 3.3 (SUP-11, SUP-12, SUP-13, SUP-14); §3.3.3 Tabla 3.9 \textbar{} Subdoc.4 §4.1.8.1 Flujo 2 \\
\addlinespace[2pt]
RF-068 & El sistema debe permitir al cliente autenticado consultar sus devoluciones. & Sí & M-09 Orquestador de Pedidos (OMS) / M-10 Atribución de Venta y Base de Comisión / M-11 Portal del Cliente & Subdoc.3 §3.2.6.3 Tabla 3.3 (SUP-11, SUP-12, SUP-13, SUP-14); §3.3.3 Tabla 3.9 \textbar{} Subdoc.4 §4.1.8.1 Flujo 2 \\
\addlinespace[2pt]
RF-069 & El sistema debe permitir al cliente autenticado consultar su estado de cuenta. & Sí & M-09 Orquestador de Pedidos (OMS) / M-10 Atribución de Venta y Base de Comisión / M-11 Portal del Cliente & Subdoc.3 §3.2.6.3 Tabla 3.3 (SUP-11, SUP-12, SUP-13, SUP-14); §3.3.3 Tabla 3.9 \textbar{} Subdoc.4 §4.1.8.1 Flujo 2 \\
\addlinespace[2pt]
RF-070 & El sistema debe permitir al cliente autenticado consultar sus documentos. & Sí & M-09 Orquestador de Pedidos (OMS) / M-10 Atribución de Venta y Base de Comisión / M-11 Portal del Cliente & Subdoc.3 §3.2.6.3 Tabla 3.3 (SUP-11, SUP-12, SUP-13, SUP-14); §3.3.3 Tabla 3.9 \textbar{} Subdoc.4 §4.1.8.1 Flujo 2 \\
\addlinespace[2pt]
RF-071 & El sistema deberá calcular la base de comisión considerando el canal de origen y el canal de cumplimiento cuando estos difieran. & Sí & M-09 Orquestador de Pedidos (OMS) / M-10 Atribución de Venta y Base de Comisión / M-11 Portal del Cliente & Subdoc.3 §3.2.6.3 Tabla 3.3 (SUP-11, SUP-12, SUP-13, SUP-14); §3.3.3 Tabla 3.9 \textbar{} Subdoc.4 §4.1.8.1 Flujo 2 \\
\addlinespace[2pt]
RF-072 & El sistema debe permitir al ejecutivo de mesón consultar el estado real del pedido. & Sí & M-09 Orquestador de Pedidos (OMS) / M-10 Atribución de Venta y Base de Comisión / M-11 Portal del Cliente & Subdoc.3 §3.2.6.3 Tabla 3.3 (SUP-11, SUP-12, SUP-13, SUP-14); §3.3.3 Tabla 3.9 \textbar{} Subdoc.4 §4.1.8.1 Flujo 2 \\
\addlinespace[2pt]
RF-073 & El sistema debe permitir al ejecutivo de mesón consultar el precio aplicado al pedido. & Sí & M-09 Orquestador de Pedidos (OMS) / M-10 Atribución de Venta y Base de Comisión / M-11 Portal del Cliente & Subdoc.3 §3.2.6.3 Tabla 3.3 (SUP-11, SUP-12, SUP-13, SUP-14); §3.3.3 Tabla 3.9 \textbar{} Subdoc.4 §4.1.8.1 Flujo 2 \\
\addlinespace[2pt]
RF-074 & El sistema debe permitir al ejecutivo de mesón ofrecer al cliente las opciones de resolución definidas en RF-059, RF-060 y RF-061. & Sí & M-09 Orquestador de Pedidos (OMS) / M-10 Atribución de Venta y Base de Comisión / M-11 Portal del Cliente & Subdoc.3 §3.2.6.3 Tabla 3.3 (SUP-11, SUP-12, SUP-13, SUP-14); §3.3.3 Tabla 3.9 \textbar{} Subdoc.4 §4.1.8.1 Flujo 2 \\
\addlinespace[2pt]
RF-075 & El sistema deberá excluir al centro de distribución de Concepción como origen de promesa de entrega en el canal digital mientras no este acreditado su sistema de gestión de almacenes. & Sí & M-09 Orquestador de Pedidos (OMS) / M-10 Atribución de Venta y Base de Comisión / M-11 Portal del Cliente & Subdoc.3 §3.2.6.3 Tabla 3.3 (SUP-11, SUP-12, SUP-13, SUP-14); §3.3.3 Tabla 3.9 \textbar{} Subdoc.4 §4.1.8.1 Flujo 2 \\
\addlinespace[2pt]
RF-076 & El sistema deberá calcular el costo total de servir de cada punto de despacho candidato, incorporando costo logístico, costo de oportunidad de sala y plazo comprometido. & Sí & M-09 Orquestador de Pedidos (OMS) / M-10 Atribución de Venta y Base de Comisión / M-11 Portal del Cliente & Subdoc.3 §3.2.6.3 Tabla 3.3 (SUP-11, SUP-12, SUP-13, SUP-14); §3.3.3 Tabla 3.9 \textbar{} Subdoc.4 §4.1.8.1 Flujo 2 \\
\addlinespace[2pt]
RF-077 & El sistema deberá seleccionar como punto de despacho aquel de menor costo total de servir. & Sí & M-09 Orquestador de Pedidos (OMS) / M-10 Atribución de Venta y Base de Comisión / M-11 Portal del Cliente & Subdoc.3 §3.2.6.3 Tabla 3.3 (SUP-11, SUP-12, SUP-13, SUP-14); §3.3.3 Tabla 3.9 \textbar{} Subdoc.4 §4.1.8.1 Flujo 2 \\
\addlinespace[2pt]
RF-078 & El sistema deberá registrar el valor de cada término del costo total de servir que fundamento la seleccion. & Sí & M-09 Orquestador de Pedidos (OMS) / M-10 Atribución de Venta y Base de Comisión / M-11 Portal del Cliente & Subdoc.3 §3.2.6.3 Tabla 3.3 (SUP-11, SUP-12, SUP-13, SUP-14); §3.3.3 Tabla 3.9 \textbar{} Subdoc.4 §4.1.8.1 Flujo 2 \\
\addlinespace[2pt]
RF-079 & El sistema deberá calcular la fecha prometida de entrega en función del punto de despacho, la capacidad de preparacion y el transportista asignado. & Sí & M-09 Orquestador de Pedidos (OMS) / M-10 Atribución de Venta y Base de Comisión / M-11 Portal del Cliente & Subdoc.3 §3.2.6.3 Tabla 3.3 (SUP-11, SUP-12, SUP-13, SUP-14); §3.3.3 Tabla 3.9 \textbar{} Subdoc.4 §4.1.8.1 Flujo 2 \\
\addlinespace[2pt]
RF-080 & El sistema deberá impedir la publicación de un plazo fijo de catálogo como fecha prometida de entrega. & Sí & M-09 Orquestador de Pedidos (OMS) / M-10 Atribución de Venta y Base de Comisión / M-11 Portal del Cliente & Subdoc.3 §3.2.6.3 Tabla 3.3 (SUP-11, SUP-12, SUP-13, SUP-14); §3.3.3 Tabla 3.9 \textbar{} Subdoc.4 §4.1.8.1 Flujo 2 \\
\addlinespace[2pt]
RF-081 & El sistema deberá calcular el cumplimiento de la promesa de entrega por pedido individual. & Sí & M-09 Orquestador de Pedidos (OMS) / M-10 Atribución de Venta y Base de Comisión / M-11 Portal del Cliente & Subdoc.3 §3.2.6.3 Tabla 3.3 (SUP-11, SUP-12, SUP-13, SUP-14); §3.3.3 Tabla 3.9 \textbar{} Subdoc.4 §4.1.8.1 Flujo 2 \\
\addlinespace[2pt]
RF-082 & El sistema deberá detectar la pérdida del enlace externo y conmutar automáticamente la tienda a modo desconectado. & Sí & M-08 Punto de Venta y Nodo de Tienda / M-17 Contingencia de Crédito Offline & Subdoc.3 §3.2.4 (SUP-07); §3.3.3 Tabla 3.9 (M-08) \textbar{} Subdoc.4 §4.2.3 Nodo de tienda, Tabla 4.31 \\
\addlinespace[2pt]
RF-083 & El sistema deberá permitir al cajero registrar una venta en modo desconectado. & Sí & M-08 Punto de Venta y Nodo de Tienda / M-17 Contingencia de Crédito Offline & Subdoc.3 §3.2.4 (SUP-07); §3.3.3 Tabla 3.9 (M-08) \textbar{} Subdoc.4 §4.2.3 Nodo de tienda, Tabla 4.31 \\
\addlinespace[2pt]
RF-084 & El sistema deberá permitir al cajero cobrar la venta en modo desconectado. & Sí & M-08 Punto de Venta y Nodo de Tienda / M-17 Contingencia de Crédito Offline & Subdoc.3 §3.2.4 (SUP-07); §3.3.3 Tabla 3.9 (M-08) \textbar{} Subdoc.4 §4.2.3 Nodo de tienda, Tabla 4.31 \\
\addlinespace[2pt]
RF-085 & El sistema deberá aplicar las promociones vigentes en modo desconectado, evaluadas según RF-032, RF-033 y RF-034. & Sí & M-08 Punto de Venta y Nodo de Tienda / M-17 Contingencia de Crédito Offline & Subdoc.3 §3.2.4 (SUP-07); §3.3.3 Tabla 3.9 (M-08) \textbar{} Subdoc.4 §4.2.3 Nodo de tienda, Tabla 4.31 \\
\addlinespace[2pt]
RF-086 & El sistema deberá emitir el documento de venta en contingencia utilizando folios previamente asignados por el sistema de gestión empresarial. & Sí & M-08 Punto de Venta y Nodo de Tienda / M-17 Contingencia de Crédito Offline & Subdoc.3 §3.2.4 (SUP-07); §3.3.3 Tabla 3.9 (M-08) \textbar{} Subdoc.4 §4.2.3 Nodo de tienda, Tabla 4.31 \\
\addlinespace[2pt]
RF-087 & El sistema deberá permitir al vendedor consultar la existencia local de la tienda en modo desconectado. & Sí & M-08 Punto de Venta y Nodo de Tienda / M-17 Contingencia de Crédito Offline & Subdoc.3 §3.2.4 (SUP-07); §3.3.3 Tabla 3.9 (M-08) \textbar{} Subdoc.4 §4.2.3 Nodo de tienda, Tabla 4.31 \\
\addlinespace[2pt]
RF-088 & El sistema deberá detectar el restablecimiento del enlace y conmutar la tienda a modo conectado. & Sí & M-08 Punto de Venta y Nodo de Tienda / M-17 Contingencia de Crédito Offline & Subdoc.3 §3.2.4 (SUP-07); §3.3.3 Tabla 3.9 (M-08) \textbar{} Subdoc.4 §4.2.3 Nodo de tienda, Tabla 4.31 \\
\addlinespace[2pt]
RF-089 & El sistema deberá permitir el otorgamiento de crédito en modo desconectado exclusivamente contra cupo preaprobado vigente almacenado en caché local. & Sí & M-08 Punto de Venta y Nodo de Tienda / M-17 Contingencia de Crédito Offline & Subdoc.3 §3.2.4 (SUP-07); §3.3.3 Tabla 3.9 (M-08) \textbar{} Subdoc.4 §4.2.3 Nodo de tienda, Tabla 4.31 \\
\addlinespace[2pt]
RF-090 & El sistema deberá aplicar el tope de monto parametrizado a cada operación de crédito cursada en modo desconectado. & Sí & M-08 Punto de Venta y Nodo de Tienda / M-17 Contingencia de Crédito Offline & Subdoc.3 §3.2.4 (SUP-07); §3.3.3 Tabla 3.9 (M-08) \textbar{} Subdoc.4 §4.2.3 Nodo de tienda, Tabla 4.31 \\
\addlinespace[2pt]
RF-091 & El sistema deberá aplicar el tope parametrizado de número de operaciones de crédito en modo desconectado. & Sí & M-08 Punto de Venta y Nodo de Tienda / M-17 Contingencia de Crédito Offline & Subdoc.3 §3.2.4 (SUP-07); §3.3.3 Tabla 3.9 (M-08) \textbar{} Subdoc.4 §4.2.3 Nodo de tienda, Tabla 4.31 \\
\addlinespace[2pt]
RF-092 & El sistema deberá impedir la apertura de una tarjeta nueva en modo desconectado. & Sí & M-08 Punto de Venta y Nodo de Tienda / M-17 Contingencia de Crédito Offline & Subdoc.3 §3.2.4 (SUP-07); §3.3.3 Tabla 3.9 (M-08) \textbar{} Subdoc.4 §4.2.3 Nodo de tienda, Tabla 4.31 \\
\addlinespace[2pt]
RF-093 & El sistema deberá impedir la ampliación de cupo en modo desconectado. & Sí & M-08 Punto de Venta y Nodo de Tienda / M-17 Contingencia de Crédito Offline & Subdoc.3 §3.2.4 (SUP-07); §3.3.3 Tabla 3.9 (M-08) \textbar{} Subdoc.4 §4.2.3 Nodo de tienda, Tabla 4.31 \\
\addlinespace[2pt]
RF-094 & El sistema deberá marcar la operación como cursada en modo desconectado para su validación posterior. & Sí & M-08 Punto de Venta y Nodo de Tienda / M-17 Contingencia de Crédito Offline & Subdoc.3 §3.2.4 (SUP-07); §3.3.3 Tabla 3.9 (M-08) \textbar{} Subdoc.4 §4.2.3 Nodo de tienda, Tabla 4.31 \\
\addlinespace[2pt]
RF-095 & El sistema debe reconciliar hacia los sistemas centrales la totalidad de las ventas registradas en modo desconectado. & Sí & M-08 Punto de Venta y Nodo de Tienda / M-17 Contingencia de Crédito Offline & Subdoc.3 §3.2.4 (SUP-07); §3.3.3 Tabla 3.9 (M-08) \textbar{} Subdoc.4 §4.2.3 Nodo de tienda, Tabla 4.31 \\
\addlinespace[2pt]
RF-096 & El sistema debe reconciliar hacia los sistemas centrales la totalidad de los documentos emitidos en modo desconectado. & Sí & M-08 Punto de Venta y Nodo de Tienda / M-17 Contingencia de Crédito Offline & Subdoc.3 §3.2.4 (SUP-07); §3.3.3 Tabla 3.9 (M-08) \textbar{} Subdoc.4 §4.2.3 Nodo de tienda, Tabla 4.31 \\
\addlinespace[2pt]
RF-097 & El sistema deberá procesar la reconciliación de forma idempotente, impidiendo la duplicación de ventas o documentos ante reintentos. & Sí & M-08 Punto de Venta y Nodo de Tienda / M-17 Contingencia de Crédito Offline & Subdoc.3 §3.2.4 (SUP-07); §3.3.3 Tabla 3.9 (M-08) \textbar{} Subdoc.4 §4.2.3 Nodo de tienda, Tabla 4.31 \\
\addlinespace[2pt]
RF-098 & El sistema deberá resolver el conflicto de existencia comprometida aplicando la regla de precedencia declarada, produciendo el mismo resultado ante repetición. & Sí & M-08 Punto de Venta y Nodo de Tienda / M-17 Contingencia de Crédito Offline & Subdoc.3 §3.2.4 (SUP-07); §3.3.3 Tabla 3.9 (M-08) \textbar{} Subdoc.4 §4.2.3 Nodo de tienda, Tabla 4.31 \\
\addlinespace[2pt]
RF-099 & El sistema deberá generar un informe de excepciones de la reconciliación, identificando cada conflicto resuelto y su regla aplicada. & Sí & M-08 Punto de Venta y Nodo de Tienda / M-17 Contingencia de Crédito Offline & Subdoc.3 §3.2.4 (SUP-07); §3.3.3 Tabla 3.9 (M-08) \textbar{} Subdoc.4 §4.2.3 Nodo de tienda, Tabla 4.31 \\
\addlinespace[2pt]
RF-100 & El sistema deberá enrutar la emisión de todo documento tributario hacia el sistema de gestión empresarial como único emisor. & Sí & M-08 Punto de Venta y Nodo de Tienda / M-17 Contingencia de Crédito Offline & Subdoc.3 §3.2.4 (SUP-07); §3.3.3 Tabla 3.9 (M-08) \textbar{} Subdoc.4 §4.2.3 Nodo de tienda, Tabla 4.31 \\
\addlinespace[2pt]
RF-101 & El sistema deberá permitir al vendedor de piso consultar la disponibilidad para vender de una referencia. & Sí & M-08 Punto de Venta y Nodo de Tienda / M-17 Contingencia de Crédito Offline & Subdoc.3 §3.2.4 (SUP-07); §3.3.3 Tabla 3.9 (M-08) \textbar{} Subdoc.4 §4.2.3 Nodo de tienda, Tabla 4.31 \\
\addlinespace[2pt]
RF-102 & El sistema debe permitir al vendedor de marketplace declarar la existencia de sus referencias. & Sí & M-12 Gobernanza de Marketplace (Seller Center) / M-14 Conciliación de Recobros B2B & Subdoc.3 §3.2.6.4 Tabla 3.4 (SUP-15, SUP-16); §3.3.4 Tabla 3.10 \textbar{} Subdoc.4 §4.1.4.2 Tabla 4.16 \\
\addlinespace[2pt]
RF-103 & El sistema debe permitir al vendedor de marketplace actualizar la existencia previamente declarada. & Sí & M-12 Gobernanza de Marketplace (Seller Center) / M-14 Conciliación de Recobros B2B & Subdoc.3 §3.2.6.4 Tabla 3.4 (SUP-15, SUP-16); §3.3.4 Tabla 3.10 \textbar{} Subdoc.4 §4.1.4.2 Tabla 4.16 \\
\addlinespace[2pt]
RNF-02 & La interfaz de declaración de existencia del vendedor de marketplace debe estar disponible de forma continua. & Sí & Nodo de tienda on-premise (Z2) + Nube Multi-AZ (Z1) + Site Principal/Secundario (Z3) & Subdoc.4 §4.2.3 (autonomía offline), §4.2.5 (Site Secundario, RTO/RPO), §4.2.7.2 (niveles de servicio) \\
\addlinespace[2pt]
RF-104 & El sistema deberá registrar cada actualización de existencia declarada con su fecha y hora. & Sí & M-12 Gobernanza de Marketplace (Seller Center) / M-14 Conciliación de Recobros B2B & Subdoc.3 §3.2.6.4 Tabla 3.4 (SUP-15, SUP-16); §3.3.4 Tabla 3.10 \textbar{} Subdoc.4 §4.1.4.2 Tabla 4.16 \\
\addlinespace[2pt]
RF-105 & El sistema deberá calcular los indicadores de nivel de servicio por vendedor externo según las reglas publicadas. & Sí & M-12 Gobernanza de Marketplace (Seller Center) / M-14 Conciliación de Recobros B2B & Subdoc.3 §3.2.6.4 Tabla 3.4 (SUP-15, SUP-16); §3.3.4 Tabla 3.10 \textbar{} Subdoc.4 §4.1.4.2 Tabla 4.16 \\
\addlinespace[2pt]
RF-106 & El sistema deberá notificar al vendedor de marketplace la recepción de la devolución en el momento en que se registra. & Sí & M-12 Gobernanza de Marketplace (Seller Center) / M-14 Conciliación de Recobros B2B & Subdoc.3 §3.2.6.4 Tabla 3.4 (SUP-15, SUP-16); §3.3.4 Tabla 3.10 \textbar{} Subdoc.4 §4.1.4.2 Tabla 4.16 \\
\addlinespace[2pt]
RF-107 & El sistema deberá permitir al ejecutivo registrar que parte asume el costo de la devolución. & Sí & M-12 Gobernanza de Marketplace (Seller Center) / M-14 Conciliación de Recobros B2B & Subdoc.3 §3.2.6.4 Tabla 3.4 (SUP-15, SUP-16); §3.3.4 Tabla 3.10 \textbar{} Subdoc.4 §4.1.4.2 Tabla 4.16 \\
\addlinespace[2pt]
RF-108 & El sistema deberá permitir al ejecutivo registrar la prestación entregada al cliente (reparación, reposición o devolución del precio). & Sí & M-12 Gobernanza de Marketplace (Seller Center) / M-14 Conciliación de Recobros B2B & Subdoc.3 §3.2.6.4 Tabla 3.4 (SUP-15, SUP-16); §3.3.4 Tabla 3.10 \textbar{} Subdoc.4 §4.1.4.2 Tabla 4.16 \\
\addlinespace[2pt]
RF-109 & El sistema deberá publicar únicamente la existencia declarada en la última actualización vigente del vendedor. & Sí & M-12 Gobernanza de Marketplace (Seller Center) / M-14 Conciliación de Recobros B2B & Subdoc.3 §3.2.6.4 Tabla 3.4 (SUP-15, SUP-16); §3.3.4 Tabla 3.10 \textbar{} Subdoc.4 §4.1.4.2 Tabla 4.16 \\
\addlinespace[2pt]
RF-110 & El sistema debe permitir al vendedor de marketplace consultar el estado de cada pedido intermediado. & Sí & M-12 Gobernanza de Marketplace (Seller Center) / M-14 Conciliación de Recobros B2B & Subdoc.3 §3.2.6.4 Tabla 3.4 (SUP-15, SUP-16); §3.3.4 Tabla 3.10 \textbar{} Subdoc.4 §4.1.4.2 Tabla 4.16 \\
\addlinespace[2pt]
RF-111 & El sistema debe permitir al vendedor de marketplace consultar las devoluciones de sus pedidos. & Sí & M-12 Gobernanza de Marketplace (Seller Center) / M-14 Conciliación de Recobros B2B & Subdoc.3 §3.2.6.4 Tabla 3.4 (SUP-15, SUP-16); §3.3.4 Tabla 3.10 \textbar{} Subdoc.4 §4.1.4.2 Tabla 4.16 \\
\addlinespace[2pt]
RF-112 & El sistema debe permitir al vendedor de marketplace consultar su evaluación de desempeño vigente. & Sí & M-12 Gobernanza de Marketplace (Seller Center) / M-14 Conciliación de Recobros B2B & Subdoc.3 §3.2.6.4 Tabla 3.4 (SUP-15, SUP-16); §3.3.4 Tabla 3.10 \textbar{} Subdoc.4 §4.1.4.2 Tabla 4.16 \\
\addlinespace[2pt]
RF-113 & El sistema debe identificar visiblemente al vendedor de la unidad en el catálogo, la ficha de producto y el proceso de compra. & Sí & M-12 Gobernanza de Marketplace (Seller Center) / M-14 Conciliación de Recobros B2B & Subdoc.3 §3.2.6.4 Tabla 3.4 (SUP-15, SUP-16); §3.3.4 Tabla 3.10 \textbar{} Subdoc.4 §4.1.4.2 Tabla 4.16 \\
\addlinespace[2pt]
RF-114 & El sistema debe identificar visiblemente a quien despacha la unidad en el catálogo, la ficha de producto y el proceso de compra. & Sí & M-12 Gobernanza de Marketplace (Seller Center) / M-14 Conciliación de Recobros B2B & Subdoc.3 §3.2.6.4 Tabla 3.4 (SUP-15, SUP-16); §3.3.4 Tabla 3.10 \textbar{} Subdoc.4 §4.1.4.2 Tabla 4.16 \\
\addlinespace[2pt]
RF-115 & El sistema debe mostrar las condiciones de devolución aplicables a la unidad. & Sí & M-12 Gobernanza de Marketplace (Seller Center) / M-14 Conciliación de Recobros B2B & Subdoc.3 §3.2.6.4 Tabla 3.4 (SUP-15, SUP-16); §3.3.4 Tabla 3.10 \textbar{} Subdoc.4 §4.1.4.2 Tabla 4.16 \\
\addlinespace[2pt]
RF-116 & El sistema debe mostrar las condiciones de garantía legal aplicables a la unidad. & Sí & M-12 Gobernanza de Marketplace (Seller Center) / M-14 Conciliación de Recobros B2B & Subdoc.3 §3.2.6.4 Tabla 3.4 (SUP-15, SUP-16); §3.3.4 Tabla 3.10 \textbar{} Subdoc.4 §4.1.4.2 Tabla 4.16 \\
\addlinespace[2pt]
RF-117 & El sistema deberá impedir que un pedido intermediado comprometa existencia propia de la compañía. & Sí & M-12 Gobernanza de Marketplace (Seller Center) / M-14 Conciliación de Recobros B2B & Subdoc.3 §3.2.6.4 Tabla 3.4 (SUP-15, SUP-16); §3.3.4 Tabla 3.10 \textbar{} Subdoc.4 §4.1.4.2 Tabla 4.16 \\
\addlinespace[2pt]
RF-118 & El sistema deberá impedir que una unidad de inventario propio sea asignada al cumplimiento de un pedido de marketplace. & Sí & M-12 Gobernanza de Marketplace (Seller Center) / M-14 Conciliación de Recobros B2B & Subdoc.3 §3.2.6.4 Tabla 3.4 (SUP-15, SUP-16); §3.3.4 Tabla 3.10 \textbar{} Subdoc.4 §4.1.4.2 Tabla 4.16 \\
\addlinespace[2pt]
RF-119 & El sistema deberá registrar el acuse de conocimiento de las reglas de evaluación por parte de cada vendedor externo. & Sí & M-12 Gobernanza de Marketplace (Seller Center) / M-14 Conciliación de Recobros B2B & Subdoc.3 §3.2.6.4 Tabla 3.4 (SUP-15, SUP-16); §3.3.4 Tabla 3.10 \textbar{} Subdoc.4 §4.1.4.2 Tabla 4.16 \\
\addlinespace[2pt]
RF-120 & El sistema deberá impedir la aplicación de una regla de evaluación a un vendedor que no tenga acuse de conocimiento previo de esa regla. & Sí & M-12 Gobernanza de Marketplace (Seller Center) / M-14 Conciliación de Recobros B2B & Subdoc.3 §3.2.6.4 Tabla 3.4 (SUP-15, SUP-16); §3.3.4 Tabla 3.10 \textbar{} Subdoc.4 §4.1.4.2 Tabla 4.16 \\
\addlinespace[2pt]
RF-121 & El sistema deberá determinar la consecuencia escalonada aplicable según la matriz de escalamiento parametrizada (advertencia, restricción de publicación, suspensión). & Sí & M-12 Gobernanza de Marketplace (Seller Center) / M-14 Conciliación de Recobros B2B & Subdoc.3 §3.2.6.4 Tabla 3.4 (SUP-15, SUP-16); §3.3.4 Tabla 3.10 \textbar{} Subdoc.4 §4.1.4.2 Tabla 4.16 \\
\addlinespace[2pt]
RF-122 & El sistema deberá requerir la ejecución de la sanción por el rol nominado en la matriz de escalamiento. & Sí & M-12 Gobernanza de Marketplace (Seller Center) / M-14 Conciliación de Recobros B2B & Subdoc.3 §3.2.6.4 Tabla 3.4 (SUP-15, SUP-16); §3.3.4 Tabla 3.10 \textbar{} Subdoc.4 §4.1.4.2 Tabla 4.16 \\
\addlinespace[2pt]
RF-123 & El sistema deberá registrar la trazabilidad de la sanción aplicada, indicando regla incumplida, consecuencia, rol ejecutor e instante. & Sí & M-12 Gobernanza de Marketplace (Seller Center) / M-14 Conciliación de Recobros B2B & Subdoc.3 §3.2.6.4 Tabla 3.4 (SUP-15, SUP-16); §3.3.4 Tabla 3.10 \textbar{} Subdoc.4 §4.1.4.2 Tabla 4.16 \\
\addlinespace[2pt]
RF-124 & El sistema deberá despublicar automáticamente la oferta cuyo stock declarado haya superado el plazo de vigencia sin actualización. & Sí & M-12 Gobernanza de Marketplace (Seller Center) / M-14 Conciliación de Recobros B2B & Subdoc.3 §3.2.6.4 Tabla 3.4 (SUP-15, SUP-16); §3.3.4 Tabla 3.10 \textbar{} Subdoc.4 §4.1.4.2 Tabla 4.16 \\
\addlinespace[2pt]
RF-125 & El sistema deberá calcular e informar al sistema de gestión empresarial la base de comisión de marketplace sobre la venta efectivamente cumplida & Sí & M-12 Gobernanza de Marketplace (Seller Center) / M-14 Conciliación de Recobros B2B & Subdoc.3 §3.2.6.4 Tabla 3.4 (SUP-15, SUP-16); §3.3.4 Tabla 3.10 \textbar{} Subdoc.4 §4.1.4.2 Tabla 4.16 \\
\addlinespace[2pt]
RF-126 & El sistema deberá notificar al sistema de gestión empresarial la anulación de la base de comisión calculada cuando ocurra una devolución o cancelación. & Sí & M-12 Gobernanza de Marketplace (Seller Center) / M-14 Conciliación de Recobros B2B & Subdoc.3 §3.2.6.4 Tabla 3.4 (SUP-15, SUP-16); §3.3.4 Tabla 3.10 \textbar{} Subdoc.4 §4.1.4.2 Tabla 4.16 \\
\addlinespace[2pt]
RF-127 & El sistema deberá calcular la existencia disponible para vender restando de la existencia registrada las reservas vigentes, el comprometido no despachado y el colchón de confianza aplicable. & Sí & M-01 Motor de Disponibilidad (ATP) / M-02 Gestión de Inventario, Conteo Cíclico y Merma / M-03 Maestro de Artículos / M-04 Gestión de Nodos Logísticos & Subdoc.3 §3.2.1 (SUP-01); §3.2.6.1 Tabla 3.1 (SUP-17, SUP-18, SUP-20); §3.3.1 Tabla 3.7 \textbar{} Subdoc.4 §4.2.9.2 Tabla 4.45 \\
\addlinespace[2pt]
RF-128 & El sistema deberá determinar el valor del colchón de confianza aplicable en función de la categoría del artículo. & Sí & M-01 Motor de Disponibilidad (ATP) / M-02 Gestión de Inventario, Conteo Cíclico y Merma / M-03 Maestro de Artículos / M-04 Gestión de Nodos Logísticos & Subdoc.3 §3.2.1 (SUP-01); §3.2.6.1 Tabla 3.1 (SUP-17, SUP-18, SUP-20); §3.3.1 Tabla 3.7 \textbar{} Subdoc.4 §4.2.9.2 Tabla 4.45 \\
\addlinespace[2pt]
RF-129 & El sistema deberá determinar el valor del colchón de confianza aplicable en función del punto de existencia. & Sí & M-01 Motor de Disponibilidad (ATP) / M-02 Gestión de Inventario, Conteo Cíclico y Merma / M-03 Maestro de Artículos / M-04 Gestión de Nodos Logísticos & Subdoc.3 §3.2.1 (SUP-01); §3.2.6.1 Tabla 3.1 (SUP-17, SUP-18, SUP-20); §3.3.1 Tabla 3.7 \textbar{} Subdoc.4 §4.2.9.2 Tabla 4.45 \\
\addlinespace[2pt]
RF-130 & El sistema deberá registrar la traza del cálculo del disponible, incluyendo el valor de cada término de la fórmula y la versión de parámetros aplicada. & Sí & M-01 Motor de Disponibilidad (ATP) / M-02 Gestión de Inventario, Conteo Cíclico y Merma / M-03 Maestro de Artículos / M-04 Gestión de Nodos Logísticos & Subdoc.3 §3.2.1 (SUP-01); §3.2.6.1 Tabla 3.1 (SUP-17, SUP-18, SUP-20); §3.3.1 Tabla 3.7 \textbar{} Subdoc.4 §4.2.9.2 Tabla 4.45 \\
\addlinespace[2pt]
RF-131 & El sistema deberá permitir parametrizar la frecuencia de conteo cíclico por categoría. & Sí & M-01 Motor de Disponibilidad (ATP) / M-02 Gestión de Inventario, Conteo Cíclico y Merma / M-03 Maestro de Artículos / M-04 Gestión de Nodos Logísticos & Subdoc.3 §3.2.1 (SUP-01); §3.2.6.1 Tabla 3.1 (SUP-17, SUP-18, SUP-20); §3.3.1 Tabla 3.7 \textbar{} Subdoc.4 §4.2.9.2 Tabla 4.45 \\
\addlinespace[2pt]
RF-132 & El sistema deberá permitir parametrizar el método de conteo por categoría. & Sí & M-01 Motor de Disponibilidad (ATP) / M-02 Gestión de Inventario, Conteo Cíclico y Merma / M-03 Maestro de Artículos / M-04 Gestión de Nodos Logísticos & Subdoc.3 §3.2.1 (SUP-01); §3.2.6.1 Tabla 3.1 (SUP-17, SUP-18, SUP-20); §3.3.1 Tabla 3.7 \textbar{} Subdoc.4 §4.2.9.2 Tabla 4.45 \\
\addlinespace[2pt]
RF-133 & El sistema deberá permitir parametrizar el criterio de gatillo de recuento extraordinario por categoría. & Sí & M-01 Motor de Disponibilidad (ATP) / M-02 Gestión de Inventario, Conteo Cíclico y Merma / M-03 Maestro de Artículos / M-04 Gestión de Nodos Logísticos & Subdoc.3 §3.2.1 (SUP-01); §3.2.6.1 Tabla 3.1 (SUP-17, SUP-18, SUP-20); §3.3.1 Tabla 3.7 \textbar{} Subdoc.4 §4.2.9.2 Tabla 4.45 \\
\addlinespace[2pt]
RF-134 & El sistema deberá generar la programación de conteos cíclicos según la frecuencia parametrizada de cada categoría. & Sí & M-01 Motor de Disponibilidad (ATP) / M-02 Gestión de Inventario, Conteo Cíclico y Merma / M-03 Maestro de Artículos / M-04 Gestión de Nodos Logísticos & Subdoc.3 §3.2.1 (SUP-01); §3.2.6.1 Tabla 3.1 (SUP-17, SUP-18, SUP-20); §3.3.1 Tabla 3.7 \textbar{} Subdoc.4 §4.2.9.2 Tabla 4.45 \\
\addlinespace[2pt]
RF-135 & El sistema debe permitir registrar el resultado del conteo cíclico. & Sí & M-01 Motor de Disponibilidad (ATP) / M-02 Gestión de Inventario, Conteo Cíclico y Merma / M-03 Maestro de Artículos / M-04 Gestión de Nodos Logísticos & Subdoc.3 §3.2.1 (SUP-01); §3.2.6.1 Tabla 3.1 (SUP-17, SUP-18, SUP-20); §3.3.1 Tabla 3.7 \textbar{} Subdoc.4 §4.2.9.2 Tabla 4.45 \\
\addlinespace[2pt]
RNF-03 & El registro del conteo cíclico no debe bloquear ni degradar la operación de venta de la sala. & Sí & Arquitectura orientada a eventos con presupuesto de latencia por componente (M-01, M-05, M-08, M-09, M-11) & Subdoc.4 §4.1.1.1 Tabla 4.11 (Latencia en operaciones) \\
\addlinespace[2pt]
RF-136 & El sistema deberá señalar automáticamente los pasillos o referencias que requieren recuento extraordinario. & Sí & M-01 Motor de Disponibilidad (ATP) / M-02 Gestión de Inventario, Conteo Cíclico y Merma / M-03 Maestro de Artículos / M-04 Gestión de Nodos Logísticos & Subdoc.3 §3.2.1 (SUP-01); §3.2.6.1 Tabla 3.1 (SUP-17, SUP-18, SUP-20); §3.3.1 Tabla 3.7 \textbar{} Subdoc.4 §4.2.9.2 Tabla 4.45 \\
\addlinespace[2pt]
RF-137 & El sistema deberá calcular la exactitud de inventario resultante por categoría. & Sí & M-01 Motor de Disponibilidad (ATP) / M-02 Gestión de Inventario, Conteo Cíclico y Merma / M-03 Maestro de Artículos / M-04 Gestión de Nodos Logísticos & Subdoc.3 §3.2.1 (SUP-01); §3.2.6.1 Tabla 3.1 (SUP-17, SUP-18, SUP-20); §3.3.1 Tabla 3.7 \textbar{} Subdoc.4 §4.2.9.2 Tabla 4.45 \\
\addlinespace[2pt]
RF-138 & El sistema deberá permitir clasificar cada diferencia detectada en uno de los componentes de merma definidos (pérdida física, daño no dado de baja, error de recepción, unidad mal ubicada, devolución mal reintegrada, error de digitación). & Sí & M-01 Motor de Disponibilidad (ATP) / M-02 Gestión de Inventario, Conteo Cíclico y Merma / M-03 Maestro de Artículos / M-04 Gestión de Nodos Logísticos & Subdoc.3 §3.2.1 (SUP-01); §3.2.6.1 Tabla 3.1 (SUP-17, SUP-18, SUP-20); §3.3.1 Tabla 3.7 \textbar{} Subdoc.4 §4.2.9.2 Tabla 4.45 \\
\addlinespace[2pt]
RF-139 & El sistema deberá impedir el cierre de un ajuste de inventario que no tenga asignado un componente de merma. & Sí & M-01 Motor de Disponibilidad (ATP) / M-02 Gestión de Inventario, Conteo Cíclico y Merma / M-03 Maestro de Artículos / M-04 Gestión de Nodos Logísticos & Subdoc.3 §3.2.1 (SUP-01); §3.2.6.1 Tabla 3.1 (SUP-17, SUP-18, SUP-20); §3.3.1 Tabla 3.7 \textbar{} Subdoc.4 §4.2.9.2 Tabla 4.45 \\
\addlinespace[2pt]
RF-140 & El sistema deberá cuantificar el monto y la proporción de cada componente de merma sobre el total del período. & Sí & M-01 Motor de Disponibilidad (ATP) / M-02 Gestión de Inventario, Conteo Cíclico y Merma / M-03 Maestro de Artículos / M-04 Gestión de Nodos Logísticos & Subdoc.3 §3.2.1 (SUP-01); §3.2.6.1 Tabla 3.1 (SUP-17, SUP-18, SUP-20); §3.3.1 Tabla 3.7 \textbar{} Subdoc.4 §4.2.9.2 Tabla 4.45 \\
\addlinespace[2pt]
RF-141 & El sistema deberá generar el informe mensual de merma cuya suma de componentes sea igual al total registrado en el período. & Sí & M-01 Motor de Disponibilidad (ATP) / M-02 Gestión de Inventario, Conteo Cíclico y Merma / M-03 Maestro de Artículos / M-04 Gestión de Nodos Logísticos & Subdoc.3 §3.2.1 (SUP-01); §3.2.6.1 Tabla 3.1 (SUP-17, SUP-18, SUP-20); §3.3.1 Tabla 3.7 \textbar{} Subdoc.4 §4.2.9.2 Tabla 4.45 \\
\addlinespace[2pt]
RF-142 & El sistema deberá exigir al analista comercial completar los atributos obligatorios de la referencia (categoría, dimensiones, fragilidad, volumen, restricción de despacho) antes de guardar el registro como apto para uso operativo. & Sí & M-01 Motor de Disponibilidad (ATP) / M-02 Gestión de Inventario, Conteo Cíclico y Merma / M-03 Maestro de Artículos / M-04 Gestión de Nodos Logísticos & Subdoc.3 §3.2.1 (SUP-01); §3.2.6.1 Tabla 3.1 (SUP-17, SUP-18, SUP-20); §3.3.1 Tabla 3.7 \textbar{} Subdoc.4 §4.2.9.2 Tabla 4.45 \\
\addlinespace[2pt]
RF-143 & El sistema deberá impedir la publicación de esa referencia en el canal digital mientras no cuente con los atributos obligatorios completos. & Sí & M-01 Motor de Disponibilidad (ATP) / M-02 Gestión de Inventario, Conteo Cíclico y Merma / M-03 Maestro de Artículos / M-04 Gestión de Nodos Logísticos & Subdoc.3 §3.2.1 (SUP-01); §3.2.6.1 Tabla 3.1 (SUP-17, SUP-18, SUP-20); §3.3.1 Tabla 3.7 \textbar{} Subdoc.4 §4.2.9.2 Tabla 4.45 \\
\addlinespace[2pt]
RF-144 & El sistema deberá generar automáticamente la propuesta diaria de reposición utilizando el disponible para vender calculado en RF-127. & Sí & M-01 Motor de Disponibilidad (ATP) / M-02 Gestión de Inventario, Conteo Cíclico y Merma / M-03 Maestro de Artículos / M-04 Gestión de Nodos Logísticos & Subdoc.3 §3.2.1 (SUP-01); §3.2.6.1 Tabla 3.1 (SUP-17, SUP-18, SUP-20); §3.3.1 Tabla 3.7 \textbar{} Subdoc.4 §4.2.9.2 Tabla 4.45 \\
\addlinespace[2pt]
RF-145 & El sistema deberá permitir al planificador de logística ajustar la propuesta de reposición antes de confirmar el envío. & Sí & M-01 Motor de Disponibilidad (ATP) / M-02 Gestión de Inventario, Conteo Cíclico y Merma / M-03 Maestro de Artículos / M-04 Gestión de Nodos Logísticos & Subdoc.3 §3.2.1 (SUP-01); §3.2.6.1 Tabla 3.1 (SUP-17, SUP-18, SUP-20); §3.3.1 Tabla 3.7 \textbar{} Subdoc.4 §4.2.9.2 Tabla 4.45 \\
\addlinespace[2pt]
RF-146 & El sistema deberá registrar el ajuste realizado indicando usuario, valor propuesto, valor confirmado e instante. & Sí & M-01 Motor de Disponibilidad (ATP) / M-02 Gestión de Inventario, Conteo Cíclico y Merma / M-03 Maestro de Artículos / M-04 Gestión de Nodos Logísticos & Subdoc.3 §3.2.1 (SUP-01); §3.2.6.1 Tabla 3.1 (SUP-17, SUP-18, SUP-20); §3.3.1 Tabla 3.7 \textbar{} Subdoc.4 §4.2.9.2 Tabla 4.45 \\
\addlinespace[2pt]
RF-147 & El sistema debe generar un reporte periódico de referencias con atributos obligatorios incompletos. & Sí & M-01 Motor de Disponibilidad (ATP) / M-02 Gestión de Inventario, Conteo Cíclico y Merma / M-03 Maestro de Artículos / M-04 Gestión de Nodos Logísticos & Subdoc.3 §3.2.1 (SUP-01); §3.2.6.1 Tabla 3.1 (SUP-17, SUP-18, SUP-20); §3.3.1 Tabla 3.7 \textbar{} Subdoc.4 §4.2.9.2 Tabla 4.45 \\
\addlinespace[2pt]
RF-148 & El sistema debe generar un reporte periódico de publicaciones fallidas indicando la causa del error. & Sí & M-01 Motor de Disponibilidad (ATP) / M-02 Gestión de Inventario, Conteo Cíclico y Merma / M-03 Maestro de Artículos / M-04 Gestión de Nodos Logísticos & Subdoc.3 §3.2.1 (SUP-01); §3.2.6.1 Tabla 3.1 (SUP-17, SUP-18, SUP-20); §3.3.1 Tabla 3.7 \textbar{} Subdoc.4 §4.2.9.2 Tabla 4.45 \\
\addlinespace[2pt]
RF-149 & El sistema debe permitir al gerente de logística definir el valor del colchón de confianza por categoría. & Sí & M-01 Motor de Disponibilidad (ATP) / M-02 Gestión de Inventario, Conteo Cíclico y Merma / M-03 Maestro de Artículos / M-04 Gestión de Nodos Logísticos & Subdoc.3 §3.2.1 (SUP-01); §3.2.6.1 Tabla 3.1 (SUP-17, SUP-18, SUP-20); §3.3.1 Tabla 3.7 \textbar{} Subdoc.4 §4.2.9.2 Tabla 4.45 \\
\addlinespace[2pt]
RF-150 & El sistema debe permitir al gerente de logística definir el valor del colchón de confianza por punto de existencia. & Sí & M-01 Motor de Disponibilidad (ATP) / M-02 Gestión de Inventario, Conteo Cíclico y Merma / M-03 Maestro de Artículos / M-04 Gestión de Nodos Logísticos & Subdoc.3 §3.2.1 (SUP-01); §3.2.6.1 Tabla 3.1 (SUP-17, SUP-18, SUP-20); §3.3.1 Tabla 3.7 \textbar{} Subdoc.4 §4.2.9.2 Tabla 4.45 \\
\addlinespace[2pt]
RF-151 & El sistema deberá registrar quien definio o modifico el colchón, cuando, el valor anterior y el valor nuevo. & Sí & M-01 Motor de Disponibilidad (ATP) / M-02 Gestión de Inventario, Conteo Cíclico y Merma / M-03 Maestro de Artículos / M-04 Gestión de Nodos Logísticos & Subdoc.3 §3.2.1 (SUP-01); §3.2.6.1 Tabla 3.1 (SUP-17, SUP-18, SUP-20); §3.3.1 Tabla 3.7 \textbar{} Subdoc.4 §4.2.9.2 Tabla 4.45 \\
\addlinespace[2pt]
RF-152 & El sistema deberá sugerir un valor de referencia del colchón calculado a partir del histórico de ventas y quiebres de la categoría. & Sí & M-01 Motor de Disponibilidad (ATP) / M-02 Gestión de Inventario, Conteo Cíclico y Merma / M-03 Maestro de Artículos / M-04 Gestión de Nodos Logísticos & Subdoc.3 §3.2.1 (SUP-01); §3.2.6.1 Tabla 3.1 (SUP-17, SUP-18, SUP-20); §3.3.1 Tabla 3.7 \textbar{} Subdoc.4 §4.2.9.2 Tabla 4.45 \\
\addlinespace[2pt]
RF-153 & El sistema deberá permitir al vendedor de piso registrar las unidades que ingresan al probador. & Sí & M-01 Motor de Disponibilidad (ATP) / M-02 Gestión de Inventario, Conteo Cíclico y Merma / M-03 Maestro de Artículos / M-04 Gestión de Nodos Logísticos & Subdoc.3 §3.2.1 (SUP-01); §3.2.6.1 Tabla 3.1 (SUP-17, SUP-18, SUP-20); §3.3.1 Tabla 3.7 \textbar{} Subdoc.4 §4.2.9.2 Tabla 4.45 \\
\addlinespace[2pt]
RF-154 & El sistema deberá excluir del disponible para vender las unidades registradas en probador. & Sí & M-01 Motor de Disponibilidad (ATP) / M-02 Gestión de Inventario, Conteo Cíclico y Merma / M-03 Maestro de Artículos / M-04 Gestión de Nodos Logísticos & Subdoc.3 §3.2.1 (SUP-01); §3.2.6.1 Tabla 3.1 (SUP-17, SUP-18, SUP-20); §3.3.1 Tabla 3.7 \textbar{} Subdoc.4 §4.2.9.2 Tabla 4.45 \\
\addlinespace[2pt]
RF-155 & El sistema debe permitir al vendedor de piso confirmar el reingreso a sala de la unidad registrada en probador. & Sí & M-01 Motor de Disponibilidad (ATP) / M-02 Gestión de Inventario, Conteo Cíclico y Merma / M-03 Maestro de Artículos / M-04 Gestión de Nodos Logísticos & Subdoc.3 §3.2.1 (SUP-01); §3.2.6.1 Tabla 3.1 (SUP-17, SUP-18, SUP-20); §3.3.1 Tabla 3.7 \textbar{} Subdoc.4 §4.2.9.2 Tabla 4.45 \\
\addlinespace[2pt]
RF-156 & El sistema debe permitir al vendedor de piso registrar la venta de la unidad registrada en probador. & Sí & M-01 Motor de Disponibilidad (ATP) / M-02 Gestión de Inventario, Conteo Cíclico y Merma / M-03 Maestro de Artículos / M-04 Gestión de Nodos Logísticos & Subdoc.3 §3.2.1 (SUP-01); §3.2.6.1 Tabla 3.1 (SUP-17, SUP-18, SUP-20); §3.3.1 Tabla 3.7 \textbar{} Subdoc.4 §4.2.9.2 Tabla 4.45 \\
\addlinespace[2pt]
RF-157 & El sistema deberá impedir que cualquier canal de venta consuma el saldo bruto de inventario para publicar o comprometer existencia. & Sí & M-01 Motor de Disponibilidad (ATP) / M-02 Gestión de Inventario, Conteo Cíclico y Merma / M-03 Maestro de Artículos / M-04 Gestión de Nodos Logísticos & Subdoc.3 §3.2.1 (SUP-01); §3.2.6.1 Tabla 3.1 (SUP-17, SUP-18, SUP-20); §3.3.1 Tabla 3.7 \textbar{} Subdoc.4 §4.2.9.2 Tabla 4.45 \\
\addlinespace[2pt]
RF-158 & El sistema debe calcular el porcentaje de error probable de la referencia en función de su categoría. & Sí & M-01 Motor de Disponibilidad (ATP) / M-02 Gestión de Inventario, Conteo Cíclico y Merma / M-03 Maestro de Artículos / M-04 Gestión de Nodos Logísticos & Subdoc.3 §3.2.1 (SUP-01); §3.2.6.1 Tabla 3.1 (SUP-17, SUP-18, SUP-20); §3.3.1 Tabla 3.7 \textbar{} Subdoc.4 §4.2.9.2 Tabla 4.45 \\
\addlinespace[2pt]
RF-159 & El sistema debe calcular el porcentaje de error probable de la referencia en función de su punto de existencia. & Sí & M-01 Motor de Disponibilidad (ATP) / M-02 Gestión de Inventario, Conteo Cíclico y Merma / M-03 Maestro de Artículos / M-04 Gestión de Nodos Logísticos & Subdoc.3 §3.2.1 (SUP-01); §3.2.6.1 Tabla 3.1 (SUP-17, SUP-18, SUP-20); §3.3.1 Tabla 3.7 \textbar{} Subdoc.4 §4.2.9.2 Tabla 4.45 \\
\addlinespace[2pt]
RF-160 & El sistema deberá mostrar al cliente el porcentaje de error probable calculado en RF-158 y RF-159 junto a la disponibilidad publicada. & Sí & M-01 Motor de Disponibilidad (ATP) / M-02 Gestión de Inventario, Conteo Cíclico y Merma / M-03 Maestro de Artículos / M-04 Gestión de Nodos Logísticos & Subdoc.3 §3.2.1 (SUP-01); §3.2.6.1 Tabla 3.1 (SUP-17, SUP-18, SUP-20); §3.3.1 Tabla 3.7 \textbar{} Subdoc.4 §4.2.9.2 Tabla 4.45 \\
\addlinespace[2pt]
RF-161 & El sistema deberá mostrar al vendedor de piso el porcentaje de error probable calculado en RF-158 y RF-159. & Sí & M-01 Motor de Disponibilidad (ATP) / M-02 Gestión de Inventario, Conteo Cíclico y Merma / M-03 Maestro de Artículos / M-04 Gestión de Nodos Logísticos & Subdoc.3 §3.2.1 (SUP-01); §3.2.6.1 Tabla 3.1 (SUP-17, SUP-18, SUP-20); §3.3.1 Tabla 3.7 \textbar{} Subdoc.4 §4.2.9.2 Tabla 4.45 \\
\addlinespace[2pt]
RF-162 & El sistema deberá generar una alerta interna dirigida al rol facultado cuando la exactitud de una categoría caiga bajo el umbral definido. & Sí & M-01 Motor de Disponibilidad (ATP) / M-02 Gestión de Inventario, Conteo Cíclico y Merma / M-03 Maestro de Artículos / M-04 Gestión de Nodos Logísticos & Subdoc.3 §3.2.1 (SUP-01); §3.2.6.1 Tabla 3.1 (SUP-17, SUP-18, SUP-20); §3.3.1 Tabla 3.7 \textbar{} Subdoc.4 §4.2.9.2 Tabla 4.45 \\
\addlinespace[2pt]
RF-163 & El sistema deberá permitir exclusivamente al rol facultado suspender manualmente la publicación de una categoría. & Sí & M-01 Motor de Disponibilidad (ATP) / M-02 Gestión de Inventario, Conteo Cíclico y Merma / M-03 Maestro de Artículos / M-04 Gestión de Nodos Logísticos & Subdoc.3 §3.2.1 (SUP-01); §3.2.6.1 Tabla 3.1 (SUP-17, SUP-18, SUP-20); §3.3.1 Tabla 3.7 \textbar{} Subdoc.4 §4.2.9.2 Tabla 4.45 \\
\addlinespace[2pt]
RF-164 & El sistema deberá abstenerse de suspender automáticamente la publicación de una categoría que no tenga una regla de degradación declarada y aprobada previamente. & Sí & M-01 Motor de Disponibilidad (ATP) / M-02 Gestión de Inventario, Conteo Cíclico y Merma / M-03 Maestro de Artículos / M-04 Gestión de Nodos Logísticos & Subdoc.3 §3.2.1 (SUP-01); §3.2.6.1 Tabla 3.1 (SUP-17, SUP-18, SUP-20); §3.3.1 Tabla 3.7 \textbar{} Subdoc.4 §4.2.9.2 Tabla 4.45 \\
\addlinespace[2pt]
RF-165 & El sistema deberá registrar cada cruce de información ejecutado entre ámbitos, indicando el dato cruzado, la finalidad, la base de licitud, la autorización nominada y el instante. & Sí & M-21 Gobierno de Frontera de Datos / M-22 Resolución de Identidad de Cliente & Subdoc.3 §3.2.3 (SUP-03, SUP-04); §3.3.6 Tabla 3.12 \textbar{} Subdoc.4 §4.1.2.2, §4.1.7 (frente 1) \\
\addlinespace[2pt]
RF-166 & El sistema deberá bloquear todo intento de cruce de información que no corresponda a una interfaz declarada en el inventario de flujos autorizados. & Sí & M-21 Gobierno de Frontera de Datos / M-22 Resolución de Identidad de Cliente & Subdoc.3 §3.2.3 (SUP-03, SUP-04); §3.3.6 Tabla 3.12 \textbar{} Subdoc.4 §4.1.2.2, §4.1.7 (frente 1) \\
\addlinespace[2pt]
RF-167 & El sistema deberá registrar el intento de cruce bloqueado con el componente origen, el dato solicitado y el instante. & Sí & M-21 Gobierno de Frontera de Datos / M-22 Resolución de Identidad de Cliente & Subdoc.3 §3.2.3 (SUP-03, SUP-04); §3.3.6 Tabla 3.12 \textbar{} Subdoc.4 §4.1.2.2, §4.1.7 (frente 1) \\
\addlinespace[2pt]
RF-168 & El sistema debe permitir al oficial de cumplimiento mantener el inventario de interfaces de cruce declaradas. & Sí & M-21 Gobierno de Frontera de Datos / M-22 Resolución de Identidad de Cliente & Subdoc.3 §3.2.3 (SUP-03, SUP-04); §3.3.6 Tabla 3.12 \textbar{} Subdoc.4 §4.1.2.2, §4.1.7 (frente 1) \\
\addlinespace[2pt]
RF-169 & El sistema debe impedir el registro de una interfaz de cruce que no declare finalidad, base de licitud y autorización nominada. & Sí & M-21 Gobierno de Frontera de Datos / M-22 Resolución de Identidad de Cliente & Subdoc.3 §3.2.3 (SUP-03, SUP-04); §3.3.6 Tabla 3.12 \textbar{} Subdoc.4 §4.1.2.2, §4.1.7 (frente 1) \\
\addlinespace[2pt]
RF-170 & El sistema deberá excluir del catálogo de atributos disponibles en el motor de campañas todo atributo de origen financiero (mora, deuda, cupo utilizado, comportamiento de pago). & Sí & M-21 Gobierno de Frontera de Datos / M-22 Resolución de Identidad de Cliente & Subdoc.3 §3.2.3 (SUP-03, SUP-04); §3.3.6 Tabla 3.12 \textbar{} Subdoc.4 §4.1.2.2, §4.1.7 (frente 1) \\
\addlinespace[2pt]
RF-171 & El sistema deberá rechazar la ejecución de toda facilidad comercial que invoque un atributo de origen financiero sin un registro de cruce autorizado asociado. & Sí & M-21 Gobierno de Frontera de Datos / M-22 Resolución de Identidad de Cliente & Subdoc.3 §3.2.3 (SUP-03, SUP-04); §3.3.6 Tabla 3.12 \textbar{} Subdoc.4 §4.1.2.2, §4.1.7 (frente 1) \\
\addlinespace[2pt]
RF-172 & El sistema deberá rechazar la ejecución de todo proceso crediticio que invoque un atributo de origen retail sin un registro de cruce autorizado asociado. & Sí & M-21 Gobierno de Frontera de Datos / M-22 Resolución de Identidad de Cliente & Subdoc.3 §3.2.3 (SUP-03, SUP-04); §3.3.6 Tabla 3.12 \textbar{} Subdoc.4 §4.1.2.2, §4.1.7 (frente 1) \\
\addlinespace[2pt]
RF-173 & El sistema deberá asignar a la misma persona un identificador distinto en el ámbito retail y en el ámbito fiscalizado. & Sí & M-21 Gobierno de Frontera de Datos / M-22 Resolución de Identidad de Cliente & Subdoc.3 §3.2.3 (SUP-03, SUP-04); §3.3.6 Tabla 3.12 \textbar{} Subdoc.4 §4.1.2.2, §4.1.7 (frente 1) \\
\addlinespace[2pt]
RF-174 & El sistema deberá resolver la correspondencia entre identificadores exclusivamente a través de una tabla de correspondencia custodiada, con acceso nominado y registrado. & Sí & M-21 Gobierno de Frontera de Datos / M-22 Resolución de Identidad de Cliente & Subdoc.3 §3.2.3 (SUP-03, SUP-04); §3.3.6 Tabla 3.12 \textbar{} Subdoc.4 §4.1.2.2, §4.1.7 (frente 1) \\
\addlinespace[2pt]
RF-175 & El sistema deberá rechazar la persistencia de la clave del ámbito financiero como columna o atributo en cualquier entidad del ámbito retail. & Sí & M-21 Gobierno de Frontera de Datos / M-22 Resolución de Identidad de Cliente & Subdoc.3 §3.2.3 (SUP-03, SUP-04); §3.3.6 Tabla 3.12 \textbar{} Subdoc.4 §4.1.2.2, §4.1.7 (frente 1) \\
\addlinespace[2pt]
RF-176 & El sistema deberá exigir el registro de la evaluación de impacto sobre la frontera de datos antes de habilitar el estado 'aprobada' de una iniciativa. & Sí & M-21 Gobierno de Frontera de Datos / M-22 Resolución de Identidad de Cliente & Subdoc.3 §3.2.3 (SUP-03, SUP-04); §3.3.6 Tabla 3.12 \textbar{} Subdoc.4 §4.1.2.2, §4.1.7 (frente 1) \\
\addlinespace[2pt]
RF-177 & El sistema debe permitir declarar previamente el orden de degradación de los servicios. & Sí & M-24 Observabilidad y Degradación Controlada & Subdoc.3 §3.2.6.6 Tabla 3.6 (SUP-23); §3.3.6 Tabla 3.12 \textbar{} Subdoc.4 §4.1.10 Política de escalamiento \\
\addlinespace[2pt]
RF-178 & El sistema debe permitir parametrizar los criterios que activan cada nivel de degradación declarado. & Sí & M-24 Observabilidad y Degradación Controlada & Subdoc.3 §3.2.6.6 Tabla 3.6 (SUP-23); §3.3.6 Tabla 3.12 \textbar{} Subdoc.4 §4.1.10 Política de escalamiento \\
\addlinespace[2pt]
RF-179 & El sistema deberá suspender la publicación de la categoría indicada por el orden de degradación declarado. & Sí & M-24 Observabilidad y Degradación Controlada & Subdoc.3 §3.2.6.6 Tabla 3.6 (SUP-23); §3.3.6 Tabla 3.12 \textbar{} Subdoc.4 §4.1.10 Política de escalamiento \\
\addlinespace[2pt]
RF-180 & El sistema deberá reducir el límite de unidades por cliente al valor declarado. & Sí & M-24 Observabilidad y Degradación Controlada & Subdoc.3 §3.2.6.6 Tabla 3.6 (SUP-23); §3.3.6 Tabla 3.12 \textbar{} Subdoc.4 §4.1.10 Política de escalamiento \\
\addlinespace[2pt]
RF-181 & El sistema deberá desactivar temporalmente los medios de pago de mayor fricción operativa declarados. & Sí & M-24 Observabilidad y Degradación Controlada & Subdoc.3 §3.2.6.6 Tabla 3.6 (SUP-23); §3.3.6 Tabla 3.12 \textbar{} Subdoc.4 §4.1.10 Política de escalamiento \\
\addlinespace[2pt]
RF-182 & El sistema deberá monitorear en tiempo real la tasa de cancelación por falta de existencia por categoría. & Sí & M-24 Observabilidad y Degradación Controlada & Subdoc.3 §3.2.6.6 Tabla 3.6 (SUP-23); §3.3.6 Tabla 3.12 \textbar{} Subdoc.4 §4.1.10 Política de escalamiento \\
\addlinespace[2pt]
RF-183 & El sistema deberá activar la acción de degradación declarada para esa categoría en RF-026.1. & Sí & M-24 Observabilidad y Degradación Controlada & Subdoc.3 §3.2.6.6 Tabla 3.6 (SUP-23); §3.3.6 Tabla 3.12 \textbar{} Subdoc.4 §4.1.10 Política de escalamiento \\
\addlinespace[2pt]
RF-184 & El sistema deberá permitir parametrizar las cinco ventanas de congelamiento (1 nov - 6 ene; evento anual y su semana previa; evento de noviembre; semana del Día de la Madre; última semana de enero a primera de marzo). & Sí & M-24 Observabilidad y Degradación Controlada & Subdoc.3 §3.2.6.6 Tabla 3.6 (SUP-23); §3.3.6 Tabla 3.12 \textbar{} Subdoc.4 §4.1.10 Política de escalamiento \\
\addlinespace[2pt]
RF-185 & El sistema deberá impedir la ejecución del despliegue o intervención en producción durante la ventana de congelamiento. & Sí & M-24 Observabilidad y Degradación Controlada & Subdoc.3 §3.2.6.6 Tabla 3.6 (SUP-23); §3.3.6 Tabla 3.12 \textbar{} Subdoc.4 §4.1.10 Política de escalamiento \\
\addlinespace[2pt]
RF-186 & El sistema deberá registrar el intento de despliegue bloqueado, indicando componente, solicitante e instante. & Sí & M-24 Observabilidad y Degradación Controlada & Subdoc.3 §3.2.6.6 Tabla 3.6 (SUP-23); §3.3.6 Tabla 3.12 \textbar{} Subdoc.4 §4.1.10 Política de escalamiento \\
\addlinespace[2pt]
RF-187 & El sistema debe permitir al ejecutivo de mesón atender íntegramente el caso de garantía legal en la tienda. & Sí & M-13 Postventa y Logística Inversa & Subdoc.3 §3.2.6.4 Tabla 3.4 (SUP-21, SUP-22); §3.3.4 Tabla 3.10 \\
\addlinespace[2pt]
RF-188 & El sistema debe impedir que el flujo de atención exija la derivación del consumidor al fabricante, al servicio técnico o al vendedor de marketplace como condición para ser atendido. & Sí & M-13 Postventa y Logística Inversa & Subdoc.3 §3.2.6.4 Tabla 3.4 (SUP-21, SUP-22); §3.3.4 Tabla 3.10 \\
\addlinespace[2pt]
RF-189 & El sistema deberá permitir al ejecutivo registrar el motivo de la devolución (talla, color, expectativa, falla). & Sí & M-13 Postventa y Logística Inversa & Subdoc.3 §3.2.6.4 Tabla 3.4 (SUP-21, SUP-22); §3.3.4 Tabla 3.10 \\
\addlinespace[2pt]
RF-190 & El sistema deberá impedir el reingreso de una unidad devuelta al inventario disponible mientras no exista una decisión de aptitud registrada. & Sí & M-13 Postventa y Logística Inversa & Subdoc.3 §3.2.6.4 Tabla 3.4 (SUP-21, SUP-22); §3.3.4 Tabla 3.10 \\
\addlinespace[2pt]
RF-191 & El sistema deberá permitir al ejecutivo registrar la decisión de aptitud de la unidad devuelta, identificando al responsable de la decisión y su instante. & Sí & M-13 Postventa y Logística Inversa & Subdoc.3 §3.2.6.4 Tabla 3.4 (SUP-21, SUP-22); §3.3.4 Tabla 3.10 \\
\addlinespace[2pt]
RF-192 & El sistema deberá presentar al ejecutivo las tres opciones de garantía legal (reparación, reposición y devolución del precio) para que el consumidor elija. & Sí & M-13 Postventa y Logística Inversa & Subdoc.3 §3.2.6.4 Tabla 3.4 (SUP-21, SUP-22); §3.3.4 Tabla 3.10 \\
\addlinespace[2pt]
RF-193 & El sistema debe registrar la opción de garantía legal ofrecida al consumidor. & Sí & M-13 Postventa y Logística Inversa & Subdoc.3 §3.2.6.4 Tabla 3.4 (SUP-21, SUP-22); §3.3.4 Tabla 3.10 \\
\addlinespace[2pt]
RF-194 & El sistema debe registrar la opción de garantía legal efectivamente elegida por el consumidor. & Sí & M-13 Postventa y Logística Inversa & Subdoc.3 §3.2.6.4 Tabla 3.4 (SUP-21, SUP-22); §3.3.4 Tabla 3.10 \\
\addlinespace[2pt]
RF-195 & El sistema deberá permitir parametrizar el plazo de garantía legal por tipo de producto, sin requerir modificacion de código. & Sí & M-13 Postventa y Logística Inversa & Subdoc.3 §3.2.6.4 Tabla 3.4 (SUP-21, SUP-22); §3.3.4 Tabla 3.10 \\
\addlinespace[2pt]
RF-196 & El sistema deberá registrar el hito de resolución al consumidor con su fecha propia. & Sí & M-13 Postventa y Logística Inversa & Subdoc.3 §3.2.6.4 Tabla 3.4 (SUP-21, SUP-22); §3.3.4 Tabla 3.10 \\
\addlinespace[2pt]
RF-197 & El sistema deberá registrar el hito de recuperación contra el tercero responsable con una fecha independiente de la resolución al consumidor. & Sí & M-13 Postventa y Logística Inversa & Subdoc.3 §3.2.6.4 Tabla 3.4 (SUP-21, SUP-22); §3.3.4 Tabla 3.10 \\
\addlinespace[2pt]
RF-198 & El sistema deberá impedir que el estado de la recuperación contra el tercero condicione el cierre de la resolución al consumidor. & Sí & M-13 Postventa y Logística Inversa & Subdoc.3 §3.2.6.4 Tabla 3.4 (SUP-21, SUP-22); §3.3.4 Tabla 3.10 \\
\addlinespace[2pt]
RF-199 & El sistema deberá exigir la entrega completa de la información precontractual antes de habilitar la evaluación de la solicitud de crédito. & Sí & M-15 Originación y Evaluación Crediticia / M-16 Gestor de Evidencia y Consentimiento / M-18 Administración de Cartera, Cobranza y Repactación / M-19 Migración y Conciliación de Cartera & Subdoc.3 §3.2.5 (SUP-25); §3.2.6.5 Tabla 3.5 (SUP-05, SUP-06); §3.3.5 Tabla 3.11 \textbar{} Subdoc.4 §4.1.5.2, §4.1.7 \\
\addlinespace[2pt]
RF-200 & El sistema deberá registrar la versión del documento precontractual entregado. & Sí & M-15 Originación y Evaluación Crediticia / M-16 Gestor de Evidencia y Consentimiento / M-18 Administración de Cartera, Cobranza y Repactación / M-19 Migración y Conciliación de Cartera & Subdoc.3 §3.2.5 (SUP-25); §3.2.6.5 Tabla 3.5 (SUP-05, SUP-06); §3.3.5 Tabla 3.11 \textbar{} Subdoc.4 §4.1.5.2, §4.1.7 \\
\addlinespace[2pt]
RF-201 & El sistema deberá registrar el instante de entrega de la información precontractual. & Sí & M-15 Originación y Evaluación Crediticia / M-16 Gestor de Evidencia y Consentimiento / M-18 Administración de Cartera, Cobranza y Repactación / M-19 Migración y Conciliación de Cartera & Subdoc.3 §3.2.5 (SUP-25); §3.2.6.5 Tabla 3.5 (SUP-05, SUP-06); §3.3.5 Tabla 3.11 \textbar{} Subdoc.4 §4.1.5.2, §4.1.7 \\
\addlinespace[2pt]
RF-202 & El sistema deberá registrar el instante de aceptación del cliente. & Sí & M-15 Originación y Evaluación Crediticia / M-16 Gestor de Evidencia y Consentimiento / M-18 Administración de Cartera, Cobranza y Repactación / M-19 Migración y Conciliación de Cartera & Subdoc.3 §3.2.5 (SUP-25); §3.2.6.5 Tabla 3.5 (SUP-05, SUP-06); §3.3.5 Tabla 3.11 \textbar{} Subdoc.4 §4.1.5.2, §4.1.7 \\
\addlinespace[2pt]
RF-203 & El sistema deberá registrar el medio por el cual se entrego y se aceptó la información precontractual. & Sí & M-15 Originación y Evaluación Crediticia / M-16 Gestor de Evidencia y Consentimiento / M-18 Administración de Cartera, Cobranza y Repactación / M-19 Migración y Conciliación de Cartera & Subdoc.3 §3.2.5 (SUP-25); §3.2.6.5 Tabla 3.5 (SUP-05, SUP-06); §3.3.5 Tabla 3.11 \textbar{} Subdoc.4 §4.1.5.2, §4.1.7 \\
\addlinespace[2pt]
RF-204 & El sistema deberá registrar el contenido exacto aceptado por el cliente o su huella de integridad verificable. & Sí & M-15 Originación y Evaluación Crediticia / M-16 Gestor de Evidencia y Consentimiento / M-18 Administración de Cartera, Cobranza y Repactación / M-19 Migración y Conciliación de Cartera & Subdoc.3 §3.2.5 (SUP-25); §3.2.6.5 Tabla 3.5 (SUP-05, SUP-06); §3.3.5 Tabla 3.11 \textbar{} Subdoc.4 §4.1.5.2, §4.1.7 \\
\addlinespace[2pt]
RF-205 & El sistema deberá dejar constancia expresa de que el cliente recibió la información antes de aceptar, mediante un mecanismo distinto de la sola firma en papel archivado. & Sí & M-15 Originación y Evaluación Crediticia / M-16 Gestor de Evidencia y Consentimiento / M-18 Administración de Cartera, Cobranza y Repactación / M-19 Migración y Conciliación de Cartera & Subdoc.3 §3.2.5 (SUP-25); §3.2.6.5 Tabla 3.5 (SUP-05, SUP-06); §3.3.5 Tabla 3.11 \textbar{} Subdoc.4 §4.1.5.2, §4.1.7 \\
\addlinespace[2pt]
RF-206 & El sistema deberá impedir el registro de la aceptación del crédito mientras no exista acreditación de entrega previa de la información precontractual completa. & Sí & M-15 Originación y Evaluación Crediticia / M-16 Gestor de Evidencia y Consentimiento / M-18 Administración de Cartera, Cobranza y Repactación / M-19 Migración y Conciliación de Cartera & Subdoc.3 §3.2.5 (SUP-25); §3.2.6.5 Tabla 3.5 (SUP-05, SUP-06); §3.3.5 Tabla 3.11 \textbar{} Subdoc.4 §4.1.5.2, §4.1.7 \\
\addlinespace[2pt]
RF-207 & El sistema deberá registrar la evidencia del consentimiento expreso e informado del cliente mediante un mecanismo verificable de integridad. & Sí & M-15 Originación y Evaluación Crediticia / M-16 Gestor de Evidencia y Consentimiento / M-18 Administración de Cartera, Cobranza y Repactación / M-19 Migración y Conciliación de Cartera & Subdoc.3 §3.2.5 (SUP-25); §3.2.6.5 Tabla 3.5 (SUP-05, SUP-06); §3.3.5 Tabla 3.11 \textbar{} Subdoc.4 §4.1.5.2, §4.1.7 \\
\addlinespace[2pt]
RF-208 & El sistema deberá impedir el registro de una modificacion de condiciones del crédito que no tenga evidencia de consentimiento asociada. & Sí & M-15 Originación y Evaluación Crediticia / M-16 Gestor de Evidencia y Consentimiento / M-18 Administración de Cartera, Cobranza y Repactación / M-19 Migración y Conciliación de Cartera & Subdoc.3 §3.2.5 (SUP-25); §3.2.6.5 Tabla 3.5 (SUP-05, SUP-06); §3.3.5 Tabla 3.11 \textbar{} Subdoc.4 §4.1.5.2, §4.1.7 \\
\addlinespace[2pt]
RF-209 & El sistema deberá reconstruir el acto de consentimiento presentando que se informó, en que versión y que aceptó el cliente. & Sí & M-15 Originación y Evaluación Crediticia / M-16 Gestor de Evidencia y Consentimiento / M-18 Administración de Cartera, Cobranza y Repactación / M-19 Migración y Conciliación de Cartera & Subdoc.3 §3.2.5 (SUP-25); §3.2.6.5 Tabla 3.5 (SUP-05, SUP-06); §3.3.5 Tabla 3.11 \textbar{} Subdoc.4 §4.1.5.2, §4.1.7 \\
\addlinespace[2pt]
RF-210 & El sistema deberá restaurar desde archivo frío los antecedentes de una operación de crédito de cualquier cohorte dentro del plazo de conservación. & Sí & M-15 Originación y Evaluación Crediticia / M-16 Gestor de Evidencia y Consentimiento / M-18 Administración de Cartera, Cobranza y Repactación / M-19 Migración y Conciliación de Cartera & Subdoc.3 §3.2.5 (SUP-25); §3.2.6.5 Tabla 3.5 (SUP-05, SUP-06); §3.3.5 Tabla 3.11 \textbar{} Subdoc.4 §4.1.5.2, §4.1.7 \\
\addlinespace[2pt]
RF-211 & El sistema debe permitir al ejecutivo de cobranza iniciar una gestión de cobranza. & Sí & M-15 Originación y Evaluación Crediticia / M-16 Gestor de Evidencia y Consentimiento / M-18 Administración de Cartera, Cobranza y Repactación / M-19 Migración y Conciliación de Cartera & Subdoc.3 §3.2.5 (SUP-25); §3.2.6.5 Tabla 3.5 (SUP-05, SUP-06); §3.3.5 Tabla 3.11 \textbar{} Subdoc.4 §4.1.5.2, §4.1.7 \\
\addlinespace[2pt]
RF-212 & El sistema debe registrar la gestión de cobranza ejecutada con su medio, su instante y su ejecutor. & Sí & M-15 Originación y Evaluación Crediticia / M-16 Gestor de Evidencia y Consentimiento / M-18 Administración de Cartera, Cobranza y Repactación / M-19 Migración y Conciliación de Cartera & Subdoc.3 §3.2.5 (SUP-25); §3.2.6.5 Tabla 3.5 (SUP-05, SUP-06); §3.3.5 Tabla 3.11 \textbar{} Subdoc.4 §4.1.5.2, §4.1.7 \\
\addlinespace[2pt]
RF-213 & El sistema debe impedir la ejecución de una gestión de cobranza fuera de los límites normativos de horario y de medio de contacto. & Sí & M-15 Originación y Evaluación Crediticia / M-16 Gestor de Evidencia y Consentimiento / M-18 Administración de Cartera, Cobranza y Repactación / M-19 Migración y Conciliación de Cartera & Subdoc.3 §3.2.5 (SUP-25); §3.2.6.5 Tabla 3.5 (SUP-05, SUP-06); §3.3.5 Tabla 3.11 \textbar{} Subdoc.4 §4.1.5.2, §4.1.7 \\
\addlinespace[2pt]
RF-214 & El sistema debe enlazar el expediente de la repactación con la gestión de cobranza que la originó. & Sí & M-15 Originación y Evaluación Crediticia / M-16 Gestor de Evidencia y Consentimiento / M-18 Administración de Cartera, Cobranza y Repactación / M-19 Migración y Conciliación de Cartera & Subdoc.3 §3.2.5 (SUP-25); §3.2.6.5 Tabla 3.5 (SUP-05, SUP-06); §3.3.5 Tabla 3.11 \textbar{} Subdoc.4 §4.1.5.2, §4.1.7 \\
\addlinespace[2pt]
RF-215 & El sistema debe enlazar el expediente de la repactación con la evidencia de consentimiento registrada en RF-207. & Sí & M-15 Originación y Evaluación Crediticia / M-16 Gestor de Evidencia y Consentimiento / M-18 Administración de Cartera, Cobranza y Repactación / M-19 Migración y Conciliación de Cartera & Subdoc.3 §3.2.5 (SUP-25); §3.2.6.5 Tabla 3.5 (SUP-05, SUP-06); §3.3.5 Tabla 3.11 \textbar{} Subdoc.4 §4.1.5.2, §4.1.7 \\
\addlinespace[2pt]
RF-216 & El sistema deberá generar un reporte de conciliación diaria de saldos durante todo el proceso de migración. & Sí & M-15 Originación y Evaluación Crediticia / M-16 Gestor de Evidencia y Consentimiento / M-18 Administración de Cartera, Cobranza y Repactación / M-19 Migración y Conciliación de Cartera & Subdoc.3 §3.2.5 (SUP-25); §3.2.6.5 Tabla 3.5 (SUP-05, SUP-06); §3.3.5 Tabla 3.11 \textbar{} Subdoc.4 §4.1.5.2, §4.1.7 \\
\addlinespace[2pt]
RF-217 & El sistema deberá permitir al cliente no autenticado consultar la información precontractual del crédito. & Sí & M-15 Originación y Evaluación Crediticia / M-16 Gestor de Evidencia y Consentimiento / M-18 Administración de Cartera, Cobranza y Repactación / M-19 Migración y Conciliación de Cartera & Subdoc.3 §3.2.5 (SUP-25); §3.2.6.5 Tabla 3.5 (SUP-05, SUP-06); §3.3.5 Tabla 3.11 \textbar{} Subdoc.4 §4.1.5.2, §4.1.7 \\
\addlinespace[2pt]
RF-218 & El sistema deberá permitir al cliente no autenticado utilizar un simulador de costo total del crédito. & Sí & M-15 Originación y Evaluación Crediticia / M-16 Gestor de Evidencia y Consentimiento / M-18 Administración de Cartera, Cobranza y Repactación / M-19 Migración y Conciliación de Cartera & Subdoc.3 §3.2.5 (SUP-25); §3.2.6.5 Tabla 3.5 (SUP-05, SUP-06); §3.3.5 Tabla 3.11 \textbar{} Subdoc.4 §4.1.5.2, §4.1.7 \\
\addlinespace[2pt]
RF-219 & El sistema deberá mantener cargada la Tasa Máxima Convencional vigente por tipo y tramo de operación, con su fecha de vigencia. & Sí & M-15 Originación y Evaluación Crediticia / M-16 Gestor de Evidencia y Consentimiento / M-18 Administración de Cartera, Cobranza y Repactación / M-19 Migración y Conciliación de Cartera & Subdoc.3 §3.2.5 (SUP-25); §3.2.6.5 Tabla 3.5 (SUP-05, SUP-06); §3.3.5 Tabla 3.11 \textbar{} Subdoc.4 §4.1.5.2, §4.1.7 \\
\addlinespace[2pt]
RF-220 & El sistema deberá impedir la originación de una operación cuya tasa supere la Tasa Máxima Convencional vigente a la fecha de la operación. & Sí & M-15 Originación y Evaluación Crediticia / M-16 Gestor de Evidencia y Consentimiento / M-18 Administración de Cartera, Cobranza y Repactación / M-19 Migración y Conciliación de Cartera & Subdoc.3 §3.2.5 (SUP-25); §3.2.6.5 Tabla 3.5 (SUP-05, SUP-06); §3.3.5 Tabla 3.11 \textbar{} Subdoc.4 §4.1.5.2, §4.1.7 \\
\addlinespace[2pt]
RF-221 & El sistema deberá determinar el cupo exclusivamente a partir de las variables de la evaluación de capacidad de pago documentada. & Sí & M-15 Originación y Evaluación Crediticia / M-16 Gestor de Evidencia y Consentimiento / M-18 Administración de Cartera, Cobranza y Repactación / M-19 Migración y Conciliación de Cartera & Subdoc.3 §3.2.5 (SUP-25); §3.2.6.5 Tabla 3.5 (SUP-05, SUP-06); §3.3.5 Tabla 3.11 \textbar{} Subdoc.4 §4.1.5.2, §4.1.7 \\
\addlinespace[2pt]
RF-222 & El sistema deberá excluir el monto de la venta en curso del conjunto de variables de asignación de cupo. & Sí & M-15 Originación y Evaluación Crediticia / M-16 Gestor de Evidencia y Consentimiento / M-18 Administración de Cartera, Cobranza y Repactación / M-19 Migración y Conciliación de Cartera & Subdoc.3 §3.2.5 (SUP-25); §3.2.6.5 Tabla 3.5 (SUP-05, SUP-06); §3.3.5 Tabla 3.11 \textbar{} Subdoc.4 §4.1.5.2, §4.1.7 \\
\addlinespace[2pt]
RF-223 & El sistema deberá excluir la solicitud del vendedor del conjunto de variables de asignación de cupo. & Sí & M-15 Originación y Evaluación Crediticia / M-16 Gestor de Evidencia y Consentimiento / M-18 Administración de Cartera, Cobranza y Repactación / M-19 Migración y Conciliación de Cartera & Subdoc.3 §3.2.5 (SUP-25); §3.2.6.5 Tabla 3.5 (SUP-05, SUP-06); §3.3.5 Tabla 3.11 \textbar{} Subdoc.4 §4.1.5.2, §4.1.7 \\
\addlinespace[2pt]
RF-224 & El sistema deberá impedir al vendedor alterar el resultado de una evaluación crediticia. & Sí & M-15 Originación y Evaluación Crediticia / M-16 Gestor de Evidencia y Consentimiento / M-18 Administración de Cartera, Cobranza y Repactación / M-19 Migración y Conciliación de Cartera & Subdoc.3 §3.2.5 (SUP-25); §3.2.6.5 Tabla 3.5 (SUP-05, SUP-06); §3.3.5 Tabla 3.11 \textbar{} Subdoc.4 §4.1.5.2, §4.1.7 \\
\addlinespace[2pt]
RF-225 & El sistema deberá impedir la repetición de la evaluación crediticia del mismo cliente dentro de la ventana de enfriamiento parametrizada. & Sí & M-15 Originación y Evaluación Crediticia / M-16 Gestor de Evidencia y Consentimiento / M-18 Administración de Cartera, Cobranza y Repactación / M-19 Migración y Conciliación de Cartera & Subdoc.3 §3.2.5 (SUP-25); §3.2.6.5 Tabla 3.5 (SUP-05, SUP-06); §3.3.5 Tabla 3.11 \textbar{} Subdoc.4 §4.1.5.2, §4.1.7 \\
\addlinespace[2pt]
RF-226 & El sistema deberá registrar cada intento de evaluación con la identidad del solicitante y el resultado obtenido. & Sí & M-15 Originación y Evaluación Crediticia / M-16 Gestor de Evidencia y Consentimiento / M-18 Administración de Cartera, Cobranza y Repactación / M-19 Migración y Conciliación de Cartera & Subdoc.3 §3.2.5 (SUP-25); §3.2.6.5 Tabla 3.5 (SUP-05, SUP-06); §3.3.5 Tabla 3.11 \textbar{} Subdoc.4 §4.1.5.2, §4.1.7 \\
\addlinespace[2pt]
RNF-04 & El estado de un pedido no debe presentar discrepancias entre canales de consulta. & Sí & M-11 Portal / M-09 OMS — estado único del pedido consultado por todos los canales & Subdoc.3 §3.3.3 Tabla 3.9 (M-09, M-11) \\
\addlinespace[2pt]
RNF-05 & La actualización del estado de un pedido debe reflejarse en todos los puntos de consulta dentro del plazo máximo definido. & Sí & Arquitectura orientada a eventos con presupuesto de latencia por componente (M-01, M-05, M-08, M-09, M-11) & Subdoc.4 §4.1.1.1 Tabla 4.11 (Latencia en operaciones) \\
\addlinespace[2pt]
RNF-06 & El tiempo de evaluación de una solicitud de crédito en el punto de venta físico no debe exceder el umbral objetivo. & Sí & Arquitectura orientada a eventos con presupuesto de latencia por componente (M-01, M-05, M-08, M-09, M-11) & Subdoc.4 §4.1.1.1 Tabla 4.11 (Latencia en operaciones) \\
\addlinespace[2pt]
RNF-07 & El tiempo total de originación de crédito en el mesón, con información precontractual acreditada, no debe superar el tiempo del proceso actual. & Sí & M-15 Originación Crediticia — evaluación ≤8s embebida en flujo POS & Subdoc.3 §3.5.2 OE-09 \textbar{} Subdoc.4 §4.1.1.1 Tabla 4.11 \\
\addlinespace[2pt]
RNF-08 & La migración de la cartera de crédito viva no debe generar pérdida de datos. & Sí & M-19 Migración y Conciliación de Cartera (olas de coexistencia, circuit breaker) & Subdoc.3 §3.2.5 (SUP-25) \textbar{} Subdoc.4 §4.3 ADR-06 \\
\addlinespace[2pt]
RNF-09 & La migración de la cartera de crédito viva no debe interrumpir el servicio de cobro. & Sí & M-19 Migración y Conciliación de Cartera (olas de coexistencia, circuit breaker) & Subdoc.3 §3.2.5 (SUP-25) \textbar{} Subdoc.4 §4.3 ADR-06 \\
\addlinespace[2pt]
RNF-10 & La migración de la cartera de crédito viva no debe producir divergencia de saldos. & Sí & M-19 Migración y Conciliación de Cartera (olas de coexistencia, circuit breaker) & Subdoc.3 §3.2.5 (SUP-25) \textbar{} Subdoc.4 §4.3 ADR-06 \\
\addlinespace[2pt]
RNF-11 & La separación lógica entre los datos del retail y los del negocio financiero debe estar implementada. & Sí & Límites de Contexto (DDD) + separación física Z3 (ADR-11, ADR-12) & Subdoc.4 §4.1.2 (Límites de contexto), §4.2.4.2 (Tabla 4.32, ADR-12) \\
\addlinespace[2pt]
RNF-12 & La separación lógica entre ambos ámbitos debe estar documentada. & Sí & Límites de Contexto (DDD) + separación física Z3 (ADR-11, ADR-12) & Subdoc.4 §4.1.2 (Límites de contexto), §4.2.4.2 (Tabla 4.32, ADR-12) \\
\addlinespace[2pt]
RNF-13 & La separación lógica entre ambos ámbitos debe ser verificada por auditoría independiente. & Sí & Límites de Contexto (DDD) + separación física Z3 (ADR-11, ADR-12) & Subdoc.4 §4.1.2 (Límites de contexto), §4.2.4.2 (Tabla 4.32, ADR-12) \\
\addlinespace[2pt]
RNF-14 & La separación física entre la infraestructura del retail y la de la filial emisora debe estar implementada y acreditada. & Sí & Límites de Contexto (DDD) + separación física Z3 (ADR-11, ADR-12) & Subdoc.4 §4.1.2 (Límites de contexto), §4.2.4.2 (Tabla 4.32, ADR-12) \\
\addlinespace[2pt]
RNF-15 & Toda acción relevante ejecutada en el sistema debe quedar registrada indicando usuario individual, acción e instante. & Sí & M-21 Gobierno de Frontera de Datos / M-23 IAM — registro inalterable de acciones & Subdoc.3 §3.2.3; §3.3.6 Tabla 3.12 \textbar{} Subdoc.4 §4.1.7 (frente 5) \\
\addlinespace[2pt]
RNF-16 & La consulta de disponibilidad en la ficha de producto del canal digital debe responder dentro del umbral definido. & Sí & Arquitectura orientada a eventos con presupuesto de latencia por componente (M-01, M-05, M-08, M-09, M-11) & Subdoc.4 §4.1.1.1 Tabla 4.11 (Latencia en operaciones) \\
\addlinespace[2pt]
RNF-17 & La confirmación de un pedido durante el evento anual debe completarse dentro del umbral definido. & Sí & Arquitectura orientada a eventos con presupuesto de latencia por componente (M-01, M-05, M-08, M-09, M-11) & Subdoc.4 §4.1.1.1 Tabla 4.11 (Latencia en operaciones) \\
\addlinespace[2pt]
RNF-18 & La venta completa en caja con medio de pago externo debe completarse dentro del umbral definido. & Sí & Arquitectura orientada a eventos con presupuesto de latencia por componente (M-01, M-05, M-08, M-09, M-11) & Subdoc.4 §4.1.1.1 Tabla 4.11 (Latencia en operaciones) \\
\addlinespace[2pt]
RNF-19 & La propagación de un cambio de precio a las 380 líneas de caja y al canal digital debe completarse dentro del umbral definido. & Sí & Arquitectura orientada a eventos con presupuesto de latencia por componente (M-01, M-05, M-08, M-09, M-11) & Subdoc.4 §4.1.1.1 Tabla 4.11 (Latencia en operaciones) \\
\addlinespace[2pt]
RNF-20 & El registro de una devolución en el mesón debe completarse dentro del umbral definido. & Sí & Arquitectura orientada a eventos con presupuesto de latencia por componente (M-01, M-05, M-08, M-09, M-11) & Subdoc.4 §4.1.1.1 Tabla 4.11 (Latencia en operaciones) \\
\addlinespace[2pt]
RNF-21 & La consulta de disponibilidad desde una terminal compartida del piso de venta debe responder dentro del umbral definido. & Sí & Arquitectura orientada a eventos con presupuesto de latencia por componente (M-01, M-05, M-08, M-09, M-11) & Subdoc.4 §4.1.1.1 Tabla 4.11 (Latencia en operaciones) \\
\addlinespace[2pt]
RNF-22 & La plataforma debe soportar la concurrencia y el volumen del peak digital del evento anual sin degradar los umbrales de NFR-06.x. & Sí & Escalamiento elástico Z1 (M-01, M-09, M-11) y capacidad fija Z2/Z3 & Subdoc.4 §4.0.4.1 Tabla 4.3, §4.2.9 Dimensionamiento \\
\addlinespace[2pt]
RNF-23 & La plataforma debe soportar la concurrencia y el volumen del peak presencial de la campaña de noviembre y diciembre sin degradar los umbrales de NFR-06.x. & Sí & Escalamiento elástico Z1 (M-01, M-09, M-11) y capacidad fija Z2/Z3 & Subdoc.4 §4.0.4.1 Tabla 4.3, §4.2.9 Dimensionamiento \\
\addlinespace[2pt]
RNF-24 & Una tienda debe poder operar sin enlace externo durante el período mínimo definido. & Sí & Nodo de tienda on-premise (Z2) + Nube Multi-AZ (Z1) + Site Principal/Secundario (Z3) & Subdoc.4 §4.2.3 (autonomía offline), §4.2.5 (Site Secundario, RTO/RPO), §4.2.7.2 (niveles de servicio) \\
\addlinespace[2pt]
RNF-25 & El centro de distribución principal debe poder operar sin enlace externo durante el período mínimo definido. & Sí & Nodo de tienda on-premise (Z2) + Nube Multi-AZ (Z1) + Site Principal/Secundario (Z3) & Subdoc.4 §4.2.3 (autonomía offline), §4.2.5 (Site Secundario, RTO/RPO), §4.2.7.2 (niveles de servicio) \\
\addlinespace[2pt]
RNF-26 & La sincronización tras la reconexión no debe superar el plazo definido. & Sí & Nodo de tienda on-premise (Z2) + Nube Multi-AZ (Z1) + Site Principal/Secundario (Z3) & Subdoc.4 §4.2.3 (autonomía offline), §4.2.5 (Site Secundario, RTO/RPO), §4.2.7.2 (niveles de servicio) \\
\addlinespace[2pt]
RNF-27 & La sincronización tras la reconexión no debe producir pérdida de ventas ni de documentos. & Sí & Nodo de tienda on-premise (Z2) + Nube Multi-AZ (Z1) + Site Principal/Secundario (Z3) & Subdoc.4 §4.2.3 (autonomía offline), §4.2.5 (Site Secundario, RTO/RPO), §4.2.7.2 (niveles de servicio) \\
\addlinespace[2pt]
RNF-28 & La operación desconectada autónoma on-premise debe sostenerse durante el período mínimo exigido por las bases técnicas. & Sí & Nodo de tienda on-premise (Z2) + Nube Multi-AZ (Z1) + Site Principal/Secundario (Z3) & Subdoc.4 §4.2.3 (autonomía offline), §4.2.5 (Site Secundario, RTO/RPO), §4.2.7.2 (niveles de servicio) \\
\addlinespace[2pt]
RNF-29 & La operación de tiendas y centros de distribución debe estar disponible en la ventana horaria definida. & Sí & Nodo de tienda on-premise (Z2) + Nube Multi-AZ (Z1) + Site Principal/Secundario (Z3) & Subdoc.4 §4.2.3 (autonomía offline), §4.2.5 (Site Secundario, RTO/RPO), §4.2.7.2 (niveles de servicio) \\
\addlinespace[2pt]
RNF-30 & El canal digital debe estar disponible de forma continua. & Sí & Nodo de tienda on-premise (Z2) + Nube Multi-AZ (Z1) + Site Principal/Secundario (Z3) & Subdoc.4 §4.2.3 (autonomía offline), §4.2.5 (Site Secundario, RTO/RPO), §4.2.7.2 (niveles de servicio) \\
\addlinespace[2pt]
RNF-31 & Los servicios financieros que afectan pagos, estados de cuenta y bloqueo de tarjetas deben estar disponibles de forma continua. & Sí & Nodo de tienda on-premise (Z2) + Nube Multi-AZ (Z1) + Site Principal/Secundario (Z3) & Subdoc.4 §4.2.3 (autonomía offline), §4.2.5 (Site Secundario, RTO/RPO), §4.2.7.2 (niveles de servicio) \\
\addlinespace[2pt]
RNF-32 & La solución debe cumplir los objetivos de recuperación ante desastre exigidos. & Sí & Site Principal / Site Secundario (RTO≤4h, RPO≤15min) — M-19 restauración desde archivo frío & Subdoc.4 §4.2.5.4 Tabla 4.37 (objetivos de recuperación) \\
\addlinespace[2pt]
RNF-33 & Los datos de identificación de clientes, la cartera de crédito, el comportamiento de pago y los antecedentes de evaluación crediticia deben cifrarse a nivel de campo. & Sí & Arquitectura Zero Trust (NIST SP 800-207): M-21/M-23, cifrado de campo, tokenización, HSM & Subdoc.4 §4.1.7 (arquitectura de seguridad, frentes 1-3) \\
\addlinespace[2pt]
RNF-34 & Los datos de medios de pago no deben almacenarse en claro, sino tokenizarse. & Sí & Arquitectura Zero Trust (NIST SP 800-207): M-21/M-23, cifrado de campo, tokenización, HSM & Subdoc.4 §4.1.7 (arquitectura de seguridad, frentes 1-3) \\
\addlinespace[2pt]
RNF-35 & La red de cajas, la administrativa, la de videovigilancia y la inalámbrica de clientes deben estar segmentadas en las 22 tiendas. & Sí & Arquitectura Zero Trust (NIST SP 800-207): M-21/M-23, cifrado de campo, tokenización, HSM & Subdoc.4 §4.1.7 (arquitectura de seguridad, frentes 1-3) \\
\addlinespace[2pt]
RNF-36 & La red y el ámbito de sistemas de la filial emisora deben estar separados de forma acreditada respecto del retail. & Sí & Arquitectura Zero Trust (NIST SP 800-207): M-21/M-23, cifrado de campo, tokenización, HSM & Subdoc.4 §4.1.7 (arquitectura de seguridad, frentes 1-3) \\
\addlinespace[2pt]
RNF-37 & Debe existir segregación de funciones verificable entre originación, aprobación, modificacion de condiciones y cobranza. & Sí & Arquitectura Zero Trust (NIST SP 800-207): M-21/M-23, cifrado de campo, tokenización, HSM & Subdoc.4 §4.1.7 (arquitectura de seguridad, frentes 1-3) \\
\addlinespace[2pt]
RNF-38 & Debe registrarse todo acceso a la cartera de crédito, al comportamiento de pago y a los antecedentes de evaluación crediticia. & Sí & M-21 Gobierno de Frontera de Datos / M-23 IAM — registro inalterable de acciones & Subdoc.3 §3.2.3; §3.3.6 Tabla 3.12 \textbar{} Subdoc.4 §4.1.7 (frente 5) \\
\addlinespace[2pt]
RNF-39 & Los actos de crédito de apertura de tarjeta y de repactación deben contar con un mecanismo de consentimiento de integridad verificable. & Sí & M-16 Gestor de Evidencia (WORM) — firma electrónica avanzada, Ley N°19.799 & Subdoc.4 §4.1.7 (evidencia y consentimiento) \\
\addlinespace[2pt]
RNF-40 & La conformidad de recepción de mercadería y la constancia de devolución deben contar con un mecanismo de consentimiento de integridad verificable. & Sí & M-16 Gestor de Evidencia (WORM) — firma electrónica avanzada, Ley N°19.799 & Subdoc.4 §4.1.7 (evidencia y consentimiento) \\
\addlinespace[2pt]
RNF-41 & Los documentos tributarios y antecedentes de venta deben conservarse por el plazo definido. & Sí & M-16 Gestor de Evidencia (WORM) / M-13 Postventa (garantía legal) / M-21 Gobierno de Frontera de Datos & Subdoc.4 §4.1.7 (frentes 3 y 6); §4.1.6.4 (retención y contratos) \\
\addlinespace[2pt]
RNF-42 & Los antecedentes del crédito deben conservarse por el plazo del crédito más el período posterior exigido. & Sí & M-16 Gestor de Evidencia (WORM) / M-13 Postventa (garantía legal) / M-21 Gobierno de Frontera de Datos & Subdoc.4 §4.1.7 (frentes 3 y 6); §4.1.6.4 (retención y contratos) \\
\addlinespace[2pt]
RNF-43 & La evidencia de consentimiento de originación y repactación debe conservarse por el mismo plazo que el crédito. & Sí & M-16 Gestor de Evidencia (WORM) / M-13 Postventa (garantía legal) / M-21 Gobierno de Frontera de Datos & Subdoc.4 §4.1.7 (frentes 3 y 6); §4.1.6.4 (retención y contratos) \\
\addlinespace[2pt]
RNF-44 & La trazabilidad del precio publicado debe conservarse por el plazo definido. & Sí & M-16 Gestor de Evidencia (WORM) / M-13 Postventa (garantía legal) / M-21 Gobierno de Frontera de Datos & Subdoc.4 §4.1.7 (frentes 3 y 6); §4.1.6.4 (retención y contratos) \\
\addlinespace[2pt]
RNF-45 & Los movimientos de inventario y conteos cíclicos deben conservarse por el plazo definido. & Sí & M-16 Gestor de Evidencia (WORM) / M-13 Postventa (garantía legal) / M-21 Gobierno de Frontera de Datos & Subdoc.4 §4.1.7 (frentes 3 y 6); §4.1.6.4 (retención y contratos) \\
\addlinespace[2pt]
RNF-46 & Los datos de fidelización deben conservarse por el plazo definido. & Sí & M-16 Gestor de Evidencia (WORM) / M-13 Postventa (garantía legal) / M-21 Gobierno de Frontera de Datos & Subdoc.4 §4.1.7 (frentes 3 y 6); §4.1.6.4 (retención y contratos) \\
\addlinespace[2pt]
RNF-47 & Los registros de videovigilancia deben conservarse por el plazo definido. & Sí & M-16 Gestor de Evidencia (WORM) / M-13 Postventa (garantía legal) / M-21 Gobierno de Frontera de Datos & Subdoc.4 §4.1.7 (frentes 3 y 6); §4.1.6.4 (retención y contratos) \\
\addlinespace[2pt]
RNF-48 & La solución debe admitir la apertura de una tienda nueva por parametrización, sin desarrollo a medida. & Sí & Nodo de tienda como unidad estandarizada replicable (RT-02.12-C) & Subdoc.4 §4.1.8.2 (Replicación a una unidad nueva) \\
\addlinespace[2pt]
RNF-49 & Las interfaces de sala, caja y mesón financiero deben ser operables por personal recién incorporado con capacitación mínima. & Sí & M-08 POS (PWA) — interfaces operables con capacitación mínima & Subdoc.4 §4.2.10.3 Tabla 4.53 (PWA sin instalación por dispositivo) \\
\addlinespace[2pt]
RNF-50 & Toda comunicación relativa al crédito debe cumplir criterios verificables de lenguaje claro, validados con personas usuarias reales. & Sí & M-15/M-16 — simulador de costo total en lenguaje normalizado & Subdoc.3 §3.3.5 Tabla 3.11 (M-15) \\
\addlinespace[2pt]
RNF-51 & Las notificaciones definidas funcionalmente deben emitirse y registrarse en el 100\% de los eventos que las gatillan. & Sí & M-11 Portal (notificaciones transaccionales) / M-14 Conciliación de Recobros B2B & Subdoc.3 §3.3.3/3.3.4 \textbar{} Subdoc.4 §4.1.5.4 (catálogo de eventos de negocio) \\
\addlinespace[2pt]
RNF-52 & La solución debe adoptar un estándar identificado y justificado para el catálogo e identificación de producto. & Sí & M-20 Bus de Eventos y API Gateway — catálogo de contratos versionados (OpenAPI/AsyncAPI) & Subdoc.4 §4.1.6.4 (gobernanza de contratos), §4.1.4 (catálogo de interfaces) \\
\addlinespace[2pt]
RNF-53 & La solución debe adoptar un estándar identificado y justificado para el intercambio de órdenes y avisos de despacho con los proveedores. & Sí & M-20 Bus de Eventos y API Gateway — catálogo de contratos versionados (OpenAPI/AsyncAPI) & Subdoc.4 §4.1.6.4 (gobernanza de contratos), §4.1.4 (catálogo de interfaces) \\
\addlinespace[2pt]
RNF-54 & La solución debe adoptar un estándar identificado y justificado para el intercambio con los transportistas de última milla. & Sí & M-20 Bus de Eventos y API Gateway — catálogo de contratos versionados (OpenAPI/AsyncAPI) & Subdoc.4 §4.1.6.4 (gobernanza de contratos), §4.1.4 (catálogo de interfaces) \\
\addlinespace[2pt]
RNF-55 & La solución debe adoptar un estándar identificado y justificado para la sincronización de catálogo, existencia y pedido con los vendedores de marketplace. & Sí & M-20 Bus de Eventos y API Gateway — catálogo de contratos versionados (OpenAPI/AsyncAPI) & Subdoc.4 §4.1.6.4 (gobernanza de contratos), §4.1.4 (catálogo de interfaces) \\
\addlinespace[2pt]
RNF-56 & La solución debe adoptar el formato normativo vigente de documento tributario electrónico. & Sí & M-20 Bus de Eventos y API Gateway — catálogo de contratos versionados (OpenAPI/AsyncAPI) & Subdoc.4 §4.1.6.4 (gobernanza de contratos), §4.1.4 (catálogo de interfaces) \\
\addlinespace[2pt]
RNF-57 & La evidencia completa de consentimiento de una operación arbitraria debe recuperarse dentro del plazo objetivo, a diez años. & Sí & M-16 Gestor de Evidencia (WORM) — recuperación de evidencia de consentimiento & Subdoc.3 §3.2.6.5 Tabla 3.5 (SUP-05) \textbar{} Subdoc.4 §4.2.5.4 \\
\addlinespace[2pt]
RNF-58 & La recuperación del precio publicado de una referencia en una fecha, hora y canal arbitrarios debe completarse dentro del plazo objetivo. & Sí & Arquitectura orientada a eventos con presupuesto de latencia por componente (M-01, M-05, M-08, M-09, M-11) & Subdoc.4 §4.1.1.1 Tabla 4.11 (Latencia en operaciones) \\
\addlinespace[2pt]
RNF-59 & La revocación total de accesos de un trabajador debe completarse dentro del plazo máximo de la política interna. & Sí & Arquitectura Zero Trust (NIST SP 800-207): M-21/M-23, cifrado de campo, tokenización, HSM & Subdoc.4 §4.1.7 (arquitectura de seguridad, frentes 1-3) \\
\addlinespace[2pt]
RNF-60 & El traspaso de sesión nominativa en una terminal compartida debe completarse dentro del plazo objetivo. & Sí & M-23 IAM — traspaso rápido de sesión (hot-swapping) & Subdoc.3 §3.2.6.6 Tabla 3.6 (SUP-24) \\
\addlinespace[2pt]
RNF-61 & La tasa de cancelación de pedidos por falta de existencia no debe superar el umbral comprometido. & Sí & M-01 ATP / M-09 OMS — indicadores comprometidos como Objetivos Específicos & Subdoc.3 §3.5.2 Tabla 3.13 (OE-01 a OE-12), §3.11.3 Tabla 3.31 \\
\addlinespace[2pt]
RNF-62 & El cumplimiento de la promesa de entrega publicada debe alcanzar el umbral comprometido, medido por pedido. & Sí & M-01 ATP / M-09 OMS — indicadores comprometidos como Objetivos Específicos & Subdoc.3 §3.5.2 Tabla 3.13 (OE-01 a OE-12), §3.11.3 Tabla 3.31 \\
\addlinespace[2pt]
RNF-63 & La diferencia de inventario por categoría debe reducirse al umbral comprometido. & Sí & M-01 ATP / M-09 OMS — indicadores comprometidos como Objetivos Específicos & Subdoc.3 §3.5.2 Tabla 3.13 (OE-01 a OE-12), §3.11.3 Tabla 3.31 \\
\addlinespace[2pt]
RNF-64 & El tiempo entre la detección de un quiebre y el aviso al cliente no debe superar el objetivo declarado. & Sí & M-11 Portal (notificaciones transaccionales) / M-14 Conciliación de Recobros B2B & Subdoc.3 §3.3.3/3.3.4 \textbar{} Subdoc.4 §4.1.5.4 (catálogo de eventos de negocio) \\
\addlinespace[2pt]
RNF-65 & No debe mantenerse un cobro capturado sobre un pedido declarado no cumplible. & Sí & M-11 Portal (notificaciones transaccionales) / M-14 Conciliación de Recobros B2B & Subdoc.3 §3.3.3/3.3.4 \textbar{} Subdoc.4 §4.1.5.4 (catálogo de eventos de negocio) \\
\addlinespace[2pt]
RNF-66 & La latencia entre la recepción física de una devolución en tienda y la notificación al vendedor de marketplace no debe superar el objetivo declarado. & Sí & M-11 Portal (notificaciones transaccionales) / M-14 Conciliación de Recobros B2B & Subdoc.3 §3.3.3/3.3.4 \textbar{} Subdoc.4 §4.1.5.4 (catálogo de eventos de negocio) \\
\addlinespace[2pt]
RNF-67 & No debe permanecer mercadería de marketplace en bodega sin aviso a su propietario. & Sí & M-11 Portal (notificaciones transaccionales) / M-14 Conciliación de Recobros B2B & Subdoc.3 §3.3.3/3.3.4 \textbar{} Subdoc.4 §4.1.5.4 (catálogo de eventos de negocio) \\
\addlinespace[2pt]
RNF-68 & La restauración desde archivo frío de una operación de crédito de la cohorte más antigua debe completarse dentro del plazo objetivo. & Sí & Site Principal / Site Secundario (RTO≤4h, RPO≤15min) — M-19 restauración desde archivo frío & Subdoc.4 §4.2.5.4 Tabla 4.37 (objetivos de recuperación) \\
\addlinespace[2pt]
RNF-69 & La atención de garantía legal en mesón no debe derivar al consumidor a un tercero como condición de atención, en ninguna tienda. & Sí & M-16 Gestor de Evidencia (WORM) / M-13 Postventa (garantía legal) / M-21 Gobierno de Frontera de Datos & Subdoc.4 §4.1.7 (frentes 3 y 6); §4.1.6.4 (retención y contratos) \\
\addlinespace[2pt]
RNF-70 & El ciclo de desarrollo debe producir un inventario de componentes (SBOM) por artefacto desplegado. & Sí & Cadena DevSecOps (SBOM, SLSA, OWASP ASVS, Zero Trust) & Subdoc.4 §4.1.7 (quinto frente, seguridad en el desarrollo), Tabla 4.53 (stack DevSecOps) \\
\addlinespace[2pt]
RNF-71 & La cadena de suministro de software debe alcanzar el nivel SLSA 3 o superior. & Sí & Cadena DevSecOps (SBOM, SLSA, OWASP ASVS, Zero Trust) & Subdoc.4 §4.1.7 (quinto frente, seguridad en el desarrollo), Tabla 4.53 (stack DevSecOps) \\
\addlinespace[2pt]
RNF-72 & La capa expuesta debe cumplir el estándar OWASP ASVS en el nivel comprometido. & Sí & Cadena DevSecOps (SBOM, SLSA, OWASP ASVS, Zero Trust) & Subdoc.4 §4.1.7 (quinto frente, seguridad en el desarrollo), Tabla 4.53 (stack DevSecOps) \\
\addlinespace[2pt]
RNF-73 & La arquitectura debe implementar un modelo Zero Trust conforme a NIST SP 800-207. & Sí & Cadena DevSecOps (SBOM, SLSA, OWASP ASVS, Zero Trust) & Subdoc.4 §4.1.7 (quinto frente, seguridad en el desarrollo), Tabla 4.53 (stack DevSecOps) \\
\addlinespace[2pt]
RNF-74 & Todo tratamiento de datos personales debe declarar su finalidad. & Sí & M-16 Gestor de Evidencia (WORM) / M-13 Postventa (garantía legal) / M-21 Gobierno de Frontera de Datos & Subdoc.4 §4.1.7 (frentes 3 y 6); §4.1.6.4 (retención y contratos) \\
\addlinespace[2pt]
RNF-75 & Todo tratamiento de datos personales debe invocar una de las bases de licitud de la Ley N° 21.719. & Sí & M-16 Gestor de Evidencia (WORM) / M-13 Postventa (garantía legal) / M-21 Gobierno de Frontera de Datos & Subdoc.4 §4.1.7 (frentes 3 y 6); §4.1.6.4 (retención y contratos) \\
\addlinespace[2pt]
RNF-76 & La compañía debe mantener el registro de actividades de tratamiento de datos personales. & Sí & M-16 Gestor de Evidencia (WORM) / M-13 Postventa (garantía legal) / M-21 Gobierno de Frontera de Datos & Subdoc.4 §4.1.7 (frentes 3 y 6); §4.1.6.4 (retención y contratos) \\
\addlinespace[2pt]
RT-02.01 & La solución se organizará en las ocho capas del numeral 2.1. El PROPONENTE presentará el diagrama de la arquitectura lógica identificando cada capa, sus componentes y las interfaces entre ellas. & Sí & Estructura por Capas (8 capas), con módulos Ancoa y emplazamiento híbrido asignado a cada una & Subdocumento 4, sección 4.1.3, Tabla 4.14 \\
\addlinespace[2pt]
RT-02.02 & La arquitectura será modular, con límites de contexto explícitos y acoplamiento débil. Se rechazará toda arquitectura monolítica que no permita desplegar de forma independiente sus componentes críticos. & Sí & Arquitectura Modular de Microservicios Sin Estado agrupada en Límites de Contexto (Domain-Driven Design), ADR-01 & Subdocumento 4, sección 4.1.1 y 4.3 (ADR-01); Subdocumento 3, línea 362 \\
\addlinespace[2pt]
RT-02.03 & La descripción de la arquitectura se ajustará a ISO/IEC/IEEE 42010, con vistas lógica, de procesos, de despliegue, de datos y de seguridad. & Sí & Cinco vistas de arquitectura conforme a ISO/IEC/IEEE 42010 (lógica, procesos, despliegue, datos, seguridad) & Subdocumento 4, sección 4.1.8, Tabla 4.26 \\
\addlinespace[2pt]
RT-02.04 & El PROPONENTE mantendrá un registro de decisiones de arquitectura (ADR) fechado, con la alternativa escogida, las alternativas descartadas y el criterio de decisión. El registro es entregable contractual y se actualizará durante toda la ejecución. & Sí & Registro de Decisiones de Arquitectura (ADR-01 a ADR-12), fechado, con alternativa descartada y criterio & Subdocumento 4, sección 4.3, Tabla de ADR \\
\addlinespace[2pt]
RT-02.05 & La capa de servicios de negocio será sin estado. El estado de sesión y el estado de proceso residirán en almacenes externos con alta disponibilidad. & Sí & Servicios de negocio sin estado (Go/Java); estado de sesión externalizado en Azure Cache for Redis & Subdocumento 4, secciones 4.1.3 (Tabla 4.14) y 4.2.10.3 (Tabla 4.53) \\
\addlinespace[2pt]
RT-02.06 & Toda operación de escritura expuesta a reintentos será idempotente, con clave de idempotencia declarada por el cliente y ventana de deduplicación documentada. & Sí & Clave de idempotencia provista por el cliente en el API Gateway; registro de resultado de la primera ejecución & Subdocumento 4, sección 4.1.6.5, Tabla 4.21 \\
\addlinespace[2pt]
RT-02.07 & Los flujos de eventos garantizarán entrega al menos una vez, con deduplicación en el consumidor y orden garantizado dentro de la partición o del agregado cuando el proceso lo exija. & Sí & Bus de eventos (M-20, Azure Event Hubs compatible Kafka) con entrega al menos una vez, deduplicación en consumidor y orden por partición & Subdocumento 4, secciones 4.1.6.2 y 4.2.9.3, Tabla 4.46 \\
\addlinespace[2pt]
RT-02.08 & La solución implementará patrones de resiliencia demostrables: reintento con retroceso exponencial y variación aleatoria, cortacircuitos, mamparos de aislamiento, límites de tasa y tiempo de espera explícito en toda llamada remota. No se admiten llamadas remotas sin tiempo de espera. & Sí & Patrones de resiliencia (límites de tasa, tiempos de espera, reintento con retroceso exponencial, mamparos, cortacircuitos) & Subdocumento 4, sección 4.1.6.6, Tabla 4.22 \\
\addlinespace[2pt]
RT-02.09 & La solución degradará de forma elegante: ante la indisponibilidad de un componente no crítico deberá continuar operando en modo reducido, informando la degradación a la persona usuaria, y nunca fallar de forma total. & Sí & Régimen de degradación elegante técnica y restrictiva comercial, gobernado por M-24 & Subdocumento 4, secciones 4.1.6.6 y 4.1.10 \\
\addlinespace[2pt]
RT-02.10 & Las capas de aplicación e integración escalarán horizontalmente de forma automática, con umbrales, límites superiores y costo asociado declarados en la oferta. & Sí & Escalamiento automático en AKS con umbral de expansión/contracción y límite superior por servicio & Subdocumento 4, sección 4.1.10 y 4.2.9.2, Tabla 4.45 \\
\addlinespace[2pt]
RT-02.11 & El PROPONENTE declarará explícitamente los puntos únicos de falla que subsistan en su arquitectura y justificará por qué son aceptables. Omitir esta declaración cuando existan puntos únicos de falla se evaluará como observación grave. & Sí & Declaración de seis puntos únicos de falla (ERP, enlace Coyhaique, Site Principal, firma electrónica, enlace de 14 tiendas, M-21) con justificación y mitigación & Subdocumento 4, sección 4.1.9, Tabla 4.29 \\
\addlinespace[2pt]
RT-02.12 & La solución admitirá su replicación a nuevas unidades, sitios, sucursales o filiales del CLIENTE sin rediseño arquitectónico, mediante parametrización o multi-tenencia. & Sí & Nodo de tienda como unidad empaquetada y estandarizada, replicable por aprovisionamiento desde imagen & Subdocumento 4, sección 4.1.8.2 \\
\addlinespace[2pt]
RT-02.13 & El PROPONENTE presentará un modelo de dominio del negocio del caso, con las entidades principales, sus relaciones y los eventos de negocio que las modifican. & Sí & Modelo de dominio por Límite de Contexto (retail y financiero), con catálogo de eventos de negocio & Subdocumento 4, sección 4.1.5, Tablas 4.17, 4.18 y 4.19 \\
\addlinespace[2pt]
RT-02.14 & Se valorará la aplicación documentada de patrones de arquitectura evolutiva que permitan sustituir un componente sin reescribir la solución: capa anticorrupción frente a sistemas heredados, estrangulamiento progresivo y abstracción de proveedores. & Sí & Patrón Estrangulador (Strangler Fig) con capa anticorrupción frente a sistemas heredados & Subdocumento 4, sección 4.1.1.3 \\
\addlinespace[2pt]
RT-03.01 & El PROPONENTE declarará el proveedor de nube pública, la región primaria y la región secundaria utilizadas. El proveedor deberá contar con presencia de región o zona en Chile o en Sudamérica. & No & - & - \\
\addlinespace[2pt]
RT-03.02 & Todos los componentes con requisito de alta disponibilidad se desplegarán en al menos dos zonas de disponibilidad. No se aceptará un diseño en una sola zona. & Sí & Despliegue Multi-AZ en Z1 (nube pública) para borde, servicios de negocio retail, bus de eventos y bases de datos & Subdocumento 4, secciones 4.1.3, 4.2.9.4 y 4.2.7.2 \\
\addlinespace[2pt]
RT-03.03 & La totalidad de la infraestructura se definirá como código, versionada en el repositorio del CLIENTE, revisable y reproducible. No se admite infraestructura creada manualmente por consola, salvo la cuenta raíz inicial, cuya creación deberá documentarse. & Sí & Infraestructura como código con OpenTofu (las tres zonas) y Ansible (nodo de tienda) & Subdocumento 4, sección 4.2.10.3, Tabla 4.53 \\
\addlinespace[2pt]
RT-03.04 & La red se segmentará por capas, con subredes privadas para aplicación y datos, y exposición pública restringida a la capa de borde. Ningún componente de datos será alcanzable desde Internet. & Sí & Segmentación de red por zonas (Z1/Z2/Z3), superficie nula declarada para Z2 y Z3; solo la capa de borde publica hacia Internet & Subdocumento 4, sección 4.1.7, Tablas 4.24 y 4.25 \\
\addlinespace[2pt]
RT-03.05 & El PROPONENTE privilegiará servicios administrados por sobre servicios autoadministrados cuando ello reduzca el riesgo operacional, y justificará cada excepción. & Sí & Criterio de selección N°2: ``servicio gestionado por sobre autoadministrado'', con excepciones justificadas (Z3) & Subdocumento 4, sección 4.2.10.1 \\
\addlinespace[2pt]
RT-03.06 & Se aplicarán prácticas FinOps: etiquetado obligatorio de todos los recursos por ambiente, módulo y centro de costo; presupuestos con alertas de desviación; y reporte mensual de consumo desglosado entregado al CLIENTE. & No & - & - \\
\addlinespace[2pt]
RT-03.07 & El PROPONENTE declarará su estrategia de reversibilidad y de mitigación del bloqueo por proveedor, identificando qué componentes son portables, cuáles no lo son y cuál sería el esfuerzo estimado de una migración. & Sí & Declaración de modelo de consumo y portabilidad por componente (Alta/Media-Alta/Total/Baja) & Subdocumento 4, sección 4.2.10.5, Tabla 4.54 \\
\addlinespace[2pt]
RT-03.08 & La solución empleará instancias reservadas, planes de ahorro o capacidad comprometida cuando el perfil de carga lo justifique, y lo reflejará en la estructura de costos. & No & - & - \\
\addlinespace[2pt]
RT-03.09 & Se valorará el uso de cómputo sin servidor o de contenedores administrados para las cargas de perfil variable, con el análisis comparativo de costo frente a instancias permanentes. & No & - & - \\
\addlinespace[2pt]
RT-03.10 & El componente on-premise operará de forma autónoma y degradada ante la pérdida total del enlace con la nube, durante un período mínimo de 24 horas continuas o el mayor que fije el caso. & Sí & Nodo de tienda con autonomía de 8 h (valor del caso, RT-03.10-C) y centro de distribución con 4 h & Subdocumento 4, sección 4.2.3 \\
\addlinespace[2pt]
RT-03.11 & Durante la operación desconectada, la solución continuará registrando las transacciones operacionales críticas de forma local, con integridad garantizada y sin pérdida de datos. & Sí & Registro local con outbox transaccional sobre PostgreSQL local del nodo de tienda & Subdocumento 4, sección 4.2.3 \\
\addlinespace[2pt]
RT-03.12 & Restablecido el enlace, la sincronización será automática, con reconciliación determinista de conflictos, regla de resolución documentada y bitácora auditable de las decisiones aplicadas. & Sí & Drenaje idempotente por identificador de transacción local, con resolución determinista y bitácora auditable & Subdocumento 4, secciones 4.1.6.5 y 4.2.3 \\
\addlinespace[2pt]
RT-03.13 & El PROPONENTE declarará qué funciones NO estarán disponibles en modo desconectado y qué procedimiento manual las suple. La ausencia de esta declaración se evaluará como observación grave. & Sí & Tabla de funciones no disponibles offline y procedimiento manual de respaldo por función & Subdocumento 4, sección 4.2.3, Tabla 4.31 \\
\addlinespace[2pt]
RT-03.14 & Los equipos on-premise críticos serán redundantes. El almacenamiento local tolerará la falla de al menos un disco; el PROPONENTE declarará el nivel RAID escogido y lo justificará frente a las alternativas. & Sí & Arreglo NVMe con RAID 10 (o erasure coding en objetos) con monitoreo predictivo de salud & Subdocumento 4, sección 4.2.11.1, Tabla 4.55 \\
\addlinespace[2pt]
RT-03.15 & Los sistemas on-premise se endurecerán conforme a los CIS Benchmarks aplicables, con gestión centralizada de parches y ventana de aplicación acordada con el CLIENTE. & No & - & - \\
\addlinespace[2pt]
RT-03.16 & El monitoreo del componente on-premise se integrará a la misma plataforma de observabilidad que la nube, con alertamiento unificado. & Sí & Observabilidad y Degradación Controlada (M-24), OpenTelemetry unificado en nube y nodos de tienda & Subdocumento 4, secciones 4.1.3 y 4.2.10.3 \\
\addlinespace[2pt]
RT-03.17 & El enlace entre el sitio on-premise y la nube será redundante, con caminos físicos y proveedores distintos, y conmutación automática con tiempo de conmutación declarado. & Sí & Enlaces primario y de respaldo por sitio, de proveedores distintos; conmutación automática SD-WAN en Coyhaique & Subdocumento 4, secciones 4.2.6 y 4.2.9.8, Tabla 4.38 \\
\addlinespace[2pt]
RT-03.18 & Los dispositivos de borde y de terreno se administrarán de forma remota y centralizada: inventario, configuración, actualización de firmware y de aplicación, bloqueo y borrado remoto. & Sí & Plataforma de gestión centralizada de flota, unificando dispositivos móviles, de terreno y estaciones de trabajo & Subdocumento 4, sección 4.2.11.3 \\
\addlinespace[2pt]
RT-03.19 & Se valorará el procesamiento en el borde de las cargas que lo admitan —filtrado, agregación previa, inferencia local— reduciendo el volumen transferido y la dependencia del enlace. & No & - & - \\
\addlinespace[2pt]
RT-03.20 & El PROPONENTE dimensionará el ancho de banda requerido por sitio, en régimen normal y en peak, y lo justificará con el cálculo de volumen de transacciones y de datos. & Sí & Dimensionamiento de ancho de banda por enlace, derivado de la volumetría de sistema & Subdocumento 4, secciones 4.2.6 y 4.2.9.8, Tabla 4.51 \\
\addlinespace[2pt]
RT-03.21 & La conexión entre la red del CLIENTE y la nube se establecerá mediante enlace privado dedicado o red privada virtual con cifrado, según lo que el volumen y la criticidad justifiquen. & Sí & VPN IPsec para tiendas; enlace privado dedicado (Azure ExpressRoute) para Site Principal & Subdocumento 4, sección 4.2.6 \\
\addlinespace[2pt]
RT-03.22 & El acceso remoto de las personas trabajadoras del CLIENTE, incluido el trabajo desde el hogar, se resolverá con acceso a la red de confianza cero, con verificación de postura del dispositivo. No se admite exponer servicios internos directamente a Internet. & Sí & Acceso a red de confianza cero (ZTNA) con verificación de postura del dispositivo & Subdocumento 4, sección 4.1.7, Tabla 4.25 \\
\addlinespace[2pt]
RT-03.23 & La red inalámbrica de los sitios operacionales, cuando el caso la requiera, contará con segmentación por tipo de dispositivo, autenticación por certificado o credencial de empresa y cobertura verificada mediante estudio de sitio. & No & - & - \\
\addlinespace[2pt]
RT-03.24 & El PROPONENTE declarará la calidad de servicio y la priorización de tráfico aplicada a las transacciones operacionales críticas frente al tráfico administrativo. & Sí & Clasificación de tráfico por marcado DSCP (reenvío acelerado para transacción crítica) & Subdocumento 4, sección 4.2.6 \\
\addlinespace[2pt]
RT-04.01 & Los cinco ambientes del numeral 4.1 estarán habilitados y operativos como condición del hito H3 del Formulario E-25. & Sí & Los cinco ambientes obligatorios (Desarrollo, QA, Preproducción, Producción, DR) declarados con su emplazamiento & Subdocumento 4, sección 4.2.8.1, Tabla 4.42 \\
\addlinespace[2pt]
RT-04.02 & Preproducción será equivalente a producción en topología, versiones de componentes y configuración. Las diferencias que subsistan por costo se declararán expresamente y se justificarán. & Sí & Topología equivalente a Producción, con diferencia declarada y justificada (separación física reducida en Preproducción) & Subdocumento 4, sección 4.2.8.1 \\
\addlinespace[2pt]
RT-04.03 & El código residirá en un sistema de control de versiones con ramas protegidas, revisión obligatoria por pares y prohibición de escritura directa sobre la rama principal. & Sí & Rama principal protegida, revisión obligatoria de todo cambio, incluidos los de infraestructura como código & Subdocumento 4, sección 4.2.8.2, Tabla 4.43 \\
\addlinespace[2pt]
RT-04.04 & Existirá trazabilidad completa entre requerimiento, incidencia, cambio de código, prueba ejecutada y despliegue realizado. & Sí & Trazabilidad apoyada en el catálogo de servicios versionado junto con el código & Subdocumento 4, sección 4.2.8.2, Tabla 4.43 \\
\addlinespace[2pt]
RT-04.05 & El flujo de integración continua ejecutará, como mínimo: compilación, pruebas unitarias, análisis estático de código, análisis de composición de software, escaneo de secretos y escaneo de imágenes de contenedor, con criterios de bloqueo automático del despliegue. & Sí & Cadena DevSecOps (Semgrep, ZAP, Trivy, OWASP Dependency-Check, Syft, Cosign) con bloqueo automático & Subdocumento 4, secciones 4.1.7 y 4.2.10.3, Tabla 4.53 \\
\addlinespace[2pt]
RT-04.06 & Los despliegues serán automatizados y reproducibles, con reversión automatizada y sin intervención manual en el paso a producción. & Sí & Despliegue automatizado (Argo CD, GitOps) para Z1/Z2; punto de aprobación manual adicional solo en Z3 (dominio financiero) & Subdocumento 4, sección 4.2.8.2, Tabla 4.43 \\
\addlinespace[2pt]
RT-04.07 & La estrategia de despliegue permitirá liberar sin interrupción del servicio: azul-verde, canario o despliegue progresivo. Se declarará cuál se emplea y se demostrará en Preproducción antes de cada paso a producción. & Sí & Despliegue azul-verde para Z1/Z2; actualización progresiva de instancias para Z3 & Subdocumento 4, sección 4.2.8.2, Tabla 4.43 \\
\addlinespace[2pt]
RT-04.08 & Toda configuración estará externalizada del artefacto y gestionada por ambiente. Un mismo artefacto deberá poder promoverse de QA a Preproducción y a Producción sin recompilación. & Sí & Gestor de configuración centralizado por ambiente, separado del artefacto de aplicación & Subdocumento 4, sección 4.2.8.2, Tabla 4.43 \\
\addlinespace[2pt]
RT-04.09 & Los secretos residirán en un gestor de secretos con rotación automática y auditoría de acceso. Queda prohibida toda credencial embebida en código, imágenes o archivos de configuración. & Sí & Azure Key Vault (aplicación) + HSM dedicado (filial), con rotación declarada de claves y credenciales & Subdocumento 4, secciones 4.1.7 y 4.2.10.3, Tabla 4.53 \\
\addlinespace[2pt]
RT-04.10 & Las migraciones de esquema de base de datos serán versionadas, reversibles y ejecutadas de forma automatizada, con estrategia de compatibilidad que permita convivir dos versiones de la aplicación durante el despliegue. & Sí & Migraciones de esquema coherentes con la regla de evolución de contratos por adición sin sustracción & Subdocumento 4, sección 4.2.8.2, Tabla 4.43 \\
\addlinespace[2pt]
RT-04.11 & La cobertura de pruebas automatizadas del código de lógica de negocio será de al menos 70\%, con umbral bloqueante en el flujo de integración continua. & Sí & Cobertura mínima de 70\% declarada como criterio de aceptación técnica bloqueante en CI & Subdocumento 4, sección 4.2.8.2, Tabla 4.43 \\
\addlinespace[2pt]
RT-04.12 & El PROPONENTE declarará su frecuencia de despliegue objetivo, su tiempo desde el compromiso de código hasta producción, su tasa de cambios fallidos y su tiempo de restauración, y los medirá durante la Operación. & Sí & Los cuatro indicadores DORA (frecuencia de despliegue, lead time, tasa de cambios fallidos, tiempo de restauración) & Subdocumento 4, sección 4.2.8.2, Tabla 4.43 \\
\addlinespace[2pt]
RT-04.13 & Los ambientes no productivos se apagarán o reducirán fuera del horario de uso, con el ahorro reflejado en la estructura de costos. & Sí & Desarrollo y QA reducidos fuera de la ventana 08:00–20:00; Preproducción activa solo en ventanas de prueba & Subdocumento 4, sección 4.2.8.2, Tabla 4.43 \\
\addlinespace[2pt]
RT-04.14 & Se valorará la existencia de ambientes efímeros por rama o por incidencia, creados y destruidos automáticamente. & Sí & Ambientes efímeros por rama para los servicios sin estado de Z1, creados/destruidos por la canalización de CI & Subdocumento 4, sección 4.2.8.2, Tabla 4.43 \\
\addlinespace[2pt]
RT-05.01 & El PROPONENTE entregará el modelo de datos documentado y un diccionario de datos con el nombre, el tipo, el dominio de valores, la obligatoriedad, el propietario y la sensibilidad de cada atributo. & Sí & Modelo de datos por dominio (entidades, dueño, módulo, naturaleza del dato) con ejemplo de atributos versionados del Precio Vigente & Subdocumento 5, secciones 5.1 (Tablas 5.1-5.3) y 5.5 \\
\addlinespace[2pt]
RT-05.02 & El PROPONENTE justificará la selección del paradigma y del motor de persistencia: relacional o no relacional, garantías transaccionales, y la posición escogida entre consistencia y disponibilidad conforme al teorema CAP, para cada dominio de datos. & Sí & Selección de paradigma/motor por entidad con postura CAP explícita (CP retail transaccional, AP vitrina, CP financiero) & Subdocumento 5, sección 5.3, Tablas 5.4 y 5.5 \\
\addlinespace[2pt]
RT-05.03 & Toda operación de negocio será trazable: la solución permitirá reconstruir quién, qué, cuándo, desde qué dispositivo y con qué valores anteriores y posteriores, para cualquier registro y en cualquier momento del período de retención. & No & - & - \\
\addlinespace[2pt]
RT-05.04 & La calidad de datos se gestionará conforme a ISO/IEC 25012, con validación en el punto de captura, indicadores de completitud, exactitud y consistencia, y tablero de calidad disponible para el CLIENTE. & Sí & Reglas de calidad de datos (completitud ``publicable'', clasificación de merma, inmutabilidad del precio) y tablero de calidad & Subdocumento 5, sección 5.8.1; Subdocumento 4, Tabla 4.40 \\
\addlinespace[2pt]
RT-05.05 & El almacenamiento transaccional y el analítico estarán separados. Ninguna consulta analítica podrá degradar el desempeño de la operación. & Sí & Separación OLTP/OLAP, comunicación exclusivamente asíncrona vía bus de eventos, prohibición de lecturas analíticas sobre producción & Subdocumento 5, sección 5.10 \\
\addlinespace[2pt]
RT-05.06 & La solución permitirá exportar la totalidad de la información del CLIENTE en formatos abiertos y documentados, en cualquier momento del Contrato, sin costo adicional y sin intervención del ADJUDICATARIO. & No & - & - \\
\addlinespace[2pt]
RT-05.07 & El PROPONENTE declarará la política de retención, archivado y eliminación por cada dominio de datos, coherente con la normativa aplicable al caso, e implementará un procedimiento verificable de eliminación segura. & Sí & Política de retención por dominio de datos, archivado en frío y eliminación segura (borrado criptográfico/físico con certificado) & Subdocumento 5, secciones 5.8.2, 5.8.3 y 5.8.4 \\
\addlinespace[2pt]
RT-05.08 & Los datos personales se tratarán conforme al Artículo 85° de las Bases Administrativas, con seudonimización o cifrado a nivel de campo para las categorías sensibles que el caso identifique. & Sí & Cifrado de campo y tokenización de identificación del cliente, cartera, comportamiento de pago y evaluación crediticia & Subdocumento 4, sección 4.1.7 \\
\addlinespace[2pt]
RT-05.09 & El PROPONENTE presentará una estrategia de gestión de datos maestros que evite la duplicación de entidades compartidas entre módulos y con sistemas externos. & Sí & Principio de propiedad del dato (Data Ownership) y Maestro de Artículos (M-03) como dueño único de la entidad Catálogo & Subdocumento 5, sección 5.1 \\
\addlinespace[2pt]
RT-05.10 & Se valorará la implementación de un catálogo de datos con linaje automatizado, que permita rastrear el origen de cada indicador de negocio hasta su fuente. & Sí & Catálogo de datos con linaje automatizado (origen, transformaciones Spark, consumo por tablero/modelo) & Subdocumento 5, sección 5.11.3 \\
\addlinespace[2pt]
RT-05.11 & El PROPONENTE presentará un plan de migración con alcance, origen, volumen, reglas de transformación, criterios de calidad, estrategia de ejecución y plan de reversión. & Sí & Estrategia de migración en cuatro fases (extracción, saneamiento, coexistencia/conciliación, corte) con circuit breaker y rollback & Subdocumento 5, sección 5.7 \\
\addlinespace[2pt]
RT-05.12 & La migración incluirá una etapa de perfilado y saneamiento previo, con informe de los defectos detectados en los datos de origen y la decisión adoptada sobre cada uno. & Sí & Fase 2 - Saneamiento: reconciliación de stock, deduplicación de SKU, reconstrucción de historial de precio, normalización de cartera & Subdocumento 5, sección 5.7 (Fase 2) \\
\addlinespace[2pt]
RT-05.13 & Se ejecutarán al menos dos ensayos completos de migración sobre Preproducción antes de la migración definitiva, con medición del tiempo total y del resultado de la conciliación. & No & - & - \\
\addlinespace[2pt]
RT-05.14 & La conciliación posterior a la migración será cuantitativa y verificable: recuentos, sumas de control y muestreo dirigido. Toda diferencia deberá quedar explicada. & Sí & Conteos de control (row counts), checksums monetarias/de stock y conciliación cruzada diaria con circuit breaker & Subdocumento 5, sección 5.7 (Fase 3) \\
\addlinespace[2pt]
RT-05.15 & Los datos históricos que no se migren quedarán accesibles en un repositorio de consulta durante el período de retención que fije el caso. & No & - & - \\
\addlinespace[2pt]
RT-05.16 & Los servicios síncronos se documentarán en OpenAPI 3.1 y los flujos dirigidos por eventos en AsyncAPI 2.6 o superior. La documentación se generará desde el código y se mantendrá actualizada automáticamente. & Sí & Catálogo central de contratos versionado junto con el código, en OpenAPI 3.1 y AsyncAPI 2.6 & Subdocumento 4, sección 4.1.6.4 \\
\addlinespace[2pt]
RT-05.17 & Los contratos de interfaz se versionarán semánticamente, con compatibilidad hacia atrás y política de obsolescencia con preaviso mínimo de seis meses. & Sí & Reglas de evolución de contratos (adición sin sustracción); versión antigua se retira solo con cero consumidores registrados & Subdocumento 4, sección 4.1.6.4, Tabla 4.20 \\
\addlinespace[2pt]
RT-05.18 & La autenticación entre sistemas empleará OAuth 2.1 con credenciales de cliente o autenticación mutua TLS. Queda prohibida la autenticación por clave estática en la ruta de la dirección web. & No & - & - \\
\addlinespace[2pt]
RT-05.19 & Toda integración registrará la transacción de entrada y de salida, con identificador de correlación común que permita seguir una operación de negocio a través de todos los sistemas involucrados. & No & - & - \\
\addlinespace[2pt]
RT-05.20 & Las integraciones con sistemas heredados o de terceros se aislarán mediante una capa anticorrupción, de modo que un cambio en el sistema externo no propague su modelo al núcleo de la solución. & Sí & Capa anticorrupción frente a cada sistema heredado que se conserva & Subdocumento 4, sección 4.1.1.3 \\
\addlinespace[2pt]
RT-05.21 & El PROPONENTE declarará, por cada integración, el modo (síncrono o asíncrono), el volumen esperado, la ventana de disponibilidad del sistema contraparte y el comportamiento de la solución cuando ese sistema no responde. & Sí & Catálogo de interfaces internas y externas con modo, volumen, ventana y comportamiento ante no respuesta & Subdocumento 4, sección 4.1.4, Tablas 4.15 y 4.16 \\
\addlinespace[2pt]
RT-05.22 & La solución soportará la carga y descarga masiva de información en formatos abiertos, con validación previa, informe de errores por registro y procesamiento parcial. & No & - & - \\
\addlinespace[2pt]
RT-05.23 & Se emplearán los estándares sectoriales de intercambio que las Bases Técnicas del caso identifiquen. & No & - & - \\
\addlinespace[2pt]
RT-05.24 & Se valorará la publicación de un portal de servicios para developers con documentación navegable, ambiente de pruebas y credenciales de prueba autoservidas. & No & - & - \\
\addlinespace[2pt]
RT-05.25 & La solución proveerá una capa analítica con tableros operacionales y de gestión, construidos sobre los indicadores que las Bases Técnicas del caso definan. & Sí & Capa de consumo (gold) con modelos dimensionales por dominio expuestos a tableros gerenciales & Subdocumento 5, sección 5.11.1 \\
\addlinespace[2pt]
RT-05.26 & Los tableros permitirán filtrar por período, unidad organizacional y dimensiones propias del caso, y profundizar desde el indicador agregado hasta la transacción de origen. & No & — & — \\
\addlinespace[2pt]
RT-05.27 & El CLIENTE podrá construir sus propios informes sin intervención del ADJUDICATARIO, mediante una herramienta de autoservicio con modelo semántico documentado. & Sí & BI de autoservicio sobre la capa de consumo (gold), con control de acceso basado en roles & Subdocumento 5, sección 5.11.1 \\
\addlinespace[2pt]
RT-05.28 & Todo informe será exportable en formatos abiertos y programable para envío automático por calendario. & No & - & - \\
\addlinespace[2pt]
RT-05.29 & La latencia máxima entre la ocurrencia de una transacción y su disponibilidad en la capa analítica será la que fije el caso y, en su defecto, no superará las 4 horas. & No & - & - \\
\addlinespace[2pt]
RT-05.30 & Se valorará la incorporación de analítica predictiva pertinente al proceso del caso, con el modelo, sus variables, su métrica de desempeño y su plan de reentrenamiento documentados. & No & - & - \\
\addlinespace[2pt]
RT-06.01 & El espacio asignado será de uso exclusivo de la solución y estará aislado de otras dependencias del CLIENTE, con acceso independiente. & Sí & Recinto de uso exclusivo, aislado, con acceso independiente (tipología Sala técnica principal) & Subdocumento 4, sección 4.2.4.2 \\
\addlinespace[2pt]
RT-06.02 & Los muros no estructurales del recinto contarán con blindaje perimetral; el PROPONENTE especificará el material y la resistencia. & Sí & Blindaje perimetral en muros no estructurales de la zona financiera y de la zona neutral & Subdocumento 4, sección 4.2.4.2 \\
\addlinespace[2pt]
RT-06.03 & El PROPONENTE entregará el plano de distribución interna del recinto, con la separación de las zonas de generadores, baterías, climatización, servidores, comunicaciones, trabajo y respaldo. & Sí & Plano de distribución interna con cinco zonas declaradas & Subdocumento 4, sección 4.2.4.2, Tabla 4.32 \\
\addlinespace[2pt]
RT-06.04 & El piso técnico, la canalización, el cableado estructurado y el etiquetado se ejecutarán conforme a norma, con documentación de la certificación de cada enlace. & Sí & Piso técnico, canalización y cableado estructurado ejecutados y certificados conforme a norma & Subdocumento 4, sección 4.2.4.2 \\
\addlinespace[2pt]
RT-06.05 & Los racks de servidores serán independientes de los racks de equipos de comunicación. Se declarará la ocupación proyectada de cada rack y su margen de crecimiento. & Sí & Racks de servidores independientes de los de comunicaciones, con ocupación proyectada y margen de 40\% & Subdocumento 4, sección 4.2.4.2 \\
\addlinespace[2pt]
RT-06.06 & La obra civil de separación de las instalaciones es de cargo del CLIENTE; su especificación técnica y su coordinación son de cargo del PROPONENTE. & Sí & Especificación técnica y coordinación de la obra civil de separación de zonas a cargo del PROPONENTE & Subdocumento 4, sección 4.2.4.2 \\
\addlinespace[2pt]
RT-06.07 & El suministro eléctrico de los equipos será ininterrumpido, con sistema de alimentación ininterrumpida dimensionado para una autonomía mínima de 30 minutos a plena carga. & Sí & UPS APC/Schneider Electric (línea Galaxy) con 30 minutos de autonomía a plena carga & Subdocumento 4, sección 4.2.4.3 \\
\addlinespace[2pt]
RT-06.08 & La capacidad de generación autónoma asegurará un rango mínimo de 24 horas continuas de operación, con estanque de combustible dimensionado y contrato de reabastecimiento declarado. & Sí & Grupo electrógeno con 24 h de autonomía, estanque dimensionado y contrato de reabastecimiento & Subdocumento 4, sección 4.2.4.3 \\
\addlinespace[2pt]
RT-06.09 & La instalación eléctrica del recinto será independiente de la del resto del edificio y cumplirá la normativa eléctrica chilena vigente, incluida la NCh Elec. 2777 sobre sistemas de puesta a tierra. & Sí & Acometida propia del recinto, independiente del resto del edificio, conforme a NCh Elec. 2777 & Subdocumento 4, sección 4.2.4.3 \\
\addlinespace[2pt]
RT-06.10 & Se efectuará revisión y medición semestral de las instalaciones eléctricas del recinto, con informe entregable al CLIENTE. & Sí & Revisión semestral programada en el calendario de mantención & Subdocumento 4, sección 4.2.4.3 \\
\addlinespace[2pt]
RT-06.11 & El PROPONENTE declarará la carga eléctrica proyectada en kW, el factor de potencia y la eficiencia en el uso de la energía (PUE) estimada del recinto. & Sí & Carga eléctrica proyectada de 18 kW, PUE objetivo de 1,6 & Subdocumento 4, sección 4.2.4.3 \\
\addlinespace[2pt]
RT-06.12 & Se valorará la redundancia de alimentación en configuración 2N o N+1 con doble acometida y transferencia automática. & Sí & Doble acometida con transferencia automática para la zona financiera & Subdocumento 4, sección 4.2.4.3 \\
\addlinespace[2pt]
RT-06.13 & El recinto contará con climatización de precisión para operación continua, redundante en configuración N+1, con control de temperatura y de humedad relativa dentro de los rangos que recomienda el fabricante del equipamiento. & Sí & Climatización de precisión N+1 (referenciada en 4.2.4.4) & Subdocumento 4, sección 4.2.4.4 \\
\addlinespace[2pt]
RT-06.14 & Se monitorearán en línea la temperatura, la humedad y la presencia de agua, con alertamiento integrado a la plataforma de observabilidad. & Sí & Monitoreo en línea de temperatura, humedad y presencia de agua, integrado a M-24 & Subdocumento 4, sección 4.2.4.4 \\
\addlinespace[2pt]
RT-06.15 & El PROPONENTE declarará la estrategia de contención de pasillo frío o caliente y su efecto en la eficiencia energética. & Sí & Contención de pasillo frío/caliente referenciada en el cálculo de PUE & Subdocumento 4, sección 4.2.4.3 y 4.2.4.4 \\
\addlinespace[2pt]
RT-06.16 & El recinto contará con detección temprana por aspiración de aire con tecnología láser, tipo AnaLASER o equivalente de prestaciones iguales o superiores. & Sí & Detección temprana por aspiración de aire (referenciada en 4.2.4.5) & Subdocumento 4, sección 4.2.4.5 \\
\addlinespace[2pt]
RT-06.17 & La extinción será automática mediante agente limpio tipo FM-200 o equivalente, con aprobación UL e instalación conforme a norma NFPA. & Sí & Extinción automática por agente limpio (referenciada en 4.2.4.5) & Subdocumento 4, sección 4.2.4.5 \\
\addlinespace[2pt]
RT-06.18 & Se proveerá un sistema secundario de extintores portátiles habilitados, con mantención y certificación vigentes. & Sí & Sistema secundario de extintores portátiles (referenciado en 4.2.4.5) & Subdocumento 4, sección 4.2.4.5 \\
\addlinespace[2pt]
RT-06.19 & El sistema de detección y extinción se integrará al monitoreo en línea y notificará al NOC y a la contraparte del CLIENTE. & Sí & Integración del sistema de detección y extinción al monitoreo en línea & Subdocumento 4, sección 4.2.4.5 \\
\addlinespace[2pt]
RT-06.20 & El ingreso al recinto se controlará mediante seguridad física y control de acceso biométrico basado principalmente en biometría facial, con AFIS como respaldo. Se admite proponer sistemas de mayor seguridad. & Sí & Control de acceso biométrico independiente por zona (referenciado en 4.2.4.2 y 4.2.4.6) & Subdocumento 4, secciones 4.2.4.2 y 4.2.4.6 \\
\addlinespace[2pt]
RT-06.21 & Todo ingreso y egreso quedará registrado en una bitácora auditable, con identificación de la persona, fecha, hora y motivo, conservada por el período de retención declarado. & Sí & Bitácora auditable de ingreso/egreso (referenciada en 4.2.4.6) & Subdocumento 4, sección 4.2.4.6 \\
\addlinespace[2pt]
RT-06.22 & Entre el acceso principal y el término del pasillo de la zona de control se dispondrá un espacio para la atención de personas en proceso de enrolamiento. Se evaluará mejor la existencia de una estación de enrolamiento fuera de las instalaciones del recinto técnico. & Sí & Espacio de enrolamiento (referenciado en 4.2.4.6) & Subdocumento 4, sección 4.2.4.6 \\
\addlinespace[2pt]
RT-06.23 & Al término del pasillo se instalará un acceso que impida el paso de más de una persona a la vez, con nueva verificación de identidad previa al ingreso. & Sí & Control de esclusa de una persona con verificación de identidad (referenciado en 4.2.4.6) & Subdocumento 4, sección 4.2.4.6 \\
\addlinespace[2pt]
RT-06.24 & El recinto contará con videovigilancia y monitoreo IP, con imágenes en línea y disponibles para visualización de al menos los últimos 30 días. Las grabaciones anteriores se respaldarán en un medio secundario recuperable y auditable. & Sí & Videovigilancia y monitoreo IP con retención de 30 días (referenciado en 4.2.4.6) & Subdocumento 4, sección 4.2.4.6 \\
\addlinespace[2pt]
RT-06.25 & El PROPONENTE declarará el procedimiento de acceso de terceros —fabricantes, mantenedores, auditores— con acompañamiento obligatorio y registro. & Sí & Procedimiento de acceso de terceros con acompañamiento y registro (referenciado en 4.2.4.6) & Subdocumento 4, sección 4.2.4.6 \\
\addlinespace[2pt]
RT-06.26 & Se habilitará un servicio de custodia de medios de respaldo para el sitio primario, en un medio físico transportable a otro lugar cuando el CLIENTE lo determine. Se admite proponer una solución más segura y eficiente, debidamente justificada. & Sí & Servicio de custodia de medios de respaldo en formato físico transportable & Subdocumento 4, sección 4.2.4.7 \\
\addlinespace[2pt]
RT-06.27 & El recinto de custodia cumplirá exigencias de luminosidad, humedad, ventilación y cualquier otro factor que pueda afectar la calidad y la disponibilidad de los medios. & Sí & Condiciones de luminosidad, humedad y ventilación adecuadas para la custodia & Subdocumento 4, sección 4.2.4.7 \\
\addlinespace[2pt]
RT-06.28 & Se llevará un inventario de medios con rotación, verificación periódica de legibilidad y registro de todo movimiento de entrada y de salida. & Sí & Inventario de medios con rotación y verificación periódica de legibilidad & Subdocumento 4, sección 4.2.4.7 \\
\addlinespace[2pt]
RT-06.29 & El PROPONENTE habilitará el espacio físico necesario para el personal encargado de la operación y administración de la plataforma, con estaciones de trabajo, telefonía, conexión a Internet y todo elemento que permita realizar la labor en condiciones adecuadas. & Sí & Espacio de operación con estaciones de trabajo, telefonía y conectividad, ≈20 m² & Subdocumento 4, sección 4.2.4.8 \\
\addlinespace[2pt]
RT-06.30 & El espacio de operación estará separado de la sala de equipos y no requerirá el ingreso al recinto técnico para las labores habituales de operación. & Sí & Espacio de operación separado de la sala de equipos, sin necesidad de ingreso al recinto técnico & Subdocumento 4, sección 4.2.4.8 \\
\addlinespace[2pt]
RT-06.31 & Las instalaciones sanitarias, las zonas de seguridad ante emergencia y las áreas exteriores existentes en el edificio del CLIENTE podrán utilizarse y no deben implementarse nuevamente. El PROPONENTE declarará de cuáles hará uso. & Sí & Uso de instalaciones sanitarias y zonas de emergencia ya existentes en el edificio de casa matriz & Subdocumento 4, sección 4.2.4.8 \\
\addlinespace[2pt]
RT-06.32 & El acceso a las redes de comunicaciones estará provisto a través de rutas físicas distintas, con ingreso al edificio por puntos separados. & Sí & Rutas físicas distintas con ingreso al edificio por puntos separados & Subdocumento 4, sección 4.2.4.9 \\
\addlinespace[2pt]
RT-06.33 & El PROPONENTE proveerá toda la conectividad, la seguridad y las canalizaciones, cañerías o ductos requeridos para cumplir los niveles de servicio comprometidos. & Sí & Provisión íntegra de conectividad, seguridad y canalizaciones & Subdocumento 4, sección 4.2.4.9 \\
\addlinespace[2pt]
RT-06.34 & Se privilegiará al PROPONENTE que ofrezca o provea especificaciones nuevas o mejores que las aquí establecidas, debidamente fundamentadas. & No & - & - \\
\addlinespace[2pt]
RT-07.01 & El PROPONENTE declarará la modalidad escogida, activo-activo o activo-pasivo, y la justificará frente al costo, al RTO comprometido y a la complejidad operacional que introduce. & Sí & Modalidad activo-pasivo, justificada frente a costo, RTO y complejidad operacional & Subdocumento 4, sección 4.2.5.1, Tabla 4.34 \\
\addlinespace[2pt]
RT-07.02 & El sitio secundario estará emplazado a una distancia suficiente del principal para no verse afectado por el mismo evento de fuerza mayor. El PROPONENTE declarará la distancia y el análisis de amenazas comunes considerado. & No & - & - \\
\addlinespace[2pt]
RT-07.03 & La replicación de datos será continua, con medición y alertamiento del retraso de replicación. & Sí & Replicación continua y asíncrona (Patroni streaming replication) con medición y alertamiento de rezago & Subdocumento 4, sección 4.2.5.3, Tabla 4.35 \\
\addlinespace[2pt]
RT-07.04 & El objetivo de tiempo de recuperación (RTO) no superará 4 horas y el objetivo de punto de recuperación (RPO) no superará 15 minutos para los servicios críticos, salvo exigencia superior del caso. & Sí & RTO ≤ 4 h y RPO ≤ 15 min comprometidos por componente, con segundo HSM en Site Secundario & Subdocumento 4, sección 4.2.5.4, Tabla 4.37 \\
\addlinespace[2pt]
RT-07.05 & El procedimiento de conmutación estará documentado, automatizado en la mayor medida posible y ejecutable por el personal del CLIENTE tras la transferencia de conocimiento. & Sí & Procedimiento de conmutación en 8 pasos, con automatización de los pasos técnicos & Subdocumento 4, sección 4.2.5.5 \\
\addlinespace[2pt]
RT-07.06 & Existirá un procedimiento de retorno al sitio principal igualmente documentado y probado, con reconciliación de los datos generados durante la contingencia. & Sí & Procedimiento de retorno al Site Principal con reconciliación por drenaje idempotente & Subdocumento 4, sección 4.2.5.6 \\
\addlinespace[2pt]
RT-07.07 & El plan de recuperación ante desastres se probará al menos dos veces al año mediante conmutación real, con informe de resultados, medición del RTO y del RPO efectivamente alcanzados y plan de corrección de las brechas detectadas. & Sí & Conmutación real dos veces al año, con medición de RTO/RPO efectivos e informe de brechas & Subdocumento 4, sección 4.2.5.7 \\
\addlinespace[2pt]
RT-07.08 & Se valorará que la conmutación sea automática ante la detección de indisponibilidad, con criterio de disparo declarado y protección contra conmutación innecesaria. & Sí & Conmutación semi-automática (detección y ejecución técnica automatizadas, decisión de disparo humana) con protección contra split-brain & Subdocumento 4, sección 4.2.5.8 \\
\addlinespace[2pt]
RT-07.09 & La política de respaldo seguirá el esquema 3-2-1-1-0: tres copias, en dos medios distintos, una fuera de sitio, una inmutable o fuera de línea y cero errores de verificación de restauración. & Sí & Política de respaldo 3-2-1-1-0 declarada por dominio de datos & Subdocumento 4, sección 4.2.7.1, Tabla 4.39 \\
\addlinespace[2pt]
RT-07.10 & Los respaldos estarán cifrados en reposo y en tránsito, con clave gestionada de forma independiente de la infraestructura respaldada. & Sí & Respaldos cifrados con clave gestionada de forma independiente (Azure Key Vault / HSM de la filial) & Subdocumento 4, sección 4.2.7.1 \\
\addlinespace[2pt]
RT-07.11 & Las copias inmutables estarán protegidas contra borrado y contra modificación durante su período de retención, incluso frente a credenciales administrativas comprometidas. & Sí & Snapshot con retención bloqueada (object lock); política WORM inmutable por diseño & Subdocumento 4, sección 4.2.7.1, Tabla 4.39 \\
\addlinespace[2pt]
RT-07.12 & Se ejecutará y documentará una prueba de restauración al menos mensual, sobre una muestra representativa, con medición del tiempo efectivo de restauración. & Sí & Prueba de restauración mensual sobre muestra representativa de cada dominio & Subdocumento 4, sección 4.2.7.1 \\
\addlinespace[2pt]
RT-07.13 & El PROPONENTE declarará, por cada dominio de datos, la frecuencia de respaldo, el período de retención y el tiempo estimado de restauración completa. & Sí & Política de respaldo declarada por dominio de datos (frecuencia, retención, tiempo de restauración) & Subdocumento 4, sección 4.2.7.1, Tabla 4.39 \\
\addlinespace[2pt]
RT-07.14 & Los respaldos permitirán la restauración granular: un registro, una tabla, un módulo o el sistema completo. & Sí & Restauración granular nativa en los cuatro dominios (respaldo lógico y snapshots) & Subdocumento 4, sección 4.2.7.1 \\
\addlinespace[2pt]
RT-08.01 & El PROPONENTE especificará el equipamiento de cómputo con su marca, modelo de referencia, procesador, memoria, almacenamiento local, interfaces y consumo, junto con el cálculo de dimensionamiento que lo sustenta. & Sí & Servidor Dell PowerEdge R660 / HPE ProLiant DL380 Gen11 o equivalente, con dimensionamiento derivado de 4.2.9.6 & Subdocumento 4, sección 4.2.11.1, Tabla 4.55 \\
\addlinespace[2pt]
RT-08.02 & El almacenamiento será redundante, con tolerancia declarada a la falla de discos, control de errores y monitoreo predictivo de salud de los medios. & Sí & Arreglo NVMe RAID 10 / erasure coding con monitoreo predictivo S.M.A.R.T. / NVMe health telemetry & Subdocumento 4, sección 4.2.11.1, Tabla 4.55 \\
\addlinespace[2pt]
RT-08.03 & Los conmutadores de núcleo, los cortafuegos y los balanceadores de carga estarán en configuración de alta disponibilidad, sin punto único de falla. & Sí & Conmutadores, cortafuegos y balanceadores clase empresarial en configuración de alta disponibilidad activo-pasivo o clúster & Subdocumento 4, sección 4.2.11.1, Tabla 4.55 \\
\addlinespace[2pt]
RT-08.04 & Todo el equipamiento contará con fuentes de poder redundantes y conexión a circuitos eléctricos distintos. & Sí & Fuente de poder dual por servidor, cada una a un circuito eléctrico distinto & Subdocumento 4, sección 4.2.11.1, Tabla 4.55 \\
\addlinespace[2pt]
RT-08.05 & El PROPONENTE declarará el margen de crecimiento del dimensionamiento propuesto, expresado como porcentaje sobre la carga proyectada del caso, y el procedimiento de ampliación. & Sí & Margen de crecimiento de 40\% con procedimiento de ampliación por adición de unidades a racks de reserva & Subdocumento 4, sección 4.2.11.1 \\
\addlinespace[2pt]
RT-08.06 & El equipamiento será nuevo, sin uso previo, con garantía de fábrica vigente desde la recepción conforme. & Sí & Equipamiento nuevo, sin uso previo, con garantía de fábrica vigente & Subdocumento 4, sección 4.2.11.1 \\
\addlinespace[2pt]
RT-08.07 & El PROPONENTE especificará las estaciones de trabajo requeridas para la operación y la administración de la plataforma, en la cantidad que determine el dimensionamiento del caso, con monitores duales. & Sí & Estación de trabajo de clase empresarial con dos monitores, dimensionada sobre 46 puestos & Subdocumento 4, sección 4.2.11.2, Tabla 4.56 \\
\addlinespace[2pt]
RT-08.08 & Los puestos de trabajo cumplirán las condiciones ergonómicas de la NCh 2527 y los equipos contarán con certificación de eficiencia energética. & Sí & Mobiliario y equipos conformes a NCh 2527, con certificación ENERGY STAR o equivalente & Subdocumento 4, sección 4.2.11.2, Tabla 4.56 \\
\addlinespace[2pt]
RT-08.09 & Las estaciones estarán gestionadas de forma centralizada, con cifrado de disco, control de dispositivos extraíbles, antivirus con detección y respuesta y actualización automatizada. & Sí & Gestión centralizada con cifrado de disco completo, control de extraíbles, EDR y actualización automatizada & Subdocumento 4, sección 4.2.11.2, Tabla 4.56 \\
\addlinespace[2pt]
RT-08.10 & El PROPONENTE especificará cada dispositivo de terreno con marca, modelo de referencia, cantidad, características mínimas, accesorios, consumibles y costo unitario estimado, aun cuando su compra sea de cargo del CLIENTE. & Sí & Especificación completa de dispositivos de terreno (terminal móvil, lector, terminal de pago, impresora, lector de CD), costo remitido a la Oferta Económica & Subdocumento 4, sección 4.2.11.3, Tabla 4.57 \\
\addlinespace[2pt]
RT-08.11 & La especificación considerará las condiciones reales de uso del caso: intemperie, humedad, polvo, vibración, temperatura, uso con guantes, luminosidad y autonomía de batería requerida por turno. & Sí & Condiciones reales de uso declaradas por dispositivo (IP54, resistencia a caída, uso con guantes, batería de repuesto) & Subdocumento 4, sección 4.2.11.3, Tabla 4.57 \\
\addlinespace[2pt]
RT-08.12 & Los dispositivos declararán su grado de protección contra polvo y agua y su resistencia a caídas, coherentes con el entorno de operación. & Sí & Grado de protección IP54 y resistencia a caída desde 1,2 m declarados & Subdocumento 4, sección 4.2.11.3, Tabla 4.57 \\
\addlinespace[2pt]
RT-08.13 & El PROPONENTE indicará el ciclo de vida esperado de cada dispositivo, la disponibilidad de repuestos y el plan de reposición durante los 56 meses del Contrato. & Sí & Ciclo de vida (4-5 años dispositivos móviles, 5-7 servidores) con renovación proyectada al mes 40 & Subdocumento 4, sección 4.2.11.5 \\
\addlinespace[2pt]
RT-08.14 & Los dispositivos se integrarán a la gestión centralizada de flota exigida en RT-03.18. & Sí & Incorporación de todos los dispositivos móviles y de terreno a la plataforma de gestión centralizada de flota & Subdocumento 4, sección 4.2.11.3 \\
\addlinespace[2pt]
RT-08.15 & El PROPONENTE proveerá una unidad de cada tipo de dispositivo especificado para pruebas de aceptación por parte del CLIENTE, antes de la compra masiva. & Sí & Provisión de una unidad de cada tipo de dispositivo de terreno para prueba de aceptación & Subdocumento 4, sección 4.2.11.3 \\
\addlinespace[2pt]
RT-08.16 & El PROPONENTE presentará el plan de ciclo de vida del equipamiento: recepción, puesta en servicio, mantención, actualización, retiro y disposición final. & Sí & Plan de ciclo de vida completo (recepción, puesta en servicio, mantención, actualización, retiro, disposición) & Subdocumento 4, sección 4.2.11.5 \\
\addlinespace[2pt]
RT-08.17 & Todo medio de almacenamiento que salga de servicio será borrado de forma segura y verificable, con certificado de destrucción o de sanitización entregado al CLIENTE. & Sí & Borrado seguro conforme a NIST SP 800-88, con certificado de destrucción/sanitización & Subdocumento 4, sección 4.2.11.5 \\
\addlinespace[2pt]
RT-08.18 & La disposición final de equipamiento electrónico se realizará con gestor autorizado, conforme a la normativa de residuos aplicable, con certificado de disposición. & Sí & Disposición final con gestor autorizado conforme a la Ley N° 20.920, con certificado por lote & Subdocumento 4, sección 4.2.11.5 \\
\addlinespace[2pt]
RT-08.19 & Se valorará una estrategia de reacondicionamiento o de extensión de vida útil que reduzca el impacto ambiental, cuantificada en la propuesta. & Sí & Reacondicionamiento certificado del fabricante para el stock de repuestos de terminales móviles y lectores & Subdocumento 4, sección 4.2.11.5 \\
\addlinespace[2pt]
RT-09.01 & El PROPONENTE presentará el cálculo de capacidad que sustenta su dimensionamiento, con los supuestos de usuarios concurrentes, transacciones por segundo, volumen de datos y crecimiento anual, tomados de la volumetría del caso. & Sí & Derivación de la volumetría de sistema (concurrencia, TPS, almacenamiento) a partir de la volumetría operacional del caso & Subdocumento 4, sección 4.0.4 (memoria de cálculo, Anexo A.2) \\
\addlinespace[2pt]
RT-09.02 & La solución soportará la concurrencia y el volumen de transacciones que fije el caso, y mantendrá los umbrales del numeral 9.1 bajo esa carga. & Sí & Dimensionamiento de cómputo, bus de eventos y motores de datos ajustado a los umbrales de latencia del caso & Subdocumento 4, sección 4.2.9, Tablas 4.45 a 4.51 \\
\addlinespace[2pt]
RT-09.03 & La solución soportará, sin rediseño, un crecimiento de al menos tres veces la volumetría inicial del caso en un horizonte de tres años. & No & - & - \\
\addlinespace[2pt]
RT-09.04 & El escalamiento de las capas de aplicación e integración será horizontal y automático, con tiempo de reacción declarado y sin pérdida de transacciones en curso. & Sí & Escalamiento automático por criterio de saturación (latencia p95, profundidad de cola) con umbrales por servicio & Subdocumento 4, sección 4.1.10 \\
\addlinespace[2pt]
RT-09.05 & EL PROPONENTE identificará el componente que primero se convertirá en cuello de botella al crecer la carga y explicará cómo lo detectará y cómo lo resolverá. & Sí & Enlace de respaldo de las 14 tiendas dependientes de red de centro comercial identificado como primer cuello de botella, con detección y mitigación & Subdocumento 4, sección 4.2.9.9 \\
\addlinespace[2pt]
RT-09.06 & Se ejecutarán pruebas de carga sobre Preproducción con un volumen equivalente a 1,5 veces el peak declarado, y pruebas de estrés hasta identificar el punto de quiebre de la solución. & Sí & Pruebas de carga a 1,5x el peak proyectado, para verificar los supuestos de volumetría & Subdocumento 4, sección 4.1.10 y Anexo A.1 (SV-01 a SV-11) \\
\addlinespace[2pt]
RT-09.07 & El informe de pruebas de carga incluirá la curva de tiempo de respuesta frente a carga, el punto de saturación, el consumo de recursos y el comportamiento durante y después del peak. & No & - & - \\
\addlinespace[2pt]
RT-09.08 & La solución degradará de forma controlada al superarse la capacidad: encolamiento, limitación de tasa y mensaje explícito a la persona usuaria, nunca error genérico ni pérdida silenciosa de transacciones. & Sí & Degradación controlada comercial (M-24): priorización de checkout sobre navegación, throttling de sesiones nuevas & Subdocumento 4, secciones 4.1.10 y 4.2.7.4 \\
\addlinespace[2pt]
RT-09.09 & Se gestionará la capacidad durante la Operación con proyección trimestral de crecimiento, alertas anticipadas de agotamiento y propuesta de ajuste de dimensionamiento y de costo. & Sí & Curva de crecimiento de almacenamiento a 3 años, reevaluada en cada revisión trimestral de capacidad & Subdocumento 4, sección 4.2.9.7, Tabla 4.50 \\
\addlinespace[2pt]
RT-09.10 & Se valorará la existencia de pruebas de carga automatizadas ejecutadas de forma periódica en el flujo de integración continua, con detección de regresiones de desempeño. & No & - & - \\
\addlinespace[2pt]
RT-10.01 & La solución alcanzará una disponibilidad mensual mínima de 99,9\% para los servicios clasificados como críticos, medida sobre la transacción de negocio de extremo a extremo. & Sí & Adopción íntegra de los niveles de disponibilidad del numeral 7.2 transversal (99,9\% transacción crítica extremo a extremo) & Subdocumento 4, sección 4.2.7.2 \\
\addlinespace[2pt]
RT-10.02 & El PROPONENTE clasificará cada servicio de la solución en crítico, alto, medio o bajo, justificando la clasificación con el impacto operacional de su indisponibilidad, y aplicará a cada uno el nivel de servicio correspondiente del Artículo 78° de las Bases Administrativas. & Sí & Clasificación de módulos en crítico/alto/medio/bajo con justificación de impacto & Subdocumento 4, sección 4.2.7.3, Tabla 4.40 \\
\addlinespace[2pt]
RT-10.03 & El plan de continuidad del negocio se elaborará conforme a ISO 22301, con análisis de impacto en el negocio, escenarios de contingencia, procedimientos manuales de respaldo y criterios de activación. & Sí & Plan de continuidad del negocio con tres escenarios de contingencia, activación y procedimiento manual de respaldo & Subdocumento 4, sección 4.2.7.4, Tabla 4.41 \\
\addlinespace[2pt]
RT-10.04 & La continuidad TIC se estructurará conforme a ISO/IEC 27031, articulada con el plan de recuperación ante desastres del Capítulo 7. & Sí & Continuidad TIC (ISO/IEC 27031) articulada con el plan de recuperación de sitio de la sección 4.2.5 & Subdocumento 4, sección 4.2.7.4 \\
\addlinespace[2pt]
RT-10.05 & Los mantenimientos programados se ejecutarán fuera de la ventana operacional crítica que defina el caso, con aviso previo mínimo de diez días hábiles. & No & - & - \\
\addlinespace[2pt]
RT-10.06 & La solución permitirá desplegar cambios sin interrupción del servicio. Las ventanas de indisponibilidad programada serán excepcionales y deberán justificarse caso a caso. & Sí & Despliegue azul-verde / actualización progresiva sin ventana de indisponibilidad & Subdocumento 4, sección 4.2.8.2, Tabla 4.43 (RT-04.07) \\
\addlinespace[2pt]
RT-10.07 & Se ejecutarán pruebas de resiliencia mediante inyección controlada de fallas —caída de instancia, de zona, de dependencia externa, latencia elevada, saturación de disco— antes de cada paso a producción y al menos una vez por semestre durante la Operación. & Sí & Ingeniería del caos dirigida sobre los cinco tipos de falla exigidos, antes de cada paso a producción y semestralmente & Subdocumento 4, sección 4.2.7.5 \\
\addlinespace[2pt]
RT-10.08 & El PROPONENTE documentará, por cada dependencia externa, el comportamiento de la solución cuando esa dependencia no responde, responde con error o responde con lentitud. & Sí & Comportamiento ante no respuesta documentado por cada integración externa & Subdocumento 4, sección 4.1.4.2, Tabla 4.16 \\
\addlinespace[2pt]
RT-10.09 & Se declarará un presupuesto de error por servicio crítico y su vinculación con el ritmo de despliegue de cambios. & Sí & Presupuesto de error mensual de 0,1\% sobre el nivel de servicio de 99,9\%, que detiene el despliegue de cambios al agotarse & Subdocumento 4, sección 4.2.7.5 \\
\addlinespace[2pt]
RT-11.01 & La arquitectura de seguridad se basará en el modelo Zero Trust conforme a NIST SP 800-207: verificación explícita de cada solicitud, privilegio mínimo y presunción de compromiso. & Sí & Arquitectura de seguridad Zero Trust (NIST SP 800-207), denegación por omisión en la Frontera de Datos (M-21) & Subdocumento 4, sección 4.1.7 \\
\addlinespace[2pt]
RT-11.02 & El PROPONENTE entregará un modelado de amenazas documentado por cada componente y por cada integración externa, con metodología declarada (STRIDE u otra), y lo actualizará ante cada cambio arquitectónico relevante. & Sí & Modelado de amenazas con metodología STRIDE por componente e integración externa & Subdocumento 4, sección 4.1.7 \\
\addlinespace[2pt]
RT-11.03 & La información del CLIENTE se clasificará por nivel de sensibilidad, con controles diferenciados por nivel documentados en una matriz. & Sí & Matriz documentada de clasificación de datos por nivel de sensibilidad & Subdocumento 4, sección 4.1.7 \\
\addlinespace[2pt]
RT-11.04 & Existirá un programa de gestión de vulnerabilidades con escaneo continuo y plazos máximos de remediación de 7 días corridos para vulnerabilidades críticas, 15 días para altas y 30 días para medias, contados desde su publicación o detección. & Sí & Escaneo continuo de vulnerabilidades con plazos de remediación de 7/15/30 días & Subdocumento 4, sección 4.1.7 \\
\addlinespace[2pt]
RT-11.05 & El PROPONENTE mantendrá una matriz de controles de seguridad trazable a ISO/IEC 27001 e ISO/IEC 27002, indicando el control, su implementación concreta en la solución y la evidencia que lo acredita. & Sí & Matriz de controles trazable a ISO/IEC 27001:2022 e ISO/IEC 27002:2022 & Subdocumento 4, sección 4.1.7 \\
\addlinespace[2pt]
RT-11.06 & Se aplicarán los controles de ISO/IEC 27017 para los servicios en nube y de ISO/IEC 27018 para el tratamiento de datos personales en nube. & Sí & Controles ISO/IEC 27017 (nube) e ISO/IEC 27018:2025 (datos personales en nube pública) & Subdocumento 4, sección 4.1.7 \\
\addlinespace[2pt]
RT-11.07 & La publicación de servicios se realizará exclusivamente a través de la capa de borde, con red de distribución de contenidos, cortafuegos de aplicaciones web con reglas gestionadas y personalizadas, y protección contra denegación de servicio distribuida en capas 3, 4 y 7. & Sí & Publicación exclusiva por la capa de borde (CDN + WAF + DDoS L3/L4/L7) & Subdocumento 4, secciones 4.1.3 y 4.1.7 \\
\addlinespace[2pt]
RT-11.08 & El cifrado en tránsito empleará TLS 1.3, con prohibición expresa de TLS 1.0 y 1.1, conjuntos de cifrado modernos, HSTS con precarga y gestión automatizada de certificados con rotación y alerta anticipada de vencimiento. & Sí & TLS 1.3 con prohibición expresa de TLS 1.0/1.1 & Subdocumento 4, sección 4.1.7 \\
\addlinespace[2pt]
RT-11.09 & La totalidad de los datos en reposo estará cifrada, con claves gestionadas en un servicio de gestión de claves o en un módulo de seguridad de hardware, política de rotación declarada y separación de funciones en la custodia de claves. & Sí & Cifrado en reposo de toda la información persistida; Azure Key Vault (nube) y HSM dedicado (filial), con rotación declarada & Subdocumento 4, sección 4.1.7 \\
\addlinespace[2pt]
RT-11.10 & Los datos de categoría sensible que el caso identifique se cifrarán adicionalmente a nivel de campo, de modo que el acceso a la base de datos no revele su contenido. & Sí & Cifrado de campo y tokenización de identificación de cliente, cartera, comportamiento de pago y evaluación crediticia & Subdocumento 4, sección 4.1.7 \\
\addlinespace[2pt]
RT-11.11 & La puerta de enlace de servicios aplicará autenticación, autorización, cuotas, límites de tasa, validación de esquema e inspección de carga útil. & Sí & API Gateway (Kong Gateway) con autenticación, autorización, cuotas y límites de tasa en la Capa 3 & Subdocumento 4, secciones 4.1.3 (Tabla 4.14) y 4.2.10.3 \\
\addlinespace[2pt]
RT-11.12 & Los puntos de entrada públicos contarán con protección contra bots y abuso automatizado, con reto progresivo que no degrade la accesibilidad ni bloquee a personas usuarias legítimas. & No & - & - \\
\addlinespace[2pt]
RT-11.13 & El PROPONENTE declarará la superficie de exposición completa de la solución: cada nombre de dominio, puerto y servicio alcanzable desde fuera de la red del CLIENTE. & Sí & Declaración de superficie de exposición (dominios, puertos, servicios) con línea base al cierre de Etapa 1 & Subdocumento 4, sección 4.1.7, Tablas 4.24 y 4.25 \\
\addlinespace[2pt]
RT-11.14 & Los eventos de seguridad se registrarán de forma centralizada e inalterable, con retención mínima de doce meses en línea y veinticuatro meses adicionales en archivo recuperable. & Sí & Registro centralizado e inalterable de eventos de seguridad, 12 meses en línea + 24 en archivo & Subdocumento 4, sección 4.1.7 \\
\addlinespace[2pt]
RT-11.15 & Los eventos se correlacionarán en una plataforma SIEM, con casos de uso de detección definidos específicamente para el proceso de negocio del caso, y no sólo genéricos de infraestructura. & Sí & SIEM (Wazuh) con casos de uso definidos para el negocio & Subdocumento 4, secciones 4.1.7 y 4.2.10.3 \\
\addlinespace[2pt]
RT-11.16 & Se implementará detección y respuesta en puntos finales y en cargas de trabajo, tanto en nube como on-premise. & Sí & Wazuh como detección en puntos finales; EDR en estaciones de trabajo & Subdocumento 4, secciones 4.2.10.3 y 4.2.11.2 \\
\addlinespace[2pt]
RT-11.17 & El PROPONENTE dispondrá de un centro de operaciones de seguridad con cobertura 24x7, propio o subcontratado, y declarará su ubicación, dotación y procedimientos. & Sí & SOC con cobertura 24x7 & Subdocumento 4, sección 4.1.7 \\
\addlinespace[2pt]
RT-11.18 & El plan de respuesta a incidentes de seguridad definirá clasificación, cadena de escalamiento, plazos, responsables y protocolo de comunicación al CLIENTE dentro de las dos horas de detectado un incidente de severidad crítica. & Sí & Plazo de 2 horas para comunicación de incidente de severidad crítica & Subdocumento 4, sección 4.1.7 \\
\addlinespace[2pt]
RT-11.19 & Toda brecha de seguridad o de datos personales se notificará al CLIENTE dentro de las 24 horas de su detección, con informe preliminar, y el análisis de causa raíz se entregará dentro de los cinco días hábiles siguientes. & Sí & Notificación de brecha dentro de 24 horas & Subdocumento 4, sección 4.1.7 \\
\addlinespace[2pt]
RT-11.20 & Se ejecutarán pruebas de intrusión por un tercero independiente del ADJUDICATARIO, anualmente y antes de cada paso a producción, con entrega íntegra del informe al CLIENTE y plan de remediación con plazos. & Sí & Pruebas de intrusión anuales y antes de cada paso a producción & Subdocumento 4, secciones 4.1.7 y 4.1.9 \\
\addlinespace[2pt]
RT-11.21 & Se realizarán ejercicios de simulación de incidente con participación del CLIENTE al menos una vez al año durante la Operación. & No & - & - \\
\addlinespace[2pt]
RT-11.22 & El flujo de integración continua incorporará análisis estático de código, análisis de composición de software, análisis dinámico y escaneo de imágenes de contenedor, con criterios de bloqueo automático del despliegue ante hallazgos críticos. & Sí & Cadena DevSecOps (Semgrep, ZAP, Trivy, OWASP Dependency-Check) con bloqueo automático & Subdocumento 4, secciones 4.1.7 y 4.2.10.3 \\
\addlinespace[2pt]
RT-11.23 & Cada versión liberada se acompañará de su inventario de componentes de software en formato CycloneDX o SPDX, entregado al CLIENTE. & Sí & SBOM en CycloneDX o SPDX (Syft) por versión liberada & Subdocumento 4, secciones 4.1.7 y 4.2.10.3 \\
\addlinespace[2pt]
RT-11.24 & Los artefactos se firmarán y su procedencia se verificará conforme a SLSA nivel 3 o superior. & Sí & Firma y procedencia SLSA nivel 3 (Sigstore Cosign) & Subdocumento 4, secciones 4.1.7 y 4.2.10.3 \\
\addlinespace[2pt]
RT-11.25 & Queda prohibido el uso de datos productivos reales en ambientes no productivos sin anonimización o seudonimización verificable. & Sí & Prohibición de datos productivos reales en ambientes no productivos; datos sintéticos en Desarrollo y QA & Subdocumento 4, secciones 4.1.7 y 4.2.8.1, Tabla 4.42 \\
\addlinespace[2pt]
RT-11.26 & El PROPONENTE declarará el proceso de aprobación de nuevas dependencias de terceros, incluyendo criterios de licencia, mantención activa y ausencia de vulnerabilidades conocidas. & No & - & - \\
\addlinespace[2pt]
RT-11.27 & Las personas desarrolladoras no tendrán acceso interactivo directo al ambiente de producción. Todo acceso excepcional será temporal, aprobado, registrado y con sesión grabada. & Sí & Sin acceso interactivo directo a producción para developers & Subdocumento 4, sección 4.1.7 \\
\addlinespace[2pt]
RT-11.28 & Se aplicará el marco OWASP SAMM o equivalente para medir y mejorar la madurez del proceso de desarrollo seguro, con evaluación inicial y reevaluación anual. & No & - & - \\
\addlinespace[2pt]
 &  &  &  &  \\
\addlinespace[2pt]
RT-12.01 & La gestión de identidad será centralizada, con federación mediante OpenID Connect y OAuth 2.1, o SAML 2.0 cuando la integración con el CLIENTE lo requiera, e integración con el directorio corporativo del CLIENTE por LDAP o su equivalente en la nube. & Sí & Módulo de Identidad y Accesos (M-23) implementado con Keycloak, dos instancias sin confianza transitiva (Z1 y Z3), con protocolo OIDC y OAuth 2.0 & Subdocumento 4, sección 4.1.7 (punto 2, Identidad y accesos) y Tabla de stack tecnológico (fila ``Identidad y accesos (M-23)'') \\
\addlinespace[2pt]
RT-12.02 & Existirá inicio de sesión único para todos los módulos de la solución, con cierre de sesión propagado a todos ellos. & No & - & - \\
\addlinespace[2pt]
RT-12.03 & La autenticación multifactor será obligatoria para personas usuarias administradoras, para todo acceso privilegiado y para todo acceso originado fuera de la red corporativa. & No & - & - \\
\addlinespace[2pt]
RT-12.04 & Se soportarán factores resistentes a la suplantación de identidad, tipo FIDO2 o claves de acceso, al menos para los perfiles administradores. & No & - & - \\
\addlinespace[2pt]
RT-12.05 & El control de acceso será basado en roles, complementado con control basado en atributos donde el proceso lo exija, con matriz de segregación de funciones documentada y verificable. & Sí & Matriz de roles y permisos con segregación de funciones entre originación, aprobación, modificación de condiciones y cobranza (dominio financiero) & Subdocumento 3, Tabla A.2.10 ``Requerimientos No Funcionales - Seguridad'' (RNF-37) \\
\addlinespace[2pt]
RT-12.06 & Los accesos privilegiados se gestionarán con elevación temporal a demanda, aprobación previa y grabación de sesión para las operaciones de mayor riesgo. & No & - & - \\
\addlinespace[2pt]
RT-12.07 & La política de sesión declarará duración máxima, caducidad por inactividad, renovación de la credencial de sesión tras la autenticación, revocación inmediata y control de sesiones concurrentes. & No & - & - \\
\addlinespace[2pt]
RT-12.08 & Las credenciales de sesión serán firmadas y de vida breve, con credencial de refresco rotatoria. Queda prohibido transportar identificadores de sesión en la ruta de la dirección web. & Sí & IAM Keycloak con OIDC y OAuth 2.0 ``con credenciales de vida corta'' & Subdocumento 4, Tabla de stack tecnológico (fila ``Identidad y accesos (M-23)'') \\
\addlinespace[2pt]
RT-12.09 & Se registrará auditoría completa del ciclo de vida de la identidad: creación, modificación, elevación, bloqueo y baja de cuentas, con no repudio y retención declarada. & No & - & - \\
\addlinespace[2pt]
RT-12.10 & El aprovisionamiento y el desaprovisionamiento estarán automatizados y ligados al ciclo de vida laboral, con baja efectiva en un plazo no superior a 24 horas desde la desvinculación. & Sí & Conector bidireccional IAM-RRHH (M-23) que revoca accesos automáticamente por evento de finiquito, con SLA interno de revocación total ≤ 4 horas & Subdocumento 3, Tabla A.2.10 ``Requerimientos No Funcionales - Seguridad'' (RNF-59); Subdocumento 4, sección 4.1.7 (punto 2) \\
\addlinespace[2pt]
RT-12.11 & El mecanismo de autenticación se adecuará al perfil operacional real descrito en el caso: entornos de terreno, uso con guantes, baja alfabetización digital, dispositivos compartidos por turno y ausencia de correo electrónico personal. & Sí & Terminal móvil compartido con traspaso rápido de sesión (hot-swapping) sin autenticación biométrica que añada tiempo al vendedor; equipamiento con grado IP54 y uso confirmado con guantes en centro de distribución & Subdocumento 4, Tabla 4.57 ``Equipamiento de terreno y dispositivos operacionales'' (sección 4.2.11.3) \\
\addlinespace[2pt]
RT-12.12 & Las personas usuarias externas del CLIENTE —clientes, proveedores, pacientes, productores, según el caso— dispondrán de un mecanismo de registro, verificación de identidad y recuperación de acceso autoservido y seguro. & No & - & - \\
\addlinespace[2pt]
RT-12.13 & El PROPONENTE declarará el procedimiento de acceso de emergencia (cuenta de último recurso), su custodia, su control y su auditoría. & No & - & - \\
\addlinespace[2pt]
RT-13.01 & Todas las interfaces destinadas a personas usuarias cumplirán WCAG 2.2 nivel AA, verificado con herramientas automatizadas y con pruebas manuales, con informe de conformidad entregable. & No & - & - \\
\addlinespace[2pt]
RT-13.02 & El diseño será responsivo y funcionará correctamente en escritorio, tableta y teléfono, con puntos de quiebre coherentes y adaptación inteligente del contenido, no simple reducción. & No & - & - \\
\addlinespace[2pt]
RT-13.03 & El PROPONENTE ejecutará investigación con personas usuarias reales del CLIENTE, prototipado y pruebas de usabilidad antes de la construcción definitiva, y documentará los hallazgos y los cambios de diseño que produjeron. & No & - & - \\
\addlinespace[2pt]
RT-13.04 & Se comprometerán indicadores de usabilidad medibles: tiempo máximo de la transacción operacional crítica, número máximo de pasos, tasa de error tolerada y tiempo de aprendizaje esperado por perfil. & Sí & Requerimiento no funcional que fija el tiempo de capacitación operativa básica en interfaces de sala, caja y mesón financiero & Subdocumento 3, Tabla A.2.14 ``Requerimientos No Funcionales - Usabilidad'' (RNF-49) \\
\addlinespace[2pt]
RT-13.05 & Ninguna funcionalidad principal requerirá más de tres interacciones desde la pantalla de inicio del perfil correspondiente. & No & - & - \\
\addlinespace[2pt]
RT-13.06 & La solución entregará retroalimentación visual clara ante cada acción, manejará los errores de forma comprensible —indicando qué ocurrió y qué hacer— y evitará mensajes técnicos dirigidos a la persona usuaria final. & No & - & - \\
\addlinespace[2pt]
RT-13.07 & La solución soportará a personas usuarias con baja alfabetización digital: alto contraste, objetivos táctiles de al menos 44x44 píxeles, iconografía acompañada de texto, y flujos guiados paso a paso. & No & - & - \\
\addlinespace[2pt]
RT-13.08 & Cuando el caso lo requiera, las interfaces de terreno operarán con guantes, a la intemperie, con luminosidad variable y sin conexión. & Sí & Terminal móvil de consulta de existencia y lector de radiofrecuencia de centro de distribución, especificados con grado de protección IP54, pantalla legible con la luminosidad de sala y uso confirmado con guantes & Subdocumento 4, Tabla 4.57 ``Equipamiento de terreno y dispositivos operacionales'' (sección 4.2.11.3, cita explícita RT-13.08) \\
\addlinespace[2pt]
RT-13.09 & Existirá un sistema de diseño documentado con paleta de colores acotada, jerarquía tipográfica de a lo más dos familias, iconografía coherente, retícula y componentes reutilizables. & No & - & - \\
\addlinespace[2pt]
RT-13.10 & El PROPONENTE declarará la matriz de navegadores y de versiones soportadas y su política de actualización. & No & - & - \\
\addlinespace[2pt]
RT-13.11 & La solución será navegable íntegramente por teclado, con orden de foco lógico y atajos para las operaciones frecuentes. & No & - & - \\
\addlinespace[2pt]
RT-13.12 & Se valorará el soporte de modo oscuro, de personalización de la interfaz por persona usuaria y de múltiples idiomas cuando el caso lo justifique. & No & - & - \\
\addlinespace[2pt]
RT-14.01 & La observabilidad será unificada para nube y on-premise, con métricas, registros y trazas distribuidas correlacionadas por un identificador único de transacción, instrumentadas conforme a OpenTelemetry. & Sí & Módulo de Observabilidad y Degradación Controlada (M-24) instrumentado con OpenTelemetry, Prometheus, Loki y Grafana, con Azure Monitor/Application Insights como backend de trazas & Subdocumento 4, Tabla de stack tecnológico (fila ``Observabilidad (M-24)'', cita explícita RT-14.01-T) \\
\addlinespace[2pt]
RT-14.02 & El CLIENTE dispondrá de acceso propio y permanente a los tableros operacionales y de negocio, con datos en tiempo real y capacidad de exportación. & No & - & - \\
\addlinespace[2pt]
RT-14.03 & Los indicadores de nivel de servicio se medirán sobre la experiencia real de la persona usuaria y no sobre pruebas sintéticas, sin perjuicio de que estas se empleen como complemento. & No & - & - \\
\addlinespace[2pt]
RT-14.04 & El alertamiento se basará en síntomas de negocio y no sólo en umbrales de infraestructura, con supresión de ruido, agrupación, escalamiento automático y turnos de disponibilidad declarados. & No & - & - \\
\addlinespace[2pt]
RT-14.05 & Existirá un libro de operación y una guía de resolución documentados para cada escenario de falla previsible, con automatización progresiva de las tareas repetitivas. & No & - & - \\
\addlinespace[2pt]
RT-14.06 & Todo incidente crítico dará lugar a un análisis de causa raíz obligatorio, con informe entregable dentro de cinco días hábiles y seguimiento de las acciones correctivas hasta su cierre. & No & - & - \\
\addlinespace[2pt]
RT-14.07 & Los registros de la solución no contendrán datos personales sensibles ni credenciales, y su acceso estará controlado y auditado. & No & - & - \\
\addlinespace[2pt]
RT-14.08 & El PROPONENTE declarará la retención de métricas, registros y trazas, y su costo asociado, distinguiendo el almacenamiento en línea del archivado. & No & - & - \\
\addlinespace[2pt]
RT-14.09 & Se valorará la detección proactiva de anomalías mediante análisis del comportamiento histórico, con alerta antes de que el incidente afecte a la operación. & No & - & - \\
\addlinespace[2pt]
RT-15.01 & El PROPONENTE dimensionará la infraestructura ajustada a la demanda real, evitando capacidad ociosa permanente, y declarará el factor de utilización proyectado. & No & - & - \\
\addlinespace[2pt]
RT-15.02 & Los ambientes no productivos se apagarán o reducirán fuera del horario de uso. & Sí & Política de reducción de ambientes de Desarrollo y QA fuera de la ventana 08:00-20:00 de días hábiles, con Preproducción activa solo durante ventanas de prueba programadas & Subdocumento 4, sección 4.2.7 (tabla de gestión de ambientes) \\
\addlinespace[2pt]
RT-15.03 & El PROPONENTE estimará la huella de carbono anual de la operación de la solución y declarará la metodología empleada. & No & - & - \\
\addlinespace[2pt]
RT-15.04 & Se declarará la eficiencia en el uso de la energía (PUE) del recinto on-premise y la intensidad de carbono de la región de nube escogida. & Sí & Cálculo de PUE objetivo de 1,6 para el recinto on-premise (Site Principal) & Subdocumento 4, sección 4.2.x (tabla ``Carga eléctrica proyectada'', cita RT-06.11) \\
\addlinespace[2pt]
RT-15.05 & Se valorará la elección de regiones de nube con menor intensidad de carbono cuando la latencia y la regulación lo permitan, con el análisis comparativo correspondiente. & No & - & - \\
\addlinespace[2pt]
RT-15.06 & Se valorará la definición de metas de reducción del consumo durante la Operación, con medición y reporte anual. & No & - & - \\
\addlinespace[2pt]
RT-15.07 & Las certificaciones se acreditarán con copia del certificado vigente y con el código de verificación del organismo emisor cuando exista. & No & - & - \\
\addlinespace[2pt]
RT-15.08 & Las personas certificadas formarán parte del equipo efectivamente asignado al PROYECTO, con dedicación declarada. No se aceptará acreditar personal que no participe. & Sí & Equipo de jefaturas nominado con dedicación exclusiva al proyecto y certificaciones individuales (PMP, Microsoft Certified Azure Solutions Architect Expert, CISSP/CCSP, CDMP, CTAL-TM, Microsoft Certified DevOps Engineer Expert, Prosci Certified Change Practitioner) & Subdocumento 1, sección 1.6 ``Frente de jefaturas asignado'' y Tabla 1 \\
\addlinespace[2pt]
RT-15.09 & El ADJUDICATARIO mantendrá vigentes estas certificaciones durante todo el período contractual y las repondrá ante la salida de una persona certificada, conforme al Artículo 76° de las Bases Administrativas. & No & - & - \\
\addlinespace[2pt]
RT-16.01 & La solución dispondrá de un módulo de administración que permita al CLIENTE, sin intervención del ADJUDICATARIO, gestionar personas usuarias, roles, permisos, unidades organizacionales y sus jerarquías. & No & - & - \\
\addlinespace[2pt]
RT-16.02 & Las reglas de negocio parametrizables —umbrales, plazos, montos, tolerancias, catálogos, listas de valores, textos de notificación— serán configurables desde la interfaz de administración, con control de versiones y registro de quién cambió qué y cuándo. & Sí & Múltiples requerimientos funcionales que exigen parametrización de umbrales, ventanas y catálogos de negocio (colchón de reserva, ventanas de conteo cíclico, ventanas de congelamiento, plazos de garantía), junto con registro de auditoría de cambios & Subdocumento 3, Anexo A.1 (p. ej. RF-131 a RF-134, RF-184, RF-195) y Tabla A.2.7 (RNF-15) \\
\addlinespace[2pt]
RT-16.03 & Todo cambio de parámetro con impacto operacional requerirá aprobación de un segundo perfil y quedará registrado con su justificación. & No & - & - \\
\addlinespace[2pt]
RT-16.04 & El PROPONENTE declarará expresamente qué elementos son parametrizables y cuáles requieren desarrollo. Presentar como parametrizable lo que exige desarrollo se evaluará como observación grave. & Sí & Catálogo atómico de 223 requerimientos funcionales y 75 no funcionales que identifica explícitamente, requerimiento por requerimiento, cuáles son parametrizables & Subdocumento 3, Anexo A.1 y A.2 (Listado de Requerimientos Funcionales y No Funcionales) \\
\addlinespace[2pt]
RT-16.05 & Existirá un ambiente de simulación que permita probar el efecto de un cambio de parámetro antes de aplicarlo a producción. & No & - & - \\
\addlinespace[2pt]
RT-16.06 & Toda operación que cree, modifique o elimine información quedará registrada con identificación de la persona o del sistema que la ejecutó, fecha y hora con zona horaria, origen, valores anteriores y valores posteriores. & Sí & Requerimiento no funcional de auditabilidad general: toda acción relevante ejecutada en el sistema queda registrada con usuario individual, acción e instante & Subdocumento 3, Tabla A.2.7 ``Requerimientos No Funcionales - Auditabilidad'' (RNF-15) \\
\addlinespace[2pt]
RT-16.07 & El registro de auditoría será inalterable y no podrá ser modificado ni eliminado por ningún perfil, incluido el administrador de la plataforma. & Sí & Bitácora de Gobernanza y Cruces (M-21), almacén de solo anexión (append-only) con valor forense & Subdocumento 5, sección 5.1.3 ``Entidades transversales'' \\
\addlinespace[2pt]
RT-16.08 & El CLIENTE podrá consultar y exportar la auditoría desde la interfaz, con filtros por persona, período, entidad y tipo de operación, sin requerir acceso a la base de datos. & No & - & - \\
\addlinespace[2pt]
RT-16.09 & Las consultas a información sensible quedarán registradas, no sólo las modificaciones. & Sí & Gobierno de Frontera de Datos (M-21): todo acceso a la cartera de crédito, comportamiento de pago y antecedentes de evaluación crediticia queda registrado (usuario, dato, finalidad, instante) & Subdocumento 3, Tabla A.2.7 (RNF-38); Subdocumento 4, sección 4.1.7 \\
\addlinespace[2pt]
RT-16.10 & El período de retención de la auditoría será el que fije el caso y, en su defecto, no inferior a cinco años. & Sí & Política de retención de datos gobernada centralizada por dominio (documentos tributarios y de venta: 6 años; antecedentes de crédito: plazo del crédito + 6 años; precio: 3 años), fijada explícitamente por el caso & Subdocumento 5, sección 5.8.2 ``Retención'' (Tabla 5.X) \\
\addlinespace[2pt]
RT-16.11 & La solución soportará flujos de trabajo con estados, transiciones, responsables, plazos, escalamiento automático por vencimiento y delegación por ausencia. & No & - & - \\
\addlinespace[2pt]
RT-16.12 & Los flujos serán configurables por el CLIENTE sin desarrollo, al menos en lo relativo a responsables, plazos y niveles de aprobación. & No & - & - \\
\addlinespace[2pt]
RT-16.13 & Toda solicitud pendiente será visible para su responsable en una bandeja de tareas unificada, con priorización y alerta de vencimiento. & No & - & - \\
\addlinespace[2pt]
RT-16.14 & El motor de reglas permitirá definir condiciones de negocio evaluables sin recompilación, con trazabilidad de qué regla se aplicó a cada transacción. & No & - & - \\
\addlinespace[2pt]
RT-16.15 & La solución gestionará documentos con versionado, metadatos, control de acceso, búsqueda por contenido y por metadato, y previsualización sin descarga. & No & - & - \\
\addlinespace[2pt]
RT-16.16 & Los documentos se almacenarán cifrados, con verificación de integridad y retención conforme a la política del caso. & Sí & Expediente de Evidencia (M-16), custodiado bajo política WORM, sellado criptográfico SHA-256 encadenado y almacenamiento cifrado & Subdocumento 5, sección 5.6 ``Modelo de evidencia de consentimiento'' \\
\addlinespace[2pt]
RT-16.17 & La solución soportará firma electrónica conforme a la Ley N° 19.799, con firma avanzada para los actos que el caso lo requiera, y verificación de validez del certificado al momento de la firma. & Sí & Firma electrónica avanzada (M-16) integrada con prestador acreditado por la autoridad chilena (Acepta, E-Sign o TOC Biometrics), conforme a Ley N° 19.799, con verificación de validez del certificado al firmar & Subdocumento 4, sección 4.1.7 (punto 3, cita explícita RT-16.17-T) y Tabla de stack tecnológico (fila ``Firma electrónica avanzada (M-16)'') \\
\addlinespace[2pt]
RT-16.18 & Se generará el sello de tiempo y se conservará la evidencia de firma que permita verificar el documento con posterioridad al vencimiento del certificado. & Sí & Expediente de consentimiento sellado criptográficamente (SHA-256 encadenado) y custodiado bajo WORM, con retención del plazo del crédito más 6 años & Subdocumento 5, sección 5.6 ``Modelo de evidencia de consentimiento''; Subdocumento 4, sección 4.1.7 \\
\addlinespace[2pt]
RT-16.19 & La solución generará documentos a partir de plantillas administrables por el CLIENTE, con datos de la transacción y salida en formato abierto. & No & - & - \\
\addlinespace[2pt]
RT-16.20 & La solución enviará notificaciones por al menos tres canales: correo electrónico, notificación en la aplicación y mensajería instantánea o SMS, según lo que el caso requiera. & No & - & - \\
\addlinespace[2pt]
RT-16.21 & Las plantillas de notificación serán administrables por el CLIENTE, con variables de la transacción, y versionadas. & No & - & - \\
\addlinespace[2pt]
RT-16.22 & Cada persona usuaria podrá configurar sus preferencias de canal y de frecuencia, respetando las notificaciones que el CLIENTE defina como obligatorias. & No & - & - \\
\addlinespace[2pt]
RT-16.23 & El envío será asíncrono, con reintento ante falla, control de duplicados y registro de entrega, apertura y error por cada mensaje. & No & - & - \\
\addlinespace[2pt]
RT-16.24 & El PROPONENTE declarará el proveedor de cada canal, su costo unitario, su volumen proyectado y el tratamiento del costo variable en la Oferta Económica. & No & - & - \\
\addlinespace[2pt]
RT-16.25 & Las notificaciones respetarán la normativa de comunicaciones comerciales y permitirán la baja cuando corresponda. & No & - & - \\
\addlinespace[2pt]
RT-16.26 & Se valorará la integración con canales conversacionales que permitan a la persona usuaria responder y ejecutar acciones desde el propio canal. & No & - & - \\
\addlinespace[2pt]
RT-16.27 & La solución dispondrá de búsqueda global con indexación de texto completo, tolerancia a errores de escritura, filtros facetados y respeto del control de acceso de la persona que busca. & Sí & Motor de búsqueda y catálogo con Elastic Cloud (Elasticsearch/OpenSearch), búsqueda facetada y reindexación incremental & Subdocumento 4, Tabla de stack tecnológico (fila ``Motor de búsqueda y catálogo (Z1)'') \\
\addlinespace[2pt]
RT-16.28 & Los listados serán ordenables, filtrables, paginados y exportables en formatos abiertos, con el filtro aplicado reflejado en la exportación. & No & - & - \\
\addlinespace[2pt]
RT-16.29 & Las exportaciones de gran volumen se procesarán de forma asíncrona, con notificación al completarse y sin bloquear la sesión. & No & - & - \\
\addlinespace[2pt]
RT-16.30 & Toda exportación de información sensible quedará registrada en la auditoría, con identificación de quién exportó qué y cuándo. & No & - & - \\
\addlinespace[2pt]
RT-16.31 & La solución dispondrá de un portal público con la información que el caso determine, accesible sin autenticación, con los mismos estándares de accesibilidad y desempeño que el resto de la plataforma. & Sí & Portal del Cliente (M-11): expone catálogo público con precio y disponibilidad y consulta de estado de pedido sin autenticación, servido tras CDN & Subdocumento 3, Tabla 3.9 (M-11 Portal del Cliente); Subdocumento 4, Tabla de stack tecnológico (fila ``Portal del cliente (Capa 1, Z1)'') \\
\addlinespace[2pt]
RT-16.32 & Las personas usuarias externas dispondrán de autoatención para las consultas de mayor frecuencia, evitando el contacto telefónico para operaciones simples. & Sí & Portal del Cliente (M-11): consulta de estado de pedido y notificaciones transaccionales sin necesidad de contacto reiterado a la compañía & Subdocumento 3, Tabla 3.9 (M-11 Portal del Cliente) \\
\addlinespace[2pt]
RT-16.33 & El PROPONENTE estimará la reducción esperada del volumen de atención asistida por efecto de la autoatención y comprometerá el indicador. & No & - & - \\
\addlinespace[2pt]
RT-16.34 & El portal público resistirá picos de tráfico sin degradar los servicios transaccionales internos, mediante aislamiento de recursos y caché. & Sí & Portal del cliente servido tras CDN, absorbiendo el pico del evento anual mayoritariamente en la capa de borde & Subdocumento 4, Tabla de stack tecnológico (fila ``Portal del cliente (Capa 1, Z1)'') \\
\addlinespace[2pt]
RT-17.01 & La solución proveerá una aplicación móvil para los perfiles operacionales que el caso identifique, con funcionamiento sin conexión y sincronización diferida cuando la operación en terreno lo requiera. & Sí & Aplicación web progresiva (PWA) de Punto de Venta con service workers e IndexedDB, replicada en el nodo de tienda para operar 8 horas sin conexión & Subdocumento 4, Tabla de stack tecnológico (fila ``Punto de venta (Capa 1, Z2)'') \\
\addlinespace[2pt]
RT-17.02 & El PROPONENTE declarará si la aplicación es nativa, híbrida o web progresiva, y justificará la decisión frente al requisito de operación desconectada, al acceso a periféricos y al costo de mantención. & Sí & Decisión explícita por PWA, justificada por la necesidad de ejecutar indistintamente en las 640 terminales compartidas sin instalación por dispositivo & Subdocumento 4, sección 4.0.2 (línea sobre ``Hardware compartido'') y Tabla de stack tecnológico \\
\addlinespace[2pt]
RT-17.03 & La aplicación soportará las versiones de sistema operativo móvil vigentes y las dos anteriores, con política de actualización declarada. & No & - & - \\
\addlinespace[2pt]
RT-17.04 & La aplicación se distribuirá por las tiendas oficiales o mediante gestión de flota corporativa, con firma de la aplicación y verificación de integridad. & No & - & - \\
\addlinespace[2pt]
RT-17.05 & La información almacenada en el dispositivo estará cifrada, con borrado remoto y bloqueo ante pérdida o desvinculación de la persona usuaria. & No & - & - \\
\addlinespace[2pt]
RT-17.06 & La aplicación integrará los periféricos que el caso requiera: cámara, lector de códigos, NFC, GPS, impresora de etiquetas, balanza o báscula. & Sí & Terminal móvil compartido con lector de código integrado; integración de la red de videovigilancia (1.840 cámaras) como periférico de red & Subdocumento 4, Tabla 4.57 ``Equipamiento de terreno y dispositivos operacionales'' (cita explícita RT-17.06) \\
\addlinespace[2pt]
RT-17.07 & El consumo de datos móviles y de batería se optimizará y se declarará el consumo estimado por turno de trabajo. & No & - & - \\
\addlinespace[2pt]
RT-17.08 & Se valorará la disponibilidad de una versión de la interfaz para dispositivos de bajo costo o de generaciones anteriores, ampliando la cobertura de personas usuarias. & No & - & - \\
\addlinespace[2pt]
RT-18.01 & El PROPONENTE declarará cada componente de inteligencia artificial: propósito, modelo o servicio empleado, proveedor, versión y ubicación de procesamiento de los datos. & No & - & - \\
\addlinespace[2pt]
RT-18.02 & Se garantizará contractualmente que los datos del CLIENTE no serán utilizados para entrenar modelos de terceros, salvo autorización expresa y escrita. & No & - & - \\
\addlinespace[2pt]
RT-18.03 & Se documentarán los límites de uso, los casos en que el resultado requiere validación humana previa y el procedimiento de supervisión. & No & - & - \\
\addlinespace[2pt]
RT-18.04 & Los riesgos de sesgo, alucinación, fuga de información y uso indebido se evaluarán y mitigarán conforme al NIST AI Risk Management Framework 1.0 y a la norma ISO/IEC 42001. & No & - & - \\
\addlinespace[2pt]
RT-18.05 & Toda interacción relevante con un componente de inteligencia artificial quedará registrada para efectos de auditoría y trazabilidad, incluyendo la entrada, la salida y la decisión humana posterior. & No & - & - \\
\addlinespace[2pt]
RT-18.06 & El componente podrá desactivarse sin comprometer la operación del resto de la solución, y existirá un procedimiento manual de respaldo para la función que automatiza. & No & - & - \\
\addlinespace[2pt]
RT-18.07 & Cuando el resultado se presente a una persona usuaria, se indicará expresamente que fue generado o sugerido de forma automática, y su nivel de confianza cuando el modelo lo provea. & No & - & - \\
\addlinespace[2pt]
RT-18.08 & Los modelos predictivos declararán sus variables de entrada, su métrica de desempeño, su línea base y su plan de reentrenamiento y de detección de deriva. & No & - & - \\
\addlinespace[2pt]
RT-18.09 & La responsabilidad por los resultados de los componentes de inteligencia artificial recae íntegramente en el ADJUDICATARIO, conforme al Artículo 86° de las Bases Administrativas. & No & - & - \\
\addlinespace[2pt]
RT-18.10 & Se valorará el uso de automatización robótica de procesos o de agentes para tareas repetitivas de back office, con el ahorro de horas cuantificado y reflejado en la Oferta Económica. & No & - & - \\
\addlinespace[2pt]
RT-19.01 & El ADJUDICATARIO constituirá una oficina de gestión del proyecto con metodología declarada, basada en el PMBOK del Project Management Institute e integrando prácticas ágiles donde el PROPONENTE lo justifique. & No & - & - \\
\addlinespace[2pt]
RT-19.02 & El Jefe de Proyecto tendrá dedicación exclusiva durante toda la fase de implementación y facultades para comprometer al ADJUDICATARIO en materias de ejecución. & Sí & Jefe de Proyecto (Vicente Rosales Miranda, PMP) al frente de un ``equipo experto con dedicación exclusiva al proyecto'' & Subdocumento 1, sección 1.6 ``Frente de jefaturas asignado'' y Tabla 1 \\
\addlinespace[2pt]
RT-19.03 & Existirá un procedimiento formal de control de cambios conforme al Artículo 72° de las Bases Administrativas, con registro de cambios, análisis de impacto y aprobación previa a la ejecución. & No & - & - \\
\addlinespace[2pt]
RT-19.04 & La gestión del riesgo seguirá la norma ISO 31000, con registro de riesgos vivo, revisión en cada Comité de Proyecto y planes de mitigación con responsable, plazo y disparador. & No & - & - \\
\addlinespace[2pt]
RT-19.05 & El ADJUDICATARIO habilitará un espacio colaborativo accesible al CLIENTE, con la documentación del proyecto, los entregables, las actas, el registro de riesgos y el registro de cambios siempre actualizados. & No & - & - \\
\addlinespace[2pt]
RT-19.06 & EL ADJUDICATARIO entregará un informe mensual de avance con estado del cronograma, avance físico y financiero, entregables del período, desviaciones, riesgos, incidencias y compromisos del período siguiente. & No & - & - \\
\addlinespace[2pt]
RT-19.07 & El avance se medirá con valor ganado, reportando el índice de desempeño del cronograma y del costo, y no mediante declaración cualitativa de porcentaje de avance. & No & - & - \\
\addlinespace[2pt]
RT-19.08 & El CLIENTE dispondrá de un tablero de estado del proyecto actualizado, accesible en cualquier momento. & No & - & - \\
\addlinespace[2pt]
RT-19.09 & Las actas de todos los comités del Artículo 71 de las Bases Administrativas se levantarán dentro de los dos días hábiles siguientes y registrarán acuerdos, responsables y plazos. & No & - & - \\
\addlinespace[2pt]
RT-19.10 & Toda desviación superior al 10\% en un hito se comunicará dentro de los cinco días hábiles de detectada, con plan de recuperación. & No & - & - \\
\addlinespace[2pt]
RT-20.01 & El PROPONENTE presentará un plan de implantación por sitio y por perfil, coherente con el cronograma obligatorio del Artículo 17° de las Bases Administrativas. & No & - & - \\
\addlinespace[2pt]
RT-20.02 & La estrategia de paso a producción será gradual y reversible. Se declarará el criterio de avance entre olas y el procedimiento de reversión, con su tiempo de ejecución. & No & - & - \\
\addlinespace[2pt]
RT-20.03 & Durante la marcha blanca, la solución convivirá con la operación vigente del CLIENTE, con conciliación diaria y sin doble digitación no declarada. & Sí & Condiciones de cierre de marcha blanca y estrategia de migración de cartera con conciliación diaria automatizada de saldos entre plataforma legada y nueva & Subdocumento 3, sección 3.11.2 (Tabla 3.30) y sección 3.2.5 (SUP-25) \\
\addlinespace[2pt]
RT-20.04 & El PROPONENTE definirá los indicadores diarios que se medirán durante la marcha blanca y sus umbrales de avance, conforme al Artículo 17.3 de las Bases Administrativas. & No & - & - \\
\addlinespace[2pt]
RT-20.05 & Se dispondrá de acompañamiento en terreno durante las primeras semanas de cada ola, con dotación declarada y decreciente según la curva de adopción. & No & - & - \\
\addlinespace[2pt]
RT-20.06 & Se establecerá un período de estabilización con atención reforzada tras cada paso a producción, con dotación y duración declaradas y sin costo adicional. & No & - & - \\
\addlinespace[2pt]
RT-20.07 & Existirá una definición de terminado acordada con el CLIENTE, aplicable a cada entregable, que incluya código, pruebas, documentación, seguridad y despliegue. & No & - & - \\
\addlinespace[2pt]
RT-20.08 & El protocolo de aceptación de cada hito se formalizará conforme al Formulario T-17, con criterios objetivos y verificables. & No & - & - \\
\addlinespace[2pt]
RT-21.01 & El ADJUDICATARIO dispondrá de un centro de operaciones de red con cobertura 24x7x365, propio o subcontratado, y declarará su ubicación, dotación por turno y procedimientos. & Sí & Soporte y operación de misión crítica 24x7x365 para el canal digital y los componentes financieros durante los 36 meses de la fase de Operación & Subdocumento 3, sección 3.6.5 ``Fase de Operación'' \\
\addlinespace[2pt]
RT-21.02 & Se designará un gerente de servicio dedicado como contraparte permanente del CLIENTE durante la fase de Operación. & No & - & - \\
\addlinespace[2pt]
RT-21.03 & El equipo de operación contará con especialistas por tecnología, nominados y con dedicación declarada. & No & - & - \\
\addlinespace[2pt]
RT-21.04 & Los procedimientos de operación estarán documentados y serán ejecutables por el personal del CLIENTE tras la transferencia de conocimiento. & No & - & - \\
\addlinespace[2pt]
RT-21.05 & El PROPONENTE dispondrá de un centro de atención adecuado para soportar integralmente las consultas, propio o subcontratado con un operador de nivel acreditado. & No & - & - \\
\addlinespace[2pt]
RT-21.06 & El centro de atención cumplirá un tiempo de atención al usuario final de 80\% antes de 20 segundos y una resolución al primer contacto de al menos 70\%. La tasa de abandono no superará el 5\%. & No & - & - \\
\addlinespace[2pt]
RT-21.07 & El horario mínimo de atención será de 8:00 a 20:00 en días hábiles, ampliado a 24x7 para los incidentes de severidad crítica y para la ventana operacional que defina el caso. & Sí & Atención en horario comercial extendido (09:00 a 23:00) para sucursales y centros de distribución, y cobertura 24x7x365 para canal digital y servicios financieros críticos & Subdocumento 4, tabla de ventana operacional (sección 4.0, cita explícita RT-21.06-C/RT-21.07-T) \\
\addlinespace[2pt]
RT-21.08 & La atención cubrirá orientación funcional de distintos grados de complejidad, desde preguntas simples hasta situaciones complejas, y las preguntas técnicas más frecuentes sobre la aplicación y su entorno de uso. & No & - & - \\
\addlinespace[2pt]
RT-21.09 & Se mantendrá un registro histórico de las actividades de soporte que permita monitorear y gestionar los principales requerimientos, con análisis de tendencia mensual. & No & - & - \\
\addlinespace[2pt]
RT-21.10 & Existirá un proceso de aprendizaje del centro de soporte que acumule conocimiento sobre las complejidades de la operación y las necesidades de las personas usuarias, reflejado en la base de conocimiento. & No & - & - \\
\addlinespace[2pt]
RT-21.11 & Se proveerán servicios de capacitación en línea en modalidad de autoformación, con registro de avance para efectos de monitoreo. & No & - & - \\
\addlinespace[2pt]
RT-21.12 & Se habilitarán espacios de interacción entre las entidades y personas usuarias del sistema, que permitan compartir experiencias de uso. & No & - & - \\
\addlinespace[2pt]
RT-21.13 & Los indicadores clave del proceso de soporte estarán abiertos al CLIENTE, y los indicadores específicos serán accesibles según nivel de autorización. & No & - & - \\
\addlinespace[2pt]
RT-21.14 & El PROPONENTE dimensionará la dotación del centro de atención con fundamento cuantitativo, empleando teoría de colas o el modelo Erlang C, a partir del volumen de contactos proyectado del caso. & No & - & - \\
\addlinespace[2pt]
RT-21.15 & Existirá un canal único de registro de incidentes y solicitudes, con número de ticket, clasificación por severidad y seguimiento del ciclo de vida completo hasta el cierre conforme. & No & - & - \\
\addlinespace[2pt]
RT-21.16 & Cuando el caso comprenda sitios alejados, los especialistas de niveles 2 y 3 deberán trasladarse cuando la resolución lo requiera, por el medio más rápido disponible y con disponibilidad de tiempo suficiente, sin alterar la atención normal del resto de los sitios. El costo del traslado está incluido en la oferta. & No & - & - \\
\addlinespace[2pt]
RT-21.17 & El cierre de un ticket requerirá confirmación de la persona usuaria o transcurso del plazo de confirmación automática declarado. & No & - & - \\
\addlinespace[2pt]
RT-21.18 & El ADJUDICATARIO reportará mensualmente el cumplimiento de los niveles de servicio de atención, con el detalle por severidad y el análisis de los incumplimientos. & No & - & - \\
\addlinespace[2pt]
RT-21.19 & El PROPONENTE declarará el tamaño de la bolsa anual de horas de mantención evolutiva, su composición por perfil y el procedimiento de solicitud, estimación, aprobación y liquidación. & No & - & - \\
\addlinespace[2pt]
RT-21.20 & La mantención evolutiva se someterá al mismo estándar de calidad, pruebas y seguridad que el desarrollo original. & No & - & - \\
\addlinespace[2pt]
RT-21.21 & El PROPONENTE mantendrá un registro de deuda técnica cuantificado y destinará una fracción declarada de la capacidad de la Operación a reducirla. & No & - & - \\
\addlinespace[2pt]
RT-21.22 & Las actualizaciones de versión de los componentes de base se planificarán anualmente, con ventana acordada y plan de reversión. & No & - & - \\
\addlinespace[2pt]
RT-22.01 & El plan de capacitación se estructurará por perfil: personas usuarias finales, usuarias avanzadas, administradoras, equipo técnico y soporte de niveles 1 y 2. & No & - & - \\
\addlinespace[2pt]
RT-22.02 & Las modalidades incluirán capacitación presencial en cada sitio de operación, sesiones en línea sincrónicas, autoformación y acompañamiento en puesto de trabajo durante la marcha blanca. & No & - & - \\
\addlinespace[2pt]
RT-22.03 & Todo el material de capacitación se entregará en español, en formato editable y de propiedad del CLIENTE: manuales por perfil, guías rápidas, preguntas frecuentes, videos tutoriales y base de conocimiento consultable. & No & - & - \\
\addlinespace[2pt]
RT-22.04 & La capacitación no podrá afectar la operación del CLIENTE: se programará por turnos y en horarios acordados, considerando la estacionalidad y la ventana operacional del caso. & No & - & - \\
\addlinespace[2pt]
RT-22.05 & Las personas usuarias administradoras y el equipo técnico del CLIENTE serán evaluados y certificados como condición para el cierre de cada marcha blanca. & No & - & - \\
\addlinespace[2pt]
RT-22.06 & Durante la Operación se ejecutarán al menos dos jornadas anuales de actualización y se capacitará al personal nuevo del CLIENTE sin costo adicional. & No & - & - \\
\addlinespace[2pt]
RT-22.07 & El ADJUDICATARIO proveerá acompañamiento y mentoría al equipo técnico del CLIENTE durante al menos seis meses posteriores al paso a producción de la Etapa 2. & No & - & - \\
\addlinespace[2pt]
RT-22.08 & La base de conocimiento será mantenida y actualizada durante todo el Contrato, con métricas de uso y de utilidad percibida. & No & - & - \\
\addlinespace[2pt]
RT-22.09 & Se valorará la existencia de un ambiente permanente de entrenamiento, con datos ficticios, disponible para la práctica sin riesgo. & No & - & - \\
\addlinespace[2pt]
RT-23.01 & Información institucional: descripción detallada del giro principal, historia y trayectoria de la empresa con al menos tres años de antigüedad comprobable, misión, visión y valores, certificaciones y acreditaciones vigentes, y presencia geográfica y oficinas. & Sí & Sitio web corporativo declarado, que según el proponente aloja esta información institucional (giro, historia de 10 años, misión, visión, valores, certificaciones) & Subdocumento 1, Anexo A.1 ``Presencia digital y web corporativa'' y secciones 1.1, 1.3 y 1.5 \\
\addlinespace[2pt]
RT-23.02 & Experiencia y casos de éxito: portafolio de proyectos similares de los últimos cinco años, casos documentados con métricas verificables, testimonios de clientes con autorización de publicación, industrias atendidas con énfasis en la del caso asignado, y volumen de transacciones o de personas usuarias gestionadas. & Sí & Portafolio de tres proyectos de misión crítica con métricas verificables, declarado como disponible en la plataforma oficial & Subdocumento 1, sección 1.4 ``Experiencia relevante en la industria del caso'' y Anexo A.1 \\
\addlinespace[2pt]
RT-23.03 & Equipo profesional: organigrama del equipo directivo, perfiles de socios y directores, currículos resumidos del equipo técnico clave y certificaciones profesionales del personal. & Sí & Tabla de roles estratégicos y credenciales del equipo de jefaturas, declarada como parte de la información disponible en el sitio corporativo & Subdocumento 1, sección 1.6 y Tabla 1; Anexo A.1 \\
\addlinespace[2pt]
RT-23.04 & Capacidades técnicas: servicios y soluciones ofrecidas, conjunto tecnológico dominado, alianzas tecnológicas con proveedores de nube y de plataforma, metodologías de trabajo certificadas, e infraestructura y capacidad instalada. & Sí & Capacidades instaladas (fábrica de software, célula de datos, célula de ciberseguridad, mesa de servicio) y alianzas tecnológicas con AWS y Microsoft Azure & Subdocumento 1, secciones 1.2 y 1.5; Anexo A.1 \\
\addlinespace[2pt]
RT-23.05 & El sitio será accesible conforme a WCAG 2.2 nivel AA, responsivo y con certificado TLS válido y vigente. & No & - & - \\
\addlinespace[2pt]
RT-23.06 & El sitio se mantendrá activo y disponible durante todo el proceso de licitación, desde el registro del participante hasta la adjudicación. & No & - & - \\
\addlinespace[2pt]
RT-23.07 & Material educativo digital: seminarios en línea, artículos técnicos o centro de recursos y documentación. & No & - & - \\
\addlinespace[2pt]
RT-23.08 & Demostración en línea, recorrido virtual de la solución o portal de soporte para clientes. & No & - & - \\
\addlinespace[2pt]
RT-23.09 & Métricas de disponibilidad y desempeño de los servicios del proponente publicadas en tiempo real. & No & - & - \\
\addlinespace[2pt]
RT-23.10 & Calculadora de retorno de la inversión o herramienta de estimación pertinente a la industria del caso. & No & - & - \\
\addlinespace[2pt]
RT-24.01 & Todos los integrantes clave del equipo deberán aparecer individualmente presentándose, en toma de medio cuerpo, mirando directamente a la cámara, con vestimenta formal o de negocio informal, y con una duración mínima de aparición de cros por persona. & No & - & - \\
\addlinespace[2pt]
RT-24.02 & Cada vez que aparezca un integrante deberá mostrarse en pantalla su nombre completo, su cargo en el proyecto, sus años de experiencia y su especialización relevante. & No & - & - \\
\addlinespace[2pt]
RT-24.03 & Las certificaciones principales de cada integrante podrán mostrarse en pantalla. & No & - & - \\
\addlinespace[2pt]
RT-24.04 & El fondo será el mismo para todas las tomas de personas, con iluminación consistente y encuadre similar para todos los participantes. & No & - & - \\
\addlinespace[2pt]
RT-24.05 & La paleta de colores corporativa, la tipografía y el uso de la marca serán consistentes en todo el video, con logotipo visible en todas las escenas y datos de contacto en el encabezado o pie. & No & - & - \\
\addlinespace[2pt]
RT-24.06 & Los niveles de volumen estarán normalizados, con la misma calidad de grabación y sin variaciones bruscas de sonido. & No & - & - \\
\addlinespace[2pt]
RT-25.01 & Navegación intuitiva sin necesidad de manual, con un máximo de tres interacciones para alcanzar cualquier funcionalidad principal. & No & - & - \\
\addlinespace[2pt]
RT-25.02 & Retroalimentación visual clara ante cada acción, estados de carga y transición fluidos, y manejo elegante de los errores. & No & - & - \\
\addlinespace[2pt]
RT-25.03 & Cumplimiento de WCAG 2.2 nivel AA: contraste adecuado, tamaños de fuente legibles, objetivos táctiles de al menos 44x44 píxeles, navegación completa por teclado, y textos alternativos y etiquetas de accesibilidad. & No & - & - \\
\addlinespace[2pt]
RT-25.04 & Diseño responsivo demostrado en escritorio (1920x1080 y 1366x768), tableta (768 x 1024 en vertical y horizontal) y teléfono (375x812), con puntos de quiebre coherentes. & No & - & - \\
\addlinespace[2pt]
RT-25.05 & Sistema de diseño documentado: paleta de a lo más cinco colores principales, a lo más dos familias tipográficas jerarquizadas, iconografía coherente, retícula y espaciado consistentes, y componentes reutilizables. & No & - & - \\
\addlinespace[2pt]
RT-25.06 & El prototipo se entregará como enlace navegable en línea, sin instalación ni complementos, compatible con los navegadores modernos vigentes. & No & - & - \\
\addlinespace[2pt]
RT-25.07 & El acceso al prototipo estará garantizado por un mínimo de seis meses desde su entrega. & No & - & - \\
\addlinespace[2pt]
RT-25.08 & El prototipo permitirá navegación completa entre pantallas, simulación de ingreso de datos, transiciones y animaciones básicas, estados de posado y activo, simulación de carga de archivos, elementos desplegables funcionales y simulación de validaciones. & No & - & - \\
\addlinespace[2pt]
RT-25.09 & Se entregará la arquitectura de información del prototipo: mapa del sitio completo, diagrama de navegación, taxonomía y nomenclatura, estructura de menús y jerarquía de la información. & No & - & - \\
\addlinespace[2pt]
RT-25.10 & Se valorará la posibilidad de exportar el prototipo a un documento interactivo para revisión sin conexión. & No & - & - \\
\addlinespace[2pt]
RT-26.01 & Cada innovación se ubicará explícitamente en la arquitectura: qué capa la contiene, qué componentes la implementan y qué interfaces consume o expone. & Sí & Las 5 fichas del Formulario T-19 declaran explícitamente el campo ``Dónde se inserta en la arquitectura'' (capa, zona y módulo) para cada una de las cinco innovaciones & Subdocumento 13, Anexos A.1 a A.5 (campo ``Dónde se inserta en la arquitectura'') \\
\addlinespace[2pt]
RT-26.02 & Cada innovación identificará los paquetes de la estructura de descomposición del trabajo que la ejecutan y el mes del cronograma en que se materializa. & No & - & - \\
\addlinespace[2pt]
RT-26.03 & Las innovaciones de base tecnológica declararán el nivel de madurez de la tecnología con la escala utilizada y citarán las fuentes en norma APA 7.ª edición. & Sí & Cada una de las cinco fichas declara ``Nivel de madurez y escala utilizada'' (escala TRL) y ``Fuentes citadas (APA 7.a ed.)'' con referencias reales & Subdocumento 13, Anexos A.1 a A.5 (campos ``Nivel de madurez y escala utilizada'' y ``Fuentes citadas'') \\
\addlinespace[2pt]
RT-26.04 & Cada innovación declarará su riesgo de adopción, su probabilidad, su impacto, la estrategia de mitigación y el plan de contingencia si no rinde lo esperado. & No & - & - \\
\addlinespace[2pt]
RT-26.05 & Cada innovación declarará su indicador de verificación con línea base, meta y momento de medición, y su impacto en inversión, costo operacional y beneficio esperado. & No & - & - \\
\addlinespace[2pt]
RT-26.06 & Las innovaciones que incorporen inteligencia artificial cumplirán íntegramente el Capítulo 18 de este documento. & No & - & - \\
\addlinespace[2pt]
RT-26.07 & Las innovaciones que modifiquen la arquitectura de seguridad requerirán su propio modelado de amenazas. & No & - & - \\
\addlinespace[2pt]
RT-26.08 & Se valorará que al menos una innovación sea verificable durante la marcha blanca de la Etapa 1, es decir, que su beneficio pueda medirse antes del mes 16. & No & - & - \\
\end{longtable}

\subsection{A.5. Listado de Supuestos, Exclusiones y Restricciones}\label{a-5-listado-de-supuestos-exclusiones-y-restricciones}

Este anexo consolida los supuestos declarados por Only Simple Solutions, las exclusiones del alcance y las restricciones que condicionan la propuesta. Los supuestos corresponden a decisiones propuestas por el equipo y deben mantenerse coherentes con el alcance, la arquitectura y la oferta económica.

\subsubsection{A.5.1. Listado de Supuestos}\label{a-5-1-listado-de-supuestos}

A continuación, se detalla el listado de supuestos propuestos para abordar las decisiones pendientes del numeral 16.1 del Caso 09. Ningún supuesto podrá reducir los requisitos obligatorios establecidos en las Bases.

\renewcommand{\thetable}{A.5.1}
\begin{longtable}[]{@{}
  >{\raggedright\arraybackslash}p{(\columnwidth - 6\tabcolsep) * \real{0.1010}}
  >{\raggedright\arraybackslash}p{(\columnwidth - 6\tabcolsep) * \real{0.4952}}
  >{\raggedright\arraybackslash}p{(\columnwidth - 6\tabcolsep) * \real{0.4038}}@{}}
\caption{Supuestos declarados del análisis.}\label{tab:a-5-1}\\
\toprule\noalign{}
\begin{minipage}[b]{\linewidth}\raggedright
\textbf{ID}
\end{minipage} & \begin{minipage}[b]{\linewidth}\raggedright
\textbf{Descripción del Supuesto}
\end{minipage} & \begin{minipage}[b]{\linewidth}\raggedright
\textbf{Fundamento}
\end{minipage} \\
\midrule\noalign{}
\endfirsthead
\toprule\noalign{}
\begin{minipage}[b]{\linewidth}\raggedright
\textbf{ID}
\end{minipage} & \begin{minipage}[b]{\linewidth}\raggedright
\textbf{Descripción del Supuesto}
\end{minipage} & \begin{minipage}[b]{\linewidth}\raggedright
\textbf{Fundamento}
\end{minipage} \\
\midrule\noalign{}
\endhead
\bottomrule\noalign{}
\endlastfoot
\textbf{SUP-01} & Se propone calcular el disponible para vender desde una fuente autorizada, descontando reservas y un margen de seguridad diferenciado por categoría y punto. & El registro de inventario presenta un 12,4 \% de discrepancia, por lo que la existencia registrada no basta para comprometer unidades. \\
\addlinespace[2pt]
\textbf{SUP-02} & Se propone transformar progresivamente el sistema central y sus integraciones, manteniendo la continuidad operacional y preservando el ERP. & La operación depende de nueve plataformas y catorce interfaces cuyo mapa completo debe levantarse. \\
\addlinespace[2pt]
\textbf{SUP-03} & Se propone relacionar la identidad de una misma persona mediante un identificador común, conservando separados sus registros comerciales y financieros. & Compartir la identidad del cliente no autoriza mezclar los datos de ambos negocios. \\
\addlinespace[2pt]
\textbf{SUP-04} & Se propone permitir únicamente intercambios de datos expresamente autorizados entre retail y filial emisora, con finalidad, dirección, fundamento y controles documentados. & La separación entre negocios debe resolverse antes de construir una vista unificada del cliente. \\
\addlinespace[2pt]
\textbf{SUP-05} & Se propone conservar evidencia del consentimiento vinculada a las condiciones aceptadas, íntegra y recuperable durante todo el plazo del crédito más seis años. & Existen 1.240 repactaciones sin evidencia recuperable y una práctica actual de conservación insuficiente. \\
\addlinespace[2pt]
\textbf{SUP-06} & Se propone registrar la entrega de información precontractual antes de la aceptación, impidiendo que se omita para acelerar la venta. & La rapidez comercial no puede sustituir la entrega de información ni su acreditación. \\
\addlinespace[2pt]
\textbf{SUP-07} & Se propone utilizar reglas y límites pre aprobados para el crédito de contingencia durante la desconexión, manteniendo autonomía de al menos 24 horas continuas. & La continuidad de venta y cobro debe contemplar el crédito sin depender de una evaluación remota. \\
\addlinespace[2pt]
\textbf{SUP-08} & Se propone aplicar una política que respete el precio exhibido y, ante discrepancias acreditadas de una misma oferta, cobrar el menor precio aplicable y registrar la diferencia. & El desfase entre etiquetas y caja no debe trasladarse al cliente; su consecuencia económica debe dimensionarse. \\
\addlinespace[2pt]
\textbf{SUP-09} & Se propone evaluar la actualización manual verificable de etiquetas frente a la alternativa de etiquetas electrónicas. & La instalación de etiquetas electrónicas no es obligatoria, pero su evaluación, especificación y costeo sí lo son. \\
\addlinespace[2pt]
\textbf{SUP-10} & Se propone conservar un historial del precio publicado por canal, con vigencia y evidencia de actualización física, durante tres años. & Actualmente no puede acreditarse qué precio estaba publicado en un instante determinado. \\
\addlinespace[2pt]
\textbf{SUP-11} & Se propone condicionar el cobro a la confirmación de disponibilidad y establecer un procedimiento para atender pedidos ya cobrados cuya unidad no esté disponible. & Las cancelaciones del evento evidencian compromisos y cobros sobre existencia inexistente. \\
\addlinespace[2pt]
\textbf{SUP-12} & Se propone utilizar reservas temporales de existencia, con liberación automática al vencer o cancelarse la compra. Su duración y condiciones deberán quedar parametrizadas. & Una reserva excesiva limita ventas y una insuficiente facilita la sobreventa. \\
\addlinespace[2pt]
\textbf{SUP-13} & Se propone asignar el punto de despacho considerando costo total de servir, disponibilidad confiable, capacidad y efecto sobre la operación de tienda. & La asignación actual por distancia no contempla todos los efectos operacionales. \\
\addlinespace[2pt]
\textbf{SUP-14} & Se propone reconocer de manera trazable la participación de la tienda y del vendedor en ventas que se originan en otro canal. & El despacho digital compite con las metas y comisiones de la venta presencial. \\
\addlinespace[2pt]
\textbf{SUP-15} & Se propone recibir y registrar en tienda las devoluciones de marketplace, notificando al vendedor y separando la atención al cliente de la recuperación de costos. & Actualmente no existe una regla escrita y el vendedor conoce la devolución al revisar su liquidación. \\
\addlinespace[2pt]
\textbf{SUP-16} & Se propone evaluar a los vendedores de marketplace mediante indicadores, períodos, umbrales y consecuencias conocidos y trazables. & La experiencia del cliente depende de terceros, aunque la relación comercial se presenta bajo la marca Ancoa. \\
\addlinespace[2pt]
\textbf{SUP-17} & Se propone clasificar las diferencias de inventario en pérdida física, daño no registrado, error de recepción, devolución mal reintegrada y error de digitación. & La cifra agregada de merma mezcla causas que requieren respuestas diferentes. \\
\addlinespace[2pt]
\textbf{SUP-18} & Se propone priorizar los conteos cíclicos según categoría, riesgo y rotación, definiendo frecuencias y condiciones de recuento extraordinario sin cerrar las tiendas. & La confiabilidad del disponible depende de medir y mejorar la exactitud del registro. \\
\addlinespace[2pt]
\textbf{SUP-19} & Se propone gestionar el acceso de repositores externos mediante patrocinio del proveedor, identificación y permisos limitados. & No existe vínculo laboral que permita imponerles herramientas o capacitación. \\
\addlinespace[2pt]
\textbf{SUP-20} & Se propone seleccionar la alternativa para el CD de Concepción después de comparar incorporación al WMS, solución específica y continuidad controlada de su operación actual. & Su incorporación debe evaluarse y costearse; no constituye una decisión previamente tomada por el CLIENTE. \\
\addlinespace[2pt]
\textbf{SUP-21} & Se propone que Ancoa atienda directamente la garantía legal y gestione posteriormente la recuperación correspondiente frente al fabricante o vendedor externo. & La atención no puede condicionarse a derivar al consumidor. \\
\addlinespace[2pt]
\textbf{SUP-22} & Se propone establecer un procedimiento común para el retracto, registrando solicitud, devolución de fondos y tratamiento de la unidad recibida. & La resolución caso a caso y el reingreso manual generan inconsistencias. \\
\addlinespace[2pt]
\textbf{SUP-23} & Se propone degradar primero las funciones no críticas durante alta concurrencia, preservando la integridad de pedidos, cobros y evidencia. & La saturación debe gestionarse antes de aceptar operaciones que no puedan cumplirse. \\
\addlinespace[2pt]
\textbf{SUP-24} & Se propone asignar accesos individuales asociados a función y vigencia, incluso cuando se utilicen dispositivos compartidos. & La rotación y contratación temporal aumentan el riesgo de accesos no revocados. \\
\addlinespace[2pt]
\textbf{SUP-25} & Se propone migrar la cartera mediante convivencia controlada, conciliación y un procedimiento de retorno preparado antes del corte definitivo. & Los 620.000 clientes con saldo deben migrar sin pérdida de información, interrupción del servicio ni divergencia de saldos. \\
\addlinespace[2pt]
\textbf{SUP-26} & Para el escenario de planificación se considera enero de 2027 como inicio del proyecto, manteniendo los 56 meses contractuales. & Permite ubicar hitos y congelamientos en fechas calendario; la fecha debe confirmarse antes de cerrar la planificación. \\
\addlinespace[2pt]
\textbf{SUP-27} & Para el dimensionamiento territorial se consideran supuestas las ubicaciones de Illapel, Valparaíso, Talca y Concepcion, y la concentración de once tiendas en el entorno de Santiago. & Estas localizaciones se utilizan como escenario metodológico y no como información confirmada por el mandante. \\
\end{longtable}

\textbf{Fuente:} Elaboración propia a partir del numeral 16.1 de las Bases Técnicas del Caso 09. SUP-26 y SUP-27 corresponden a supuestos metodológicos adicionales del proponente.

\subsubsection{A.5.2. Listado de Exclusiones}\label{a-5-2-listado-de-exclusiones}

A continuación, se detalla el listado de actividades excluidas del alcance, indicando las obligaciones que permanecen bajo responsabilidad del proponente.

\renewcommand{\thetable}{A.5.2}
\begin{longtable}[]{@{}
  >{\raggedright\arraybackslash}p{(\columnwidth - 6\tabcolsep) * \real{0.0865}}
  >{\raggedright\arraybackslash}p{(\columnwidth - 6\tabcolsep) * \real{0.4631}}
  >{\raggedright\arraybackslash}p{(\columnwidth - 6\tabcolsep) * \real{0.4503}}@{}}
\caption{Exclusiones del alcance.}\label{tab:a-5-2}\\
\toprule\noalign{}
\begin{minipage}[b]{\linewidth}\raggedright
\textbf{ID}
\end{minipage} & \begin{minipage}[b]{\linewidth}\raggedright
\textbf{Descripción de la Exclusión}
\end{minipage} & \begin{minipage}[b]{\linewidth}\raggedright
\textbf{Obligación que permanece dentro del alcance}
\end{minipage} \\
\midrule\noalign{}
\endfirsthead
\toprule\noalign{}
\begin{minipage}[b]{\linewidth}\raggedright
\textbf{ID}
\end{minipage} & \begin{minipage}[b]{\linewidth}\raggedright
\textbf{Descripción de la Exclusión}
\end{minipage} & \begin{minipage}[b]{\linewidth}\raggedright
\textbf{Obligación que permanece dentro del alcance}
\end{minipage} \\
\midrule\noalign{}
\endhead
\bottomrule\noalign{}
\endlastfoot
\textbf{EX-01} & Se excluye el reemplazo del ERP y de la emisión de documentos tributarios. & Integrar la solución con el ERP, manteniéndolo como único emisor tributario. \\
\addlinespace[2pt]
\textbf{EX-02} & No se exige instalar etiquetas electrónicas en las 22 tiendas. & Evaluar, especificar y costear la alternativa frente a otras opciones. \\
\addlinespace[2pt]
\textbf{EX-03} & Se excluye la adquisición de dispositivos móviles para el personal de venta. & Especificar las cantidades y características necesarias, considerando dispositivos compartidos. \\
\addlinespace[2pt]
\textbf{EX-04} & Se excluyen el desarrollo de las plataformas de los vendedores de marketplace y la operación de su logística. & Integrarlas, medir su desempeño y resolver las devoluciones en tienda. \\
\addlinespace[2pt]
\textbf{EX-05} & Se excluye la gestión de remuneraciones. & Calcular la base de comisión cuando una venta se origina en un canal y se cumple en otro. \\
\addlinespace[2pt]
\textbf{EX-06} & Se excluye la operación de cobranza judicial. & Mantener el expediente trazable de las operaciones de cobranza y repactación. \\
\addlinespace[2pt]
\textbf{EX-07} & Se excluye la sustitución de las empresas de transporte de última milla. & Integrarlas y mantener la trazabilidad del pedido hasta su entrega. \\
\addlinespace[2pt]
\textbf{EX-08} & Se excluye la ejecución de canalizaciones, obras eléctricas y cableado. & Especificar y costear las obras necesarias; su ejecución corresponde al cliente. \\
\addlinespace[2pt]
\textbf{EX-09} & Se excluye resolver la relación contractual con los administradores de centros comerciales. & Diseñar la continuidad de venta y cobro ante la indisponibilidad de sus enlaces. \\
\addlinespace[2pt]
\textbf{EX-10} & Se excluye la adquisición por el proponente del hardware cuya compra corresponde al cliente. & Especificar exactamente los equipos, cantidades y características conforme a las Bases Transversales. \\
\end{longtable}

\subsubsection{A.5.3. Listado de Restricciones}\label{a-5-3-listado-de-restricciones}

A continuación, se detalla el listado de restricciones del caso y de condiciones contractuales y técnicas que deben respetarse durante el diseño, la implementación y la operación.

\renewcommand{\thetable}{A.5.3}
\begin{longtable}[]{@{}
  >{\raggedright\arraybackslash}p{(\columnwidth - 6\tabcolsep) * \real{0.1024}}
  >{\raggedright\arraybackslash}p{(\columnwidth - 6\tabcolsep) * \real{0.6976}}
  >{\raggedright\arraybackslash}p{(\columnwidth - 6\tabcolsep) * \real{0.2000}}@{}}
\caption{Restricciones de la propuesta.}\label{tab:a-5-3}\\
\toprule\noalign{}
\begin{minipage}[b]{\linewidth}\raggedright
\textbf{ID}
\end{minipage} & \begin{minipage}[b]{\linewidth}\raggedright
\textbf{Descripción de la Restricción}
\end{minipage} & \begin{minipage}[b]{\linewidth}\raggedright
\textbf{Origen}
\end{minipage} \\
\midrule\noalign{}
\endfirsthead
\toprule\noalign{}
\begin{minipage}[b]{\linewidth}\raggedright
\textbf{ID}
\end{minipage} & \begin{minipage}[b]{\linewidth}\raggedright
\textbf{Descripción de la Restricción}
\end{minipage} & \begin{minipage}[b]{\linewidth}\raggedright
\textbf{Origen}
\end{minipage} \\
\midrule\noalign{}
\endhead
\bottomrule\noalign{}
\endlastfoot
\textbf{RES-01} & La separación de datos entre retail y filial emisora deberá quedar definida, implementada técnicamente, documentada y auditable. & Caso 09, capítulo 10, restricción 1. \\
\addlinespace[2pt]
\textbf{RES-02} & Ninguna modificación de condiciones del crédito podrá registrarse sin evidencia recuperable del consentimiento informado, conservada durante el plazo del crédito más seis años. & Caso 09, capítulo 10, restricción 2; numeral 15, RT-05.10. \\
\addlinespace[2pt]
\textbf{RES-03} & La información precontractual deberá entregarse antes de aceptar el crédito y su entrega deberá acreditarse de manera estructurada. & Caso 09, capítulo 10, restricción 3. \\
\addlinespace[2pt]
\textbf{RES-04} & El precio cobrado deberá respetar el precio exhibido y la compañía deberá acreditar el precio publicado en cada canal. & Caso 09, capítulo 10, restricción 4. \\
\addlinespace[2pt]
\textbf{RES-05} & Las tiendas deberán continuar vendiendo y cobrando ante la pérdida del enlace externo, con autonomía local y degradada durante \textbf{≥ 24 horas continuas}. & Caso 09, capítulo 10, restricción 5; Bases Transversales, RT-03.10. \\
\addlinespace[2pt]
\textbf{RES-06} & El ERP deberá mantenerse como único emisor de documentos tributarios. & Caso 09, capítulo 10, restricción 6. \\
\addlinespace[2pt]
\textbf{RES-07} & La garantía legal deberá atenderse directamente en la compañía, sin derivar al consumidor como condición de atención. & Caso 09, capítulo 10, restricción 7. \\
\addlinespace[2pt]
\textbf{RES-08} & La migración de los 620.000 clientes con saldo no admitirá pérdida de información, interrupción del servicio ni divergencia de saldos. & Caso 09, capítulo 10, restricción 8. \\
\addlinespace[2pt]
\textbf{RES-09} & Las intervenciones deberán respetar las cinco ventanas de congelamiento comercial definidas por el caso. & Caso 09, capítulo 10, restricción 9. \\
\addlinespace[2pt]
\textbf{RES-10} & La planificación deberá absorber la fecha del evento anual, anunciada con aproximadamente seis semanas de anticipación, sin desplazar hitos contractuales. & Caso 09, capítulo 10, restricción 10. \\
\addlinespace[2pt]
\textbf{RES-11} & La solución no podrá suponer un dispositivo por vendedor; deberá considerar 640 terminales compartidas para 3.820 personas. & Caso 09, capítulo 10, restricción 11. \\
\addlinespace[2pt]
\textbf{RES-12} & No podrán imponerse herramientas o capacitación a los repositores externos por la vía laboral. & Caso 09, capítulo 10, restricción 12. \\
\addlinespace[2pt]
\textbf{RES-13} & La capacitación y gestión de accesos deberán considerar una rotación anual del 62 \% y 1.900 incorporaciones de temporada durante congelamientos. & Caso 09, capítulo 10, restricción 13. \\
\addlinespace[2pt]
\textbf{RES-14} & La incorporación del CD de Concepción a un sistema de gestión de almacenes deberá evaluarse y costearse; no podrá darse por supuesta. & Caso 09, capítulo 10, restricción 14. \\
\addlinespace[2pt]
\textbf{RES-15} & El levantamiento de las catorce integraciones será parte del trabajo del proponente y no una entrega previa que deba resolver el CLIENTE. & Caso 09, capítulo 10, restricción 15. \\
\addlinespace[2pt]
\textbf{RES-16} & El despliegue deberá ser híbrido, con carga principal en nube pública y componentes on-premise. & Bases Administrativas, artículo 16. \\
\addlinespace[2pt]
\textbf{RES-17} & El contrato tendrá una duración obligatoria de 56 meses: Etapa 1 en meses 1–15 y producción en el mes 16; Etapa 2 en meses 13–20 y producción en el mes 21; operación en meses 21–56. & Bases Administrativas, artículo 17. \\
\addlinespace[2pt]
\textbf{RES-18} & La propuesta deberá incorporar cinco innovaciones obligatorias, una por tipo, trazables con arquitectura, EDT y flujo de caja. & Bases Administrativas, capítulo 5. \\
\addlinespace[2pt]
\textbf{RES-19} & La transición financiera deberá coordinarse con el fin de soporte de la plataforma y el último hito de remediación, ambos en 2029. & Caso 09, numerales 13.1–13.2. \\
\addlinespace[2pt]
\textbf{RES-20} & La solución deberá cumplir, entre otros umbrales, evaluación crediticia \textbf{≤8 s}, propagación de precios a cajas y canal digital \textbf{≤5 min} y disponibilidad publicada actualizada \textbf{≤30 s} después de una venta. & Caso 09, numeral 15, RT-09.01 y RT-05.29. \\
\end{longtable}

\ossFinalPage

\end{document}
